\documentclass[a4paper,10pt]{article}
\usepackage{geometry}
\usepackage[english]{babel}
\usepackage{hyperref}
\usepackage{amsmath,amssymb,stmaryrd,enumerate,latexsym,theorem,bbm}


%%%%%%%%%% Start TeXmacs macros
\catcode`\<=\active \def<{
\fontencoding{T1}\selectfont\symbol{60}\fontencoding{\encodingdefault}}
\newcommand{\assign}{:=}
\newcommand{\asterisk}{\mathord{*}}
\newcommand{\mathd}{\mathrm{d}}
\newcommand{\nobracket}{}
\newcommand{\tmaffiliation}[1]{\\ #1}
\newcommand{\tmem}[1]{{\em #1\/}}
\newcommand{\tmemail}[1]{\\ \textit{Email:} \texttt{#1}}
\newcommand{\tmmathbf}[1]{\ensuremath{\boldsymbol{#1}}}
\newcommand{\tmop}[1]{\ensuremath{\operatorname{#1}}}
\newcommand{\tmtextbf}[1]{{\bfseries{#1}}}
\newcommand{\tmtextit}[1]{{\itshape{#1}}}
\newcommand{\tmtextrm}[1]{{\rmfamily{#1}}}

\newcommand{\cprime}{$'$}


\newenvironment{enumeratenumeric}{\begin{enumerate}[1.] }{\end{enumerate}}
\newenvironment{itemizedot}{\begin{itemize} \renewcommand{\labelitemi}{$\bullet$}\renewcommand{\labelitemii}{$\bullet$}\renewcommand{\labelitemiii}{$\bullet$}\renewcommand{\labelitemiv}{$\bullet$}}{\end{itemize}}


\newenvironment{proof}{\noindent\textbf{Proof\ }}{\hspace*{\fill}$\Box$\medskip}
\newenvironment{proof*}[1]{\noindent\textbf{#1\ }}{\hspace*{\fill}$\Box$\medskip}
\newtheorem{corollary}{Corollary}
\newtheorem{definition}{Definition}
\newtheorem{lemma}{Lemma}
\newtheorem{proposition}{Proposition}
{\theorembodyfont{\rmfamily}\newtheorem{remark}{Remark}}
\newtheorem{theorem}{Theorem}
%%%%%%%%%% End TeXmacs macros

\newcommand{\descriptionparagraphs}[1]{\begin{description}#1\end{description}}
\newcommand{\textdots}[0]{...}


%

\newcommand{\bibliographytext}{Bibliography}
\newcommand{\itemparagraph}[1]{{\paragraphtitle{#1}}}

\geometry{
a4paper,
%total
left=30mm,
right=30mm,
top=30mm,
%bottom=20mm
}

\begin{document}

\title{
  The elliptic stochastic quantization\\
  of some two dimensional Euclidean QFTs
}

\author{
  Sergio Albeverio, Francesco C. De Vecchi and Massimiliano Gubinelli
  \tmaffiliation{Hausdorff Center of Mathematics \&\\
  Institute of Applied Mathematics, \\
  University of Bonn, Germany}
  \tmemail{albeverio@iam.uni-bonn.de}
  \tmemail{francesco.devecchi@uni-bonn.de}
  \tmemail{gubinelli@iam.uni-bonn.de}
}

\maketitle

\begin{abstract}
  We study a class of elliptic SPDEs with additive Gaussian noise on
  $\mathbb{R}^2 \times M$, with $M$ a $d$-dimensional manifold equipped with a
  positive Radon measure, and a real-valued non linearity given by the
  derivative of a smooth potential $V$, convex at infinity and growing at most
  exponentially. For quite general coefficients and a suitable regularity of
  the noise we obtain, via the dimensional reduction principle discussed
  in~{\cite{Albeverio2018elliptic}}, the identity between the law of the
  solution to the SPDE evaluated at the origin with a Gibbs type measure on
  the abstract Wiener space $L^2 (M)$. The results are then applied to the
  elliptic stochastic quantization equation for the scalar field with
  polynomial interaction over $\mathbb{T}^2$, and with exponential interaction
  over $\mathbb{R}^2$ (known also as H{\o}egh-Krohn or Liouville model in the
  literature). In particular for the exponential interaction case, the
  existence and uniqueness properties of solutions to the elliptic equation
  over $\mathbb{R}^{2 + 2}$ is derived as well as the dimensional reduction
  for all values of the ``charge parameter'' $\sigma = \frac{\alpha}{2
  \sqrt{\pi}} < \sqrt{8 \pi}$ for which the model has an Euclidean invariant
  measure (hence also permitting to get the corresponding relativistic
  invariant model on the two dimensional Minkowski space).
\end{abstract}

{\tableofcontents}

\section{Introduction}

The Euclidean approach to quantization of relativistic non-linear wave
equations requires the construction of certain probability measure supported
on distributions satisfying a set of quite strong requirements among which
invariance under the full group of rigid motions of the Euclidean space and
reflection positivity. Such constructions have succeeded in the case where the
dimension $d$ of the Euclidean space is equal to $1, 2, 3$ (that is also the
dimension of the space-time of the original relativistic equation). The
non-linearity is quite general if $d = 1$, whereas it is restricted to the
derivatives of polynomials of even degree respectively suitable superpositions
of trigonometric or exponential functions for $d = 2$, or a cubic monomial
function for $d = 3$ (the $\varphi^4_3$ model). In these cases a number of
interesting physical and mathematical properties of the quantised relativistic
fields models have been put in evidence.

In recent years new methods coming from the study of singular partial
stochastic differential equations (SSPDEs) have been developed and have opened
the possibility of new constructions of such Euclidean measures (and other
similar measures associated with problems coming from other areas of
applications, like statistical mechanics, hydrodynamics, wave propagation in
random media {\textdots}). In these approaches the relevant Euclidean
invariant measures are obtained as invariant measures for certain parabolic
semilinear SPDEs called stochastic quantization equations (SQEs).

Typically the measures to be constructed have the heuristic form
\[ \mu (\mathd \varphi) = Z^{- 1} e^{- S (\varphi)} \mathcal{D} \varphi, \]
where $\varphi = \varphi (x) \in \mathbb{R}$, $x \in \mathbb{R}^d$,
$\mathcal{D} \varphi = \Pi_{x \in \mathbb{R}^d} \mathd \varphi (x)$ is a
``flat measure'', $Z$ is a ``normalisation constant'', $S (\varphi) = S_0
(\varphi) + \lambda S_{\tmop{int}} (\varphi)$ with $S_0 (\varphi) =
\frac{1}{2} \int_{\mathbb{R}^d} | \nabla \varphi (x) |^2 \mathd x + m^2
\int_{\mathbb{R}^d} | \varphi (x) |^2 \mathd x$, and where $\lambda \geqslant
0$, $m \geqslant 0$ are parameters,
\[ S_{\tmop{int}} (\varphi) = \int_{\mathbb{R}^d} V (\varphi (x)) \mathd x, \]
for some (non-linear) measurable real valued function $V$ over $\mathbb{R}$.
Note that $\mu (\mathd \varphi)$ is heuristically invariant under Euclidean
transformations (translations and rotations and also additional reflections,
in the case where $V$ is an even function).

The associated stochastic quantization equation (first introduced
in~{\cite{parisi_perturbation_1981}}) is of the form
\[ \mathd X_{\tau} = (\Delta - m^2) X_{\tau} \mathd t - V' (X_{\tau}) \mathd
   t + \mathd W_{\tau} \]
where $\tau$ is an additional ``computer time'', $\Delta$ is the Laplacian in
$\mathbb{R}^d$, $X_{\tau} = X_{\tau} (x)$, $x \in \mathbb{R}^d$, is a random
variable taking values, for fixed $\tau$, in the space of ``(generalised)
functions'' of $x$, $\mathd W_{\tau} (x)$ is a Gaussian white noise in $\tau$
and $x$, and $V'$ in the derivative of $V$.

Following a terminology used especially when $V$ is a polynomial $P$, we call
the above SPDE the $V (\varphi)_d$--stochastic quantization equation (SQE).
Various detailed results have been obtained on $P (\varphi)_2$ SQEs, see the
introduction of~{\cite{albeverio_invariant_2017}}
and~{\cite{gubinelli_pde_2018,GH18}} for references. In particular
in~{\cite{da_prato_strong_2003}} were obtained a strong solution and a unique
ergodic measure in the case where $\mathbb{R}$ is replaced by the $2$-torus
$\mathbb{T}^2$.

As for the case of the $\varphi^4_3$ SQE, a breakthrough for local in time
solutions of the SQE on the $3$-torus $\mathbb{T}^3$ came
from~{\cite{Hairer2014}} see also~{\cite{CaCh,GuImPe}}. More recently still
for the case of the $3$-torus $\mathbb{T}^3$ a construction of invariant
solutions has been given in~{\cite{albeverio_invariant_2017}}. This can be
looked upon as a construction of the $\varphi^4_3$-measure on $\mathbb{T}^3$,
by methods different from the previous ones provided by mathematical physics
approaches, see references in~{\cite{Glimm1987}}. The methods
in~{\cite{albeverio_invariant_2017}} have been simplified and extended to the
case of $\mathbb{R}^3$ in~{\cite{gubinelli_pde_2018}}, yielding a full
alternative construction of the measures which qualify to be called
$\varphi^4_3$ Euclidean measures. A previous approach following
{\cite{Hairer2014}} is in {\cite{MoWe}}. A number of open problems has been
mentioned in~{\cite{albeverio_invariant_2017,gubinelli_pde_2018}}, the main
one concerns the uniqueness of the invariant measures.

Other approaches are however possible if one understands stochastic
quantization in the broader sense of associating to a given heuristic
``Euclidean measure'' a stochastic process having this very measure as
suitable marginal. A first approach consists in constructing a jump type
process via infinite dimensional quasi-regular jump-type Dirichlet
forms~{\cite{Yoshida,Albeverio2002,Yoshida2008,Yoshida2010}} (for such
Dirichlet forms see also~{\cite{ARu}}), that has by construction the given
Euclidean measure as invariant measure. Another approach obtains the invariant
measure of a SQE in dimension $d$ by looking at solutions of a suitable
associated elliptic SPDE in dimension $d + 2$ and restricting the solution to
the $d$ dimensional Euclidean space obtained sending to zero the two
additional variables. This latter procedure is known in the physical
literature as Parisi-Sourlas dimensional reduction and has been implemented by
using algebraic methods of supersymmetry (see~{\cite{parisi_random_1979}}).
For the case of a regularised non-linear term in the original SQE and
regularised noise it has been studied from a mathematical point of view
in~{\cite{Klein1984}}.

Elliptic stochastic PDEs related to stochastic quantization of the
$\varphi_2^4$ and $\varphi_3^4$ models respectively have been discussed both
on the torus and in all space in~{\cite{GH18}}. A systematic mathematical
investigation of the mechanism of dimensional reduction has been initiated
in~{\cite{Albeverio2018elliptic}}, for the $V (\varphi)_0$ model, for both the
cases where $V$ is a convex function and also the case where $V$ is
non-convex, where an interesting phenomenon of non-uniqueness of solutions has
been put in evidence.

In fact in~{\cite{Albeverio2018elliptic}} an explicit formula for the law of
the solution of a class of elliptic SPDE in $\mathbb{R}^2$ taken at the origin
has been obtained, by means of a rigorous version of dimensional reduction.
Besides proving an instance of elliptic stochastic quantization for the case
of Euclidean dimension $d = 0$, it also provided a realisation of the relation
between a supersymmetric quantum field model and an elliptic SPDE.

\

The present paper has two principal aims. The first one is to prove the
dimensional reduction principle for elliptic SPDEs with non singular noise for
equations in $d + 2$ dimensions with $d > 0$. The second is to extend the
result to elliptic singular SPDE having applications in stochastic
quantization program and analysing the corresponding the dimensional reduction
wrt. to removal of spatial cut-offs.

We restrict our attention to singular SPDEs with $d = 2$ and analyse
polynomial and exponential interactions. In particular, in the case of
potentials of the form $\exp (\sigma \varphi)_2$, for suitable real $\sigma$
we are able to complete the dimensional reduction picture removing all the
spatial cutoff (i.e. in both the ``fictious'' and ``real'' spatial variables).

\

Before passing to a more detailed description of our results, let us stress
the importance of the $\exp (\sigma \varphi)_2$ model (and related ones) in
relativistic and Euclidean quantum field theory. The $\exp (\sigma \varphi)_2$
model over $\mathbb{R}^2$ has been first introduced by
H{\o}egh-Krohn~{\cite{Krohn1971}} (see also~{\cite{Krohn1974}}
and~{\cite{Simon1974}}, p.~178~and~307-313) and constructed for an interaction
of the form
\[ U_g (\varphi) : = \int : e^{\sigma \varphi (z)} : g (z) \mathd z \mathd \nu
   (\sigma) \]
for any positive space cut-off function $g \in L^1 (\mathbb{R}^2) \cap L^2
(\mathbb{R}^2)$ and any positive finite measure $\nu$ with support in $(-
\sqrt{4 \pi}, \sqrt{4 \pi})$, in the sense that the interacting measure
\[ \mu_g (\mathd \varphi) : = \frac{e^{- U_g (\varphi)} \mu_0 (\mathd
   \varphi)}{\int e^{- U (\varphi)} \mu_0 (\mathd \varphi)}, \]
for $m > 0$, is absolutely continuous with respect to Nelson's free field
Gaussian measure $\mu_0$ with mean zero and covariance $(m^2 - \Delta_z)^{-
1}$, $z \in \mathbb{R}^2$, (realised as a probability measure on, for example,
$\mathcal{S}' (\mathbb{R}^2)$). See~{\cite{Krohn1974,Krohn1971}}
and~{\cite{Simon1974}} p.~178 for further properties of the measure $\mu_g$.
This model has been discussed in the infinite volume limit $g \rightarrow 1$,
that is unique, in a further paper~{\cite{Krohn1974}}, under the more
restrictive assumption that $\tmop{supp} (\nu) \subset \left[ -
\frac{4}{\sqrt{\pi}}, \frac{4}{\sqrt{\pi}} \right]$.

\

In fact in~{\cite{Krohn1974}} the authors prove that the measure of the $\exp
(\sigma \varphi)_2$ model for $| \sigma | < \sqrt{4 \pi}$ (or more generally
of suitable superpositions $\int \tmop{Ch} (\sigma \varphi)_2 \mathd \sigma$
of these models) satisfies all axioms of Euclidean quantum field theory and
leads to relativistic quantum fields over two-dimensional Minkowski space-time
satisfying all Wightman axioms with interaction and a positive, finite mass
gap at the lower end of the spectrum of the corresponding Hamiltonian.

Further properties of such models were discussed in,
e.g.,~{\cite{AHiPRoS2,AHiPRoS1,FroPa,Gie,HiKPS,Simon1974}}
and~{\cite{Albeverio2002,AlHKa,AFeHKL,AlLi,Liang}}. The case $| \sigma | <
\sqrt{8 \pi}$ was discussed by~S.~Kusuoka~in~{\cite{Ku}} and the relation with
independent work on multiplicative chaos by J. P. Kahane~{\cite{Ka}} (see also
below) was pointed out by H. Sato. In~{\cite{AHK2}} and
also~{\cite{Simon1974}} estimates in $L^p$ (for the interaction in a bounded
region) were given for $| \sigma | < \sqrt{\frac{4 \pi}{p - 1}}$. The Gaussian
character of the model for $| \sigma |$ sufficiently large or for any $\sigma$
if the Euclidean space dimension $d$ satisfies $d \geqslant 3$ \ has been
pointed out in~{\cite{AHK2}} and~{\cite{AGaHK}}. The relevance of the $\exp
(\sigma \varphi)_2$ model for Polyakov string theory has been discussed
in~{\cite{AlKrPaSc1986,AlKrPaSc1992,AJPS}} ($| \sigma | < \sqrt{4 \pi}$
corresponding to the embedding dimension $D < 13$) and rediscovered more
recently in connection with topics like Liouville model quantum gravity and
multiplicative chaos. Let us mention in this
connection~{\cite{AnKa,VaRhShe,RhVa}} see also
{\cite{BeShXi2014,DuSh2019,Gottschalk2013,Sha2016}} (for the literature for
other approaches to Liouville type models not directly connected with
probabilistic methods see, e.g.~{\cite{Bahns2018}}).

Finally a connection of the $\exp (\sigma \varphi)_2$ model with the
irreducibility of a unitary representation of the group of mappings of a
manifold into a compact Lie group has been pointed out and studied
in~{\cite{AHKMT}}, see also the recent work~{\cite{AGoVe}}.

\

The $\exp (\sigma \varphi)_2$ model is also closely connected with the $\sin
(\sigma \varphi)_2$ model (``Sine-Gordon equation'') as discussed
in~{\cite{AHK3,FrSei}}. In the latter reference all Wightman axioms are proved
again in the region $| \sigma | < \sqrt{4 \pi}$.

Both in the $\exp (\sigma \varphi)_2$ model and the $\sin (\sigma \varphi)_2$
model a replacement of the $\exp$ (respectively $\sin$) function by a Wick
ordered version is needed (see~{\cite{AHaRu}}) but also suffices for the
existence and non-triviality of the model for $| \sigma |$ up to $\sqrt{4
\pi}$. For larger values of $| \sigma |$, in the case of the $\sin (\sigma
\varphi)_2$ model, up to $\sqrt{8 \pi}$, further renormalisation by
counter-terms is required (see~{\cite{GaNi}}).

\

The study of the stochastic quantization equation associated with the latter
class of models has been initiated in~{\cite{AHaRu}} (where strong solutions
have been discussed and the necessity of renormalisation has been pointed
out). In the case where $\mathbb{R}^2$ (or $\mathbb{T}^2$) is replaced by
$\mathbb{R}$ (respectively $\mathbb{T}$), i.e. for the model $\exp (\sigma
\varphi)_1$, a deeper analysis is possible and has been pursued
in~{\cite{AKaRo}}, where the existence of solution and the strong uniqueness
of the invariant measure are proven (for the corresponding stochastic
quantization of $P (\varphi)_1$ see~{\cite{Iwata1987}}).

The case of the SQE in $d = 2$ and with a regularised noise on the torus (and
corresponding changes in the coefficients of the stochastic quantization
equation needed to keep the same invariant measure) has been discussed
in~{\cite{Mi}} (see also~{\cite{AKaMiRo}} for further development), where
existence of solutions was proven using essentially the properties of the
$\exp (\sigma \varphi)_2$ model established in~{\cite{Krohn1974}} and methods
of~{\cite{DaPraDe}}. Uniqueness problems are also discussed
in~{\cite{AKaMiRo}} in conjunction with an approach using Dirichlet forms.

Hairer and Shen~{\cite{HaiShe}} introduced powerful methods to handle the
dynamical $\sin (\sigma \varphi)_2$ model on $\mathbb{T}^2$, and via
regularity structures local existence is shown up to $\sigma^2 < \frac{16}{3}
\pi$. More recently these results has been extended in~{\cite{ChaHaiShe}} to
cover all the subcritical regime $\sigma^2 < 8 \pi$.

\

Concerning the exponential interaction, the work of
Garban~{\cite{Garban2018}} appeared recently in which the author studies the
SQE on the torus $\mathbb{T}^2$ and on the sphere $\mathbb{S}^2$. After
subtracting the solution to the linear equation he obtains an SPDE driven by a
multi-fractal and intermitted multiplicative chaos. When $| \sigma | < \left(
4 - 2 \sqrt{3} \right) \sqrt{\pi} \left( < \sqrt{4 \pi} \right)$ (a regime
correspondingly called Da Prato-Debussche phase) he shows the existence of a
strong solution and the convergence of the solution to the equation with
regularised noise to the singular one both on the torus $\mathbb{T}^2$ and on
the sphere $\mathbb{S}^2$ (see Theorem~1.7 and Theorem~1.9
in~{\cite{Garban2018}}). The method used is based on the Besov regularity of
the Gaussian multiplicative chaos (related to the theory of~{\cite{Ka}}). The
paper~{\cite{Garban2018}} also refers to other interesting relations to
quantum gravity, conformal fields theory, strings theory, multiplicative chaos
and exciting new results on random measures. A result on existence and
uniqueness for the equation (without the proof of the convergence of the
solution to the equation with regularised noise to the singular one) is also
proved (Theorem 1.11 in~{\cite{Garban2018}}) for $| \sigma | < \left( 4 - 2
\sqrt{2} \right) \sqrt{\pi} \left( < \sqrt{4 \pi} \right)$. Moreover a
comparison with the SQE for the Sine-Gordon model is provided (cfr. Section 7
in~{\cite{Garban2018}}).

\

Let us briefly review and comment the main results obtained in this paper. In
Section~\ref{section_discrete} we study the case of the following elliptic
SPDE (equation~{\eqref{eq:main1}} in Section~\ref{section_discrete})
\[ (- \Delta_x + m^2 + \mathfrak{L}) (\phi) (x, z) + \mathcal{A}^2 [f (x') g
   (z') V' (\phi (x', z'))] (x, z) = \xi^A (x, z) \]
on $\mathbb{R}^2 \times M$, with $M$ a $d$-dimensional manifold with or
without boundary, equipped with a positive Radon measure $\mathd z$ and a
real-valued non-linearity given by the derivative of a positive smooth
function $V$ defined on the real line growing at most exponentially at
infinity. We use the shorthand $f (x') g (z') V' (\phi (x', z'))$ with free
variables $(x', z')$ to mean the function $(x', z') \mapsto f (x') g (z') V'
(\phi (x', z'))$. Additional assumptions are as follows. The function $V$ is
taken to the convex or ``quasi convex'' (see Hypothesis~QC below).
$\mathfrak{L}$ and $\mathcal{A}$ are two positive self-adjoint operators
acting on $L^2 (M)$, $\mathcal{A}$ is bounded, $\mathfrak{L}$ is
not-necessarily bounded but has a dense domain in $L^2 (M)$, $\mathcal{A}$ and
$\mathfrak{L}$ commute and both have completely discrete spectrum with some
additional hypotheses on the eigenfunctions and the eigenvalues of
$\mathcal{A}$. The noise $\xi^A$ is supposed to be Gaussian with mean zero and
covariance corresponding to a regularisation given by $\mathcal{A}$. Finally
$f$ is a $C^2$ real-valued cut-off function on $\mathbb{R}^2$, decaying
exponentially at infinity and $g$ is a smooth real-valued cut-off with compact
support on $M$.

We introduce finite dimensional projections relative of~{\eqref{eq:main1}}
and prove the existence of a weak solution $\nu$ satisfying a form of the
dimensional reduction principle described in our previous
paper~{\cite{Albeverio2018elliptic}}, extended to the $d$ dimension case with
the presence of a non-trivial operator $\mathfrak{L}$. See
Theorem~\ref{theorem_main} for a precise statement.

Next we prove extensions of these results in two directions: in
Section~\ref{subsection_extension1} we consider the case $M =\mathbb{R}^d$,
$\mathfrak{L} = - \Delta_z$ (cfr. Theorem~\ref{theorem_extension1}), in
Section~\ref{subsection_extension2} we consider the problem of the removal of
the cutoffs $f, g$ in the case where $V$ is a convex function (cfr.
Theorem~\ref{theorem_extension2}).

\

The first part of the present paper contains the first example of rigorous
dimensional reduction for elliptic SPDE in $d + 2$ dimension with $d > 0$.
Indeed, the other only paper which to our knowledge study dimensional
reduction from the rigorous point of view is~{\cite{Klein1984}} whose main
result is a theorem about the dimensional reduction of an integral in a space
of functions in $d + 2$ variables to a Gibbs type measure in a space of
functions in $d$ variables. In particular the authors do not attack the
problem of the relation with the elliptic SPDE. Furthermore
in~{\cite{Klein1984}} only polynomial potentials $V$ are considered and the
regularisation of the noise is chosen to be compactly supported in Fourier
space. Ours Hypothesis~QC on the potential $V$ and Hypotheses H$\mathcal{A}$
and H$\mathcal{A} 1$ for the regularisation of the noise considered here are
quite more general. Let us stress that the results obtained under these more
general hypotheses uses in essential way the SPDE formulation of dimensional
reduction.

\

In Section~\ref{section_exponential} we extend our result to singular SPDEs.
In particular the elliptic stochastic quantization of the exponential
interaction is proven, where $V$ is a suitable renormalised exponential
function $\exp (\alpha \phi - \infty)$, for $| \alpha | < 4 \sqrt{2} \pi$
(corresponding to the reduced index $| \sigma | < \sqrt{8 \pi}$), and the
white noise is unregularized. Existence and uniqueness is first proven for
suitable regularised model using the results of the previous sections (cfr.
Theorem~\ref{theorem_regularexponential}). Subsequently the removal of the
regularisation in the noise is achieved in the Besov space $B^s_{p, p, \ell}
(\mathbb{R}^4)$, where $- 1 < s < 0$ and $1 < p \leqslant 2$ are suitable
constants depending only on $\alpha$. Using the existence, uniqueness and
convergence results for the elliptic SPDEs over $\mathbb{R}^4$ proven in
Section~\ref{subsection_extension2}, the existence and uniqueness results for
the elliptic stochastic quantization equation for exponential interaction
follows in Theorem~\ref{theorem_exponentialmain}. The dimensional reduction
result is then obtained (cfr. Theorem~\ref{theorem_dimensionalreduction1})
with uniqueness for all $| \alpha | < 4 \sqrt{2} \pi$. Our last result in this
section, Theorem~\ref{theorem_dimensionalreduction2} shows the convergence and
uniqueness of the equation when the spatial cut-off $g$ is removed. As an easy
corollary one obtains the full Euclidean invariance of the law of the
solutions.

Finally in Section~\ref{section_power} we prove the existence for singular
elliptic SPDE with Wick power non linearity in $d + 2 = 4$ dimension.
Furthermore we prove a dimensional reduction principle for some weak solutions
to SPDEs of this form providing the first example of elliptic stochastic
quantization for the polynomial interaction model $P (\varphi)_2$. Let us
stress that for polynomial interactions the problem of uniqueness of weak
solutions to the elliptic SPDE is open and this prevents a more detailed
analysis of the dimensional reduction.

\

The results contained in this second part of the paper gives the first
example of dimensional reduction for singular elliptic SPDE as originally
formulated by Parisi and Sourlas in~{\cite{parisi_random_1979}}. In particular
with respect to~{\cite{GH18}}, where the existence for singular elliptic SPDEs
with polynomial type non-linearity is proven for both $d + 2 = 4$ and, with
cubic type non-linearity, $d + 2 = 5$ dimension, in
Section~\ref{section_power} we prove, at least in the $d + 2 = 4$ dimensional
case, they can be used as stochastic quantization equation for quantum field
theory. However many problems still remains open wrt. the removal of the
spatial cutoffs.

\

Let us stress that, in Section~\ref{section_exponential}, we prove existence
and uniqueness of solutions for both elliptic and parabolic SPDE and the
relation with the Liouville measure in the full subcritical regime $| \sigma |
< \sqrt{8 \pi}$. The previous best result to our knowledge is the
paper~{\cite{Garban2018}} which considers only the case where $| \sigma | <
\left( 4 - 2 \sqrt{2} \right) \sqrt{\pi} \left( < \sqrt{4 \pi} \right)$ and
for the parabolic setting. This result is achieved using in an essential way
two properties of the exponential model: the fact that the Wick exponential of
a distribution is a positive measure (this fact is already exploited
in~{\cite{Garban2018}}) and the multifractality of Wick exponentials (see
Lemma~\ref{lemma_garban} for a precise formulation of this property). In
particular we prove the existence and uniqueness for the elliptic SPDE with
Wick exponential non-linearity using the space $B^{- s}_{p, p, \ell}
(\mathbb{R}^4)$, instead of $B^{- s'}_{\infty, \infty, \ell} (\mathbb{R}^4)$,
and only the Da Prato-Debussche trick. This is possible since for any exponent
$| \sigma | < \sqrt{8 \pi}$ the noise lives in inside $B^{- s}_{p, p, \ell}
(\mathbb{R}^4)$ for some $s < 1$ and $1 < p \leqslant 2$, instead for $|
\sigma | \geqslant \left( 4 - 2 \sqrt{2} \right) \sqrt{\pi} \left( < \sqrt{4
\pi} \right)$ the Wick power of the noise is in $B^{- s}_{\infty, \infty,
\ell} (\mathbb{R}^4)$ with $s > 2$. The other novelty is the fact that we are
able to solve the equation in the full space and we are easily able to prove
the Euclidean invariance of the law.

\

\tmtextbf{Acknowledgments.} The first named author is grateful to HCM of the
University of Bonn, to Stefania Ugolini (Universit{\`a} degli Studi di Milano)
and to the Department of Mathematics of Universit{\`a} degli Studi di Milano
for financial support. The authors are supported by the German Research
Foundation (DFG) via CRC 1060.

\section{Dimensional reduction with regularised noise}\label{section_discrete}

\subsection{Discrete spectrum}\label{section_discrete1}

Consider the following elliptic SPDE
\begin{equation}
  (- \Delta_x + m^2 + \mathfrak{L}) (\phi) (x, z) + \mathcal{A}^2 [f (x') g
  (z') \partial V (\phi (x', z'))] (x, z) = \xi^A (x, z)  \label{eq:main1}
\end{equation}
where $x \in \mathbb{R}^2$ and $z \in M$ (where $M$ is a $d$ dimensional
manifold with or without boundary, equipped with a Radon positive measure
$\mathd z$) and $V$ is a smooth function on $\mathbb{R}$ growing at most
exponentially at infinity, $\partial V$ its gradient and $\phi : \mathbb{R}^2
\times M \rightarrow \mathbb{R}$ a scalar random field.

We now list the hypotheses and assumptions on the various elements of
equation~{\eqref{eq:main1}} and of its generalisation to the vector case,
where $V$ is replaced by $V_n$ being defined on $\mathbb{R}^n$ (when
$\mathbb{R}^n =\mathbb{R}$ hereafter instead of writing $V_1$ we simply write
$V$).

{\descriptionparagraphs{\item[Hypothesis~C.] The potential $V_n :
\mathbb{R}^n \rightarrow \mathbb{R}$ is a positive smooth function such that
\[ y \in \mathbb{R}^n \mapsto V_n (y), \]
is strictly convex and it and its first and second partial derivatives grow at
most exponentially at infinity.

\item[Hypothesis~QC.] The potential $V_n : \mathbb{R}^n \rightarrow
\mathbb{R}$ is a positive smooth function, such that it and its first and
second partial derivatives grow at most exponentially at infinity and
additionally such that there exists a function $\mathfrak{H} : \mathbb{R}
\rightarrow \mathbb{R}$ with exponential growth at infinity such that we have
\[ - \langle \hat{n}, \partial V_n (y + r \hat{n}) \rangle \leqslant
   \mathfrak{H} (y), \qquad \text{$\hat{n} \in \mathbb{S}^{n - 1}, y \in
   \mathbb{R}^n$ and $r \in \mathbb{R}_+$}, \]
with $\mathbb{S}^{n - 1}$ is the $n - 1$ dimensional sphere.

\item[Hypothesis~H$f$.] The function $f : \mathbb{R}^2 \rightarrow \mathbb{R}$
has at least $C^2$ smoothness and in addition satisfies $f' \leqslant 0$, it
decays exponentially at infinity and fulfils $\Delta (f) \leqslant b^2 f$ for
$b^2 \ll m^2$ (some examples of such functions are given
in~{\cite{Klein1984}}).

\item[Hypothesis~H$g$.] The function $g : M \rightarrow \mathbb{R}$ is smooth
and with compact support.

\item[Hypothesis~H$\mathfrak{L}$.] The operator $\mathfrak{L}$ is a closed
operator (possibly unbounded) defined on the vector space $D_{\mathfrak{L}}
\subset L^2 (M)$ (where $L^2 (M)$ is the space of measurable $L^2$ functions
defined on $M$ with respect the measure $\mathd z$ on $M$), which is dense in
$L^2 (M)$. The range of $\mathfrak{L}$ is a subset of $L^2 (M)$. The operator
$\mathfrak{L}$ is positive and self-adjoint and has a completely discrete
spectrum. Furthermore we suppose that there exists at least an orthonormal
basis $H_1, \ldots, H_k, \ldots$ in $L^2 (M)$ composed by eigenfunctions of
$\mathfrak{L}$ which are $C^0_w (M)$ functions (where $C^0_w (M)$ is the space
of continuous functions $h$ defined on $M$ such that $\| h (z) w (z)
\|_{\infty} < + \infty$ where $w$ is a given positive weight continuous
function). Finally we denote by $\lambda_1 \leqslant \lambda_2 \leqslant
\cdots \leqslant \lambda_k \leqslant \cdots$ the eigenvalues of $\mathfrak{L}$
corresponding to the eigenfunctions $H_1, H_2, \ldots, H_k, \ldots_{}$.

\item[Hypothesis~H$\mathcal{A}$.] The operator $\mathcal{A} : L^2 (M)
\rightarrow L^2 (M)$ is a linear, injective, continuous, positive self-adjoint
operator commuting with $\mathfrak{L}$. If we denote by $\sigma_1, \ldots,
\sigma_k, \ldots$ the eigenvalues of $\mathcal{A}$ corresponding to the basis
$H_1, \ldots, H_k, \ldots_{}$ we suppose that
\[ \sum_{k = 1}^{\infty} \sigma_k < + \infty \hspace{3em} \sum_{k =
   1}^{\infty} \sigma_k \| H_k \|_{\infty, w}^2 < + \infty, \]
where $\| \cdot \|_{\infty, w}$ denotes the $L^{\infty}$-norm in the weighted
space $C^0_w (M)$.

\item[Hypothesis~H$\xi$.] The noise $\xi^A$ is Gaussian with mean zero and a
covariance corresponding to a regularisation provided by $\mathcal{A}$, namely
such that if $h_1, h_2 \in C^{\infty}_0 (\mathbb{R}^2 \times M)$ we have
\[ \mathbb{E} [\langle h_1, \xi^A \rangle \langle h_2^{}, \xi^A \rangle] =
   \int_{\mathbb{R}^2 \times M} (I \otimes \mathcal{A}) (h_1) (x, z) (I
   \otimes \mathcal{A}) (h_2) (x, z) \mathd x \mathd z. \]}}

\begin{remark}
  When $V_1 = V : \mathbb{R} \rightarrow \mathbb{R}$, Hypothesis~QC can be
  reformulated as
  \[ \mp V' (y \pm r) \leqslant \mathfrak{H} (y), \]
  where $r \geqslant 0$.
\end{remark}

In the rest of this Subsection~\ref{section_discrete1} we will always assume
Hypothesis~H$f$, H$g$, H$\mathfrak{L}$, H$\mathcal{A}$ and we specify the use
of Hypothesis~QC or C case by case. Moreover, hereafter we systematically use
the following abuse of notation: we denote by $\mathfrak{L}$ and $\mathcal{A}$
both the operators with the same name defined on the space $L^2 (M)$ and also
the operators $I \otimes \mathfrak{L}$ and $I \otimes \mathcal{A}$ defined on
$L^2 (\mathbb{R}^2 \times M)$.

\

We describe now the setting for equation~{\eqref{eq:main1}}. We assume that
$\xi^A$ is defined on an abstract Wiener space $(\mathcal{W}, \mathcal{H}^{},
\mu^{\mathcal{A}})$ where the Cameron-Martin space is
\[ \mathcal{H} = L^2 (\mathbb{R}^2) \otimes_H \mathcal{A} (L^2 (M)), \]
(here $\otimes_H$ is the natural tensor product of Hilbert spaces) equipped
with the following scalar product
\[ \langle h_1, h_2 \rangle = \int_{\mathbb{R}^2 \times M} \mathcal{A}^{- 1}
   (h_1) (x, z) \mathcal{A}^{- 1} (h_2) (x, z) \mathd x \mathd z, \qquad h_1,
   h_2 \in \mathcal{H} . \]
The Wiener space $\mathcal{W}$ is given by
\begin{eqnarray}
  \mathcal{W} & = & (- \Delta_x + 1) (C^0_{\ell} (\mathbb{R}^2) \cap W^{1 -,
  p}_{\ell} (\mathbb{R}^2)) \otimes_{\pi} \mathcal{A}^{1 / 2} (L^2 (M)),
  \nonumber
\end{eqnarray}
where $\ell \geqslant 0$, \ $C^0_{\ell} (\mathbb{R}^2)$ is the space of
continuous functions on $\mathbb{R}^2$ such that
\[ \| k \|_{\infty, \ell} = \sup_{x \in \mathbb{R}^2} (| k (x) | (1 + | x
   |)^{- \ell}), \]
$W^{1 -, p}_{\ell} (\mathbb{R}^2)$ denotes the Sobolev space of regularity $1
-$ and weight $r_{\ell} = (1 + | x |)^{- \ell}$, $\otimes_{\pi}$ denotes the
projective tensor product (see~{\cite{Ryan2002}}). The measure
$\mu^{\mathcal{A}}$ is the centred Gaussian measure with Cameron-Martin space
$\mathcal{H}$.

In order to prove that $(\mathcal{W}, \mathcal{H}, \mu^{\mathcal{A}})$ is
actually an abstract Wiener space it is sufficient to prove that $\mathbb{E}
[\| \xi^{\mathcal{A}} \|_{\mathcal{W}}] < + \infty$ where $\| \cdot
\|_{\mathcal{W}}$ denotes the norm of $\mathcal{W}$. We have that
\begin{equation}
  \xi^A (x, z) = \sum_{k = 1}^{\infty} \sigma_k \xi^k (x) H_k (z),
  \label{eq:serie1}
\end{equation}
where $x \in \mathbb{R}^2$, $z \in M$, $\xi^k$ are a sequence of independent
Gaussian white noises defined on $\mathbb{R}^2$ and $\sigma_k$ are defined as
in Hypothesis~H$\mathcal{A}$. We have that
\begin{eqnarray}
  \| \xi^A \|_{\mathcal{W}} & \leqslant & \sqrt{\left( \sum_{k = 1}^{\infty}
  \sigma_k \| H_k \|^2_{L^2 (M)} \right)} \cdot \sqrt{\sum^{\infty}_{k = 1}
  \sigma_k \| (- \Delta_x + 1)^{- 1} (\xi^k) \|^2_{C^0_{\ell} (\mathbb{R}^2)
  \cap W^{1 -, p}_{\ell} (\mathbb{R}^2)}} . \nonumber
\end{eqnarray}
On the other hand, by Hypothesis~H$\mathcal{A}$, we obtain
\[ \mathbb{E} \left[ \sum^{\infty}_{k = 1} \sigma_k \| (- \Delta_x + 1)^{- 1}
   (\xi^k) \|^2_{C^0_{\ell} (\mathbb{R}^2) \cap W^{1 -, p}_{\eta}
   (\mathbb{R}^2)} \right] \lesssim \sum_{k = 1}^{\infty} \sigma_k < + \infty,
\]
(where $\lesssim$ stands for $\leqslant$ modulo a multiplicative positive
constant) and $\sum_{k = 1}^{\infty} \sigma_k \| H_k \|^2_{L^2 (M)} = \sum_{k
= 1}^{\infty} \sigma_k < + \infty$, and thus $\mathbb{E} [\| \xi^A
\|_{\mathcal{W}}] < + \infty$. This proves that $(\mathcal{W}, \mathcal{H},
\mu^{\mathcal{A}})$ is an abstract Wiener space when $\xi^A : \mathcal{W}
\rightarrow \mathcal{W}$, defined as $\xi^A (w) = w$, has exactly
$\mu^{\mathcal{A}}$ as probability distribution.

For later use, it is important to note that $\mathcal{W}$ is continuously
embedded in
\[ \mathcal{W}' = (- \Delta_x + 1) (C^0_{\ell} (\mathbb{R}^2) \cap W^{1 -,
   p}_{\ell} (\mathbb{R}^2)) \otimes_{\pi} C^0_w (M), \]
since $\mathcal{A}^{1 / 2} (L^2 (M))$ is continuously embedded in $C^0_w (M)$
by Hypothesis~H$\mathcal{A}$.



Equation~{\eqref{eq:main1}} can be written as a problem defined on the
abstract Wiener space $(\mathcal{W}, \mathcal{H}, \mu^{\mathcal{A}})$. In
particular we define the map $U : \mathcal{W} \rightarrow \mathcal{H}$ by
\[ U (w) = \mathcal{A}^2 [f (x) g (z) V' (\mathcal{I} w)], \]
where $\mathcal{I} : \mathcal{W} \rightarrow (- \Delta_x + 1)^{- 1}
(\mathcal{W})$ is the linear operator given by $\mathcal{I} = (- \Delta_x +
m^2 + \mathfrak{L})^{- 1}$. We define the map $T : \mathcal{W} \rightarrow
\mathcal{W}$ as $T (w) = w + U (w)$.

\begin{definition}
  \label{definition1}It is clear that if S is any measurable map $S :
  \mathcal{W} \rightarrow \mathcal{W}$ satisfying $T \circ S =
  \tmop{id}_{\mathcal{W}}$ $\mu^{\mathcal{A}}$-almost surely we have that
  $\phi (w, x, z) = \mathcal{I} (S (w)) (x, z)$ is a (strong) solution to
  equation~{\eqref{eq:main1}}. Furthermore if $\nu$ is a probability law on
  $\mathcal{W}$ such that $T_{\asterisk} (\nu) = \mu^{\mathcal{A}}$ then a
  random variable $\phi$ taking values on $\mathcal{} \mathcal{I}
  (\mathcal{W})$ and having the same law as $\mathcal{I}_{\asterisk} (\nu)$ is
  a (weak) solution to equation~{\eqref{eq:main1}}. For this reason in the
  following we call a map $S : \mathcal{W} \rightarrow \mathcal{W}$ such that
  $T \circ S = \tmop{id}_{\mathcal{W}}$ $\mu^{\mathcal{A}}$-almost surely
  {\tmem{a strong solution to equation~{\eqref{eq:main1}}}} and we call a
  measure $\nu$ on $\mathcal{W}$ such that $T_{\asterisk} (\nu) =
  \mu^{\mathcal{A}}$ \tmtextit{{\tmem{weak solution to
  equation~{\eqref{eq:main1}}}}}.
\end{definition}

In order to solve equation~{\eqref{eq:main1}} we need to introduce an
approximation. Let $P_n$ be the orthogonal projection in $L^2 (M)$ onto the
finite dimensional subspace generated by $H_1, \ldots, H_n$. The restrictions
$\nobracket P_n |_{\mathcal{A} [L^2 (M)]}$ and $\nobracket P_n
|_{\mathcal{A}^{1 / 2} [L^2 (M)]}$ of $P_n$ on $\mathcal{A} [L^2 (M)]$ and
$\mathcal{A}^{1 / 2} [L^2 (M)]$ are the orthogonal projection on the subspace
generated by $H_1, \ldots, H_n$ too. This implies also that $I \otimes P_n$
(in the following denoted also simply by $P_n$) is a continuous linear
operator on $\mathcal{W}$.

Let $\phi_n$ be the solution to the following approximated equation
\begin{equation}
  (- \Delta_x + \mathfrak{L} + m^2) (\phi_n) (x, z) + P_n \mathcal{A}^2 [f
  (x') g (z') V' (\mathcal{} P_n (\phi_n) (x', z'))] (x, z) = \mathcal{} \xi^A
  (x, z) . \label{eq:reduced1}
\end{equation}
and let $U_n (w) = P_n (U (P_n (w)))$ and $T_n (w) = w + U_n (w)$ be the maps,
analogous to $U$ and $T$, related to equation~{\eqref{eq:reduced1}}. Using
these objects we can define weak and strong solutions to
equation~{\eqref{eq:reduced1}} as in Definition~\ref{definition1}.

\

Let us now study equation~{\eqref{eq:reduced1}}. We introduce the following
function $V_n : \mathbb{R}^n \rightarrow \mathbb{R}$
\[ V_n (y) = \int_M g (z) V \left( \sum_{k = 1}^n y^k H_k (z) \right) \mathd
   z. \]
Then, since $P_n$ commutes with $\mathfrak{L}$ and $\mathcal{A}$, being the
projection on a subset of common eigenfunctions of both $\mathfrak{L}$ and
$\mathcal{A}$, we have that $\phi_n$ is of the form
\[ \phi_n (x, z) = \mathcal{I}_{\mathfrak{}} \xi^A (x, z) + \bar{\phi}_n (x,
   z) = \mathcal{I} \xi^A (x, z) + \sum_{k = 1}^n \bar{\psi}_n^k (x) H_k (z),
\]
where $\bar{\psi}^k_n (x)$ solves the set of equations
\begin{equation}
  \begin{array}{lll}
    0 & = & (- \Delta_x + m^2 + \lambda_k) (\bar{\psi}^k_n) (x)\\
    &  & \quad + \sigma_k^2 \int_M g (z) \left. V' \middle( \sum_{j = 1}^n
    \mu_j \mathcal{I}_{\lambda_k} \xi^j (x) H_j (z) + \left. \sum_{j = 1}^n
    \bar{\psi}^j_n (x) \mathcal{} H_j (z) \right) H_k (z) \mathd z \right.\\
    & = & (- \Delta_x + m^2 + \lambda_k) (\bar{\psi}^k_n) (x) + \sigma_k^2
    [\partial_{y_k} (V_n) (\hat{\xi}_n (x) + \bar{\psi}_n)],
  \end{array} \label{eq:reduced2}
\end{equation}
with \ $\mathcal{I}_{\lambda_j} = (- \Delta_x + \lambda_j + m^2)^{- 1}$,
$\tilde{\xi}_n (x) = (\sigma_k \mathcal{I}_{\lambda_k} (\xi^k))_{k = 1,
\ldots, n} \in \mathbb{R}^n$ and we used the fact that $\mathcal{A}$ is
self-adjoint in $L^2 (M)$. If we denote by $\mathcal{A}_n$ the $n \times n$
diagonal matrix such that $\mathcal{A}^{_{i j}}_n = \sigma_j \delta^{i j}$,
$i, j = 1, \ldots, n$, and by $\mathfrak{L}_n$ the $n \times n$ matrix such
that $\mathfrak{L}_n^{i j} = \lambda_i \delta^{i j}$ we can write
equation~{\eqref{eq:reduced2}} in the following way
\begin{equation}
  (- \Delta_x + m^2 + \mathfrak{L}_n) (\psi_n) + \mathcal{A}^2_n \cdot
  \partial V_n (\psi_n) = \mathcal{A}_n \cdot \xi_n, \label{eq:reduced3}
\end{equation}
where $\nobracket \psi_n = (\bar{\psi}_n^k + \hat{\xi}_n^k) |_{k = 1, \ldots,
n} \nobracket .$ Equation~{\eqref{eq:reduced3}} is defined on the abstract
Wiener space
\[ \widehat{\mathcal{H}}_n \assign L^2 (\mathbb{R}^2 ; \mathbb{R}^n), \]
with its natural scalar product and natural norm given by $\langle h, g
\rangle = \sum_{i = 1}^n \frac{1}{\sigma_n} \int_{\mathbb{R}^2} h^i (x) g^i
(x) \mathd x$. Let $\widehat{\mathcal{W}}_n$ (in which
$\widehat{\mathcal{H}}_n$ is densely embedded) be the space
\[ \widehat{\mathcal{W}}_n \assign W^{p, - 1 -}_{\ell} (\mathbb{R}^2 ;
   \mathbb{R}^n) \cap (1 - \Delta_x) (C^0_{\ell} (\mathbb{R}^2 ;
   \mathbb{R}^n)) \]
and $\hat{\mu}_n$ is the law of the white noise $\hat{\xi}_n = (\sigma_1
\xi^1, \ldots, \sigma_n \xi^n)$. On $\widehat{\mathcal{W}}_n$ we can define
the corresponding maps $\hat{U}_n (w) (x) = f (x) \partial V^{\mathcal{A}}_n
(\mathcal{I}_{\mathfrak{L}_n} (w))$ and $\hat{T}_n (w) = w + \hat{U}_n (w)$
where
\[ V^{\mathcal{A}}_n (y) : = V_n (\sigma_1 y^1, \ldots, \sigma_n y^n), \]
$y = (y^1, \ldots, y^n)$, $y^i \in \mathbb{R}$, $i = 1, \ldots, n$ and
$\mathcal{I}_{\mathfrak{L}_n}$ is the matrix valued operator defined as
$\mathcal{I}_{\mathfrak{L}_n}^{i j} = \delta^{i j} \mathcal{I}_{\lambda_j}$.
Using $\hat{T}^n$ we can define the concept of strong and weak solutions to
equation~{\eqref{eq:reduced3}}.

\begin{remark}
  It is important to note that the relationship between $\psi_n$ and the
  strong solution $\hat{S}_n$ (defined using the map $\hat{T}_n$) is given by
  $\psi_n (w, x) = \mathcal{} \mathcal{I}_{\mathfrak{L}_n} (\hat{S}_n (w))
  (x)$, $x \in \mathbb{R}$. Furthermore if $\hat{\nu}_n$ is a weak solution to
  equation~{\eqref{eq:reduced3}} (i.e. such that $\hat{T}_{n, \asterisk}
  (\hat{\nu}_n) = \hat{\mu}_n$) we have that $\hat{\nu}^n$ is the law of $(-
  \Delta_x + m^2 + \mathfrak{L}_n) (\psi_n)$.
\end{remark}

In the following we shall denote by $\hat{\kappa}^n$ the probability law on
$\mathbb{R}^n$ such that
\[ \frac{\mathd \hat{\kappa}^n}{\mathd y} = \frac{\exp (- 2 \pi (| (m^2 +^{}
   \mathfrak{L}_n)^{1 / 2} (\mathcal{A}_n^{- 1} \cdot y) |^2 + V_n
   (y)))}{Z_{\hat{\kappa}^n}}, \]
where $Z_{\hat{\kappa}^n}$ is a suitable constant and $\mathd y$ is the
Lebesgue measure on $\mathbb{R}^n$. Finally we define $\tilde{\Upsilon}_{f,
n}$ as the random variable defined on $\widehat{\mathcal{W}}^n$ such that
\[ \tilde{\Upsilon}_{f, n} (w) : = \exp \left( 4 \int_{\mathbb{R}^2} f' (x)
   V_n (\mathcal{I}_{\mathfrak{L}_n} (w) (x)) \mathd x \right) = \exp \left( 4
   \int_{\mathbb{R}^2} f' (x) V_n^{} (\psi_n (w, x)) \mathd x \right) . \]
\begin{theorem}
  \label{theorem:reduced1}If $V^A_n$ satisfies Hypothesis~QC then there exists
  at least a weak solution $\hat{\nu}^n$ to equation~{\eqref{eq:reduced2}}
  such that for any bounded measurable function $F : \mathbb{R}^n \rightarrow
  \mathbb{R}$ we have
  \begin{equation}
    \int_{\widehat{\mathcal{W}}^n} F (\mathcal{} \mathcal{I}_{\mathfrak{L}_n}
    w (0)) \tilde{\Upsilon}_{f, n} (w) \mathd \hat{\nu}_n (w) = Z_{f, n}
    \int_{\mathbb{R}^n} F (y) \mathd \hat{\kappa}^n \label{eq:reduced4}
  \end{equation}
  where
  \[ Z_{f, n} : = \int_{\widehat{\mathcal{W}}^n} \tilde{\Upsilon}_{f, n} (w)
     \mathd \hat{\nu}_n (w) . \]
\end{theorem}

\begin{proof}
  The proof is given in~{\cite{Albeverio2018elliptic}} Theorem 1 for the case
  $\mathfrak{L}_n = 0$. The case with $\mathfrak{L}_n$ a generic positive
  diagonal matrix is a trivial extension.
\end{proof}

It is simple to prove that if the potential $V$ satisfies Hypothesis~QC the
potential $V_n^A$ satisfies Hypothesis~QC too, indeed the following
proposition holds.

\begin{proposition}
  If $V : \mathbb{R} \rightarrow \mathbb{R}$ satisfies Hypothesis~QC then
  $V_n^A$ satisfies Hypothesis~QC. Finally if $V$ is strictly convex (i.e. it
  satisfies Hypothesis~C) also $V_n^A$ and $V_n$ are convex.
\end{proposition}

\begin{proof}
  First of all we note that
  \[ \partial_{y^k} V^A_n (y) = \int_M g (z) V' \left( \mathcal{} \mathcal{A}
     \left( \sum_{j = 1}^n y^j H_j (z) \right) \right) \mathcal{A} (H_k) (z)
     \mathd z. \]
  This means that if $\tmmathbf{n} \in \mathbb{S}^n$ we have
  \begin{eqnarray}
    -\tmmathbf{n} \cdot \partial V_n^A (y + t\tmmathbf{n}) & = & - \sum_{k =
    1}^n n^k \partial_{y^k} (V_n^A) (y + t\tmmathbf{n}) \nonumber\\
    & = & - \int_M n_k g (z) V' \left( \mathcal{A} \sum_{j = 1}^n y^j H_j (z)
    + \mathcal{A} \sum_{j = 1}^n n^j H_j (z) \right) \mathcal{A} (H_k) (z)
    \mathd z \nonumber\\
    & \leqslant & \int_M g (z) \mathfrak{H} \left( \mathcal{A} \sum_{j = 1}^n
    y^j H_j (z) \right) \mathd z \nonumber\\
    & \lesssim & \mathfrak{H} \left( | y | \sum_{k = 1}^{+ \infty} \sigma_k
    \left\| H_k \mathbbm{1}_{g (z) \not{=} 0} \right\|_{\infty} \right)
    \nonumber
  \end{eqnarray}
  where we use that $g$ has compact support, the fact that $V$ satisfies
  Hypothesis~QC and the fact that $\mathfrak{H}$ is increasing. Since
  \[ \sum_{k = 1}^{+ \infty} \sigma_k \left\| H_k \mathbbm{1}_{g (z) \not{=}
     0} \right\|_{\infty} \leqslant \left\| \frac{1}{w} \mathbbm{1}_{g (z)
     \not{=} 0} \right\|_{\infty} \sum_{k = 1}^{+ \infty} \sigma_k \| H_k
     \|_{\infty, w} < + \infty \]
  by Hypothesis~H$\mathcal{A}$, the thesis is proved.
\end{proof}

Using the weak solution $\hat{\nu}^n$ to equation~{\eqref{eq:reduced3}}, given
by Theorem~\ref{theorem:reduced1}, we are able to construct a weak solution to
equation~{\eqref{eq:reduced1}} satisfying the dimensional reduction principle.

A weak solution $\nu_n$ to equation~{\eqref{eq:reduced1}} is of the form
$\nu_n = \tilde{\nu}_n \otimes \mu_n^{\mathcal{A}}$, where $\tilde{\nu}_n$ is
the law of
\[ (- \Delta_x + m^2 + \mathfrak{L}) P_n (\phi_n) (x, z) = \sum_{k = 1}^n (-
   \Delta_x + m^2 + \lambda_k) (\psi^k_n) (x) H_k (z), \]
on the subspace $\tmop{Im} (P_n) \subset \mathcal{W}$ and
$\mu_n^{\mathcal{A}}$ is the law of
\[ Q_n (\phi_n) = (- \Delta_x + m^2 + \mathfrak{L}) (I - P_n) (\phi_n) (x, z)
   = \sum_{k = n + 1}^{+ \infty} \mathcal{} \mu_k \xi^k (x) H_k (z), \]
which is the law of the Gaussian field on $\tmop{Im} (Q_n) \subset
\mathcal{W}$. Using the basis $H_1, \ldots, H_n$ we can identify the law
$\tilde{\nu}_n$ on $\tmop{Im} (P_n)$ with a probability measure $\hat{\nu}_n$
on $\widehat{\mathcal{W}}^n$. In this way it is evident that if the law
$\hat{\nu}_n$ satisfies the dimensional reduction principle on
$\widehat{\mathcal{W}}^n$ then the probability law $\nu_n = \hat{\nu}_n
\otimes \mu_n^{\mathcal{A}}$ satisfies the dimensional reduction principle on
$\mathcal{W}$, since it is the tensor product of two probability laws
satisfying the dimensional reduction principle.

More precisely, we consider the following abstract Wiener space $(\mathbb{W},
\mathbb{H}^{}, \mu^{\mathcal{A}, \mathfrak{L}})$ where
\begin{eqnarray}
  \mathbb{W} & = & \mathcal{A}^{1 / 2} (L^2 (M)), \nonumber\\
  \mathbb{H} & = & (\mathcal{\mathfrak{L}} + m^2)^{- 1 / 2} (\mathcal{A} (L^2
  (M))), \nonumber
\end{eqnarray}
(with $\mathcal{A}$ and $\mathfrak{L}$ as in~{\eqref{eq:main1}}) and
$\mathbb{W}$ is equipped with its natural norm while on $\mathbb{H}$ we define
the following scalar product
\begin{equation}
  \langle h_1, h_2 \rangle_{\mathbb{H}} = \frac{1}{4 \pi} \int_M (\mathfrak{L}
  + m^2)^{1 / 2} (\mathcal{A} (h_1)) (z) (\mathfrak{L} + m^2)^{1 / 2}
  (\mathcal{A} (h_2)) (z) \mathd z. \label{eq:covariance1}
\end{equation}
Thus $\mu^{\mathcal{A}, \mathfrak{L}}$ is the centred Gaussian measure on
$\mathbb{W}$ with Cameron-Martin space given by $\mathbb{H}$. We now denote by
$\kappa_n$ the measure on $\mathbb{W}$ such that
\begin{equation}
  \frac{\mathd \kappa_n}{\mathd \mu^{\mathcal{A} \mathfrak{, L}}} (\omega) =
  \left. \exp \left( - 4 \pi \int_M g (z) V (P_n (\mathcal{} \nobracket
  \omega) \nobracket (z)) \mathd z \right) \middle/ Z_{\kappa_n} \right.
  \label{eq:kappa1}
\end{equation}
where $\omega \in \mathbb{W}$ and $Z_{\kappa_n}$ is a renormalisation
constant.

\

We introduce also the following random variable
\[ \Upsilon_{f, n} (w) : = \exp \left( 4 \int_{\mathbb{R}^2 \times M} f' (x) g
   (z) V (\mathcal{} P_n (\mathcal{I} w) (x, z)) \mathd x \mathd z \right) \]
where $w \in \mathcal{W}$. If we use the previous identification of $\tmop{Im}
(P_n)$ with $\widehat{\mathcal{W}}^n$ we have that
\[ \tilde{\Upsilon}_{f, n} (P_n (w)) = \Upsilon_{f, n} (w) . \]
Using the previous observations the next proposition trivially follows.

\begin{proposition}
  \label{proposition_reduced1}If $V$ satisfies Hypothesis~QC then there exists
  a weak solution $\nu_n$ to equation~{\eqref{eq:reduced1}} such that for any
  measurable function $F : \mathbb{W} \rightarrow \mathbb{R}$ we have that
  \begin{equation}
    \int_{\mathcal{W}} F (\mathcal{I} w (0, \cdot)) \Upsilon_{f, n} (w) \mathd
    \nu_n (w) = Z_{f, n} \int_{\mathbb{W}} F (\omega) \mathd \kappa_n (\omega)
    . \label{eq:reduced8}
  \end{equation}
\end{proposition}

The rest of the section concerns the generalisation of
Proposition~\ref{proposition_reduced1} to solutions to
equation~{\eqref{eq:main1}}. In particular we denote by $\kappa$ the
probability measure on $\mathbb{W}$ such that
\begin{equation}
  \frac{\mathd \kappa}{\mathd \mu^{\mathfrak{\mathcal{A}, L}}} (\omega) =
  \left. \exp \left( - 4 \pi \int_M g (z) V (\omega (z)) \mathd z \right)
  \middle/ Z_{\kappa} \right. \label{eq:kappa2},
\end{equation}
where $\omega \in \mathbb{W}$.

\begin{theorem}
  \label{theorem_main}Suppose that $V$ satisfies Hypothesis~QC then there
  exists at least a weak solution $\nu$ to equation~{\eqref{eq:reduced2}} such
  that for any Borel measurable function $F : \mathbb{W} \rightarrow
  \mathbb{R}$ we have
  \begin{equation}
    \int_{\mathcal{W}} F (\mathcal{} \mathcal{I} w (0, \cdot)) \Upsilon_f (w)
    \mathd \nu (w) = Z_f \int_{\mathbb{W}} F (\omega) \mathd \kappa (\omega)
    \label{eq:main2},
  \end{equation}
  where $\kappa$ is defined as in~{\eqref{eq:kappa1}} and
  \[ Z_f = \int_{\mathcal{W}} \Upsilon_f (w) \mathd \nu (w) . \]
\end{theorem}

\begin{remark}
  Since $\kappa_n$ converges weakly to $\kappa$, the right hand side of
  equation~{\eqref{eq:reduced8}} converges to the right hand side of
  equation~{\eqref{eq:main2}}, as $n \rightarrow + \infty$. 
\end{remark}

In order to prove Theorem~\ref{theorem_main} we prove the following lemmas. In
the space of functions $C^0_{\ell} (\mathbb{R}^2, \mathbb{R}^n)$ we denote by
$\| \cdot \|^{\mathcal{A}}_{\ell, k}$ the following norm
\begin{eqnarray}
  \| h \|^{\mathcal{A}}_{\ell, k} & = & \sup_{x \in \mathbb{R}^2}
  \sqrt{\sum_{j = 1}^n \sigma_j^{2 k} (h^j (x))^2 r_{\ell} (x)^2}, \nonumber
\end{eqnarray}
where $h \in C^0 (\mathbb{R}^2, \mathbb{R}^n)$ and $k, \ell \in \mathbb{R}$.

\begin{lemma}
  \label{lemma_reduced1}Let $\bar{\psi}_n$ be a solution to the
  equation~{\eqref{eq:reduced2}} then we have
  \begin{equation}
    \| \bar{\psi}_n \|^{\mathcal{A}}_{\ell, - 1} \lesssim \left\| \exp (\alpha
    \Xi_{\ell', n} r_{\ell'} (x) w (z)) \mathbbm{1}_{g (z) \not{=} 0} f (x)
    r_{\ell} (x) \right\|_{\infty} \label{eq:inequality1}
  \end{equation}
  \begin{equation}
    \begin{array}{l}
      \| (- \Delta_x + m^2 + \mathfrak{L}_n) (\bar{\psi}_n)
      \|^{\mathcal{A}}_{\ell, - 1}\\
      \qquad \lesssim \left\| \exp ((\alpha' (\Xi_{\ell', n} w (z) + \|
      \bar{\psi}_n \|^{\mathcal{A}}_{\ell', - 1}) r_{\ell'} (x \nobracket))
      \mathbbm{1}_{g (z) \not{=} 0} f (x) r_{\ell} (x) \right\|_{\infty}
    \end{array} \label{eq:inequality2}
  \end{equation}
  where $\Xi_{\ell', n} = \| P_n \mathcal{I} \xi^{\alpha} (x, z)
  \|_{\nobracket C^0_{\ell'} (\mathbb{R}^2 \nobracket) \otimes_{\pi} C^0_w
  (M)}$, $\alpha, \alpha'$ and the implicit constants do not depend on $n$.
\end{lemma}

\begin{proof}
  We write $\tilde{\psi}_n = \left( \frac{\bar{\psi}_n^k}{\mu_k} \right)_{k =
  1, \ldots, n}$. The equation for $\tilde{\psi}_n$ reads
  \[ (- \Delta_x + m^2 + \lambda_k) (\tilde{\psi}^k_n) (x) + \mu_k f (x)
     \int_M g (z) H_k (z) V' (P_n \mathcal{I}_{\mathfrak{}} \xi^A (x, z) +
     \bar{\phi}_n (x, z) \mathd z) = 0 \]
  where $\bar{\phi}_n (x, z)$ is related to $\tilde{\psi}_n$ by
  \[ \bar{\phi}_n (x, z) = \sum_{k = 1}^n \mu_k \tilde{\psi}_k (x) H_k (z) .
  \]
  Putting $\tilde{\Psi}_n^2 (x) = (1 + \theta | x |)^{- 2 \ell} \sum_{k = 1}^n
  (\tilde{\psi}^k_n (x))^2$, writing $r_{\ell, \theta} (x) = (1 + \theta | x
  |)^{- 2 \ell}$, for some $\theta > 0$, and denoting by $\bar{x}$ the maximum
  of $\tilde{\Psi}_n$ we have
  \begin{eqnarray}
    m^2 \tilde{\Psi}_n^2 (\bar{x}) & \leqslant & - \frac{1}{2} \Delta
    (\tilde{\Psi}_n^2) (\bar{x}) + m^2 \tilde{\Psi}_n^2 (\bar{x}) +
    \tilde{\psi}_n (\bar{x}) \cdot \mathfrak{L}_n (\tilde{\psi}_n (\bar{x}))
    \nonumber\\
    & \leqslant & - r_{\ell, \theta} \tilde{\psi}_n \cdot \Delta
    \tilde{\psi}_n - r_{\ell, \theta} | \nabla \tilde{\psi}_n |^2 - \left(
    \frac{| \nabla r_{\ell, \theta} |^2 + \Delta r_{\ell, \theta}}{r_{\ell,
    \theta}^2} \right) \tilde{\Psi}^2_n (\bar{x}) \nonumber\\
    &  & \quad + m^2 \tilde{\Psi}_n^2 (\bar{x}) + r_{\ell, \theta}
    \tilde{\psi}_n (\bar{x}) \cdot \mathfrak{L}_n (\tilde{\psi}_n (\bar{x}))
    \nonumber\\
    & \leqslant & - f (\bar{x}) r_{\ell, \theta} (x) \int_M g (z)
    \bar{\phi}_n (\bar{x}, z) V' (P_n \mathcal{I}_{\mathfrak{}} \xi^A
    (\bar{x}, z) + \bar{\phi}_n (\bar{x}, z)) \mathd z \nonumber\\
    &  & \quad - \left[ \frac{| \nabla r_{\ell, \theta} |^2 + \Delta r_{\ell,
    \theta}}{r_{\ell, \theta}^2} \right] \tilde{\Psi}^2_n (\bar{x}), 
    \label{eq:norm1}
  \end{eqnarray}
  where we used that
  \[ \nabla \left[ \sum_{k = 1}^n (\tilde{\psi}^k_n (\bar{x}))^2 \right] = -
     \frac{\nabla r_{\ell, \theta} (\bar{x})}{r_{\ell, \theta} (\bar{x})}
     \sum_{k = 1}^n (\tilde{\psi}^k_n (\bar{x}))^2 \]
  since $\bar{x}$ is a stationary point for $\tilde{\Psi}_n$. From
  equation~{\eqref{eq:norm1}} if, for any fixed $\ell$, we choose a $\theta$
  in such a way that
  \[ \left[ \frac{| \nabla r_{\ell, \theta} |^2 + \Delta r_{\ell,
     \theta}}{r_{\ell, \theta}^2} \right] \leqslant m^2 \]
  we obtain
  \[ \tilde{\Psi}_n^2 (\bar{x}) \lesssim \left\| \exp (\alpha \Xi_{\eta', n}
     r_{\ell'} w (z)) \mathbbm{1}_{g (z) \not{=} 0} \right\|_{\infty} \cdot
     \left\| \bar{\phi}_n (x, z) \mathbbm{1}_{g (z) \not{=} 0}
     \right\|_{\infty} \]
  where all the constants do not depend on $n$. On the other hand
  \begin{eqnarray*}
    | \bar{\phi}_n (x, z) | & \leqslant & \sum_{k = 1}^n | \bar{\psi}_n^k (x)
    | \cdot | H_k (z) |\\
    & \lesssim & \left\| \frac{1}{w (z)} \mathbbm{1}_{g (z) \not{=} 0}
    \right\|_{\infty} \tilde{\Psi}_n (\bar{x}) \sqrt{\sum_{k = 1}^{\infty}
    \mu_k \| H_k (z) w (z) \|_{\infty}^2},
  \end{eqnarray*}
  for all $x \in \mathbb{R}^2$ and $z \in M$. Using Hypothesis~H$\mathcal{A}$
  we deduce from this that
  \[ \tilde{\Psi}_n (\bar{x}) \lesssim \left\| \exp (\alpha \Xi_{\ell', n}
     r_{\ell'} (x) w (z)) \mathbbm{1}_{g (z) \not{=} 0} \right\|_{\infty} . \]
  Since $\| \bar{\psi}_n \|^{\mathcal{A}}_{\ell, - 1} \lesssim_{\theta}
  \tilde{\Psi}_n (\bar{x})$ inequality~{\eqref{eq:inequality1}} is proved.
  Inequality~{\eqref{eq:inequality2}} follows directly from
  inequality~{\eqref{eq:inequality1}} and equation~{\eqref{eq:reduced2}}.
\end{proof}

\begin{corollary}
  \label{corollary_reduced1}Under the hypotheses and notations of
  Lemma~\ref{lemma_reduced1} we have that there exists an increasing
  continuous function $K : \mathbb{R}_+ \rightarrow \mathbb{R}_+$ such that
  \[ \| (- \Delta_x + m^2 + \mathfrak{L}) (\bar{\phi}_n) \|_{\mathcal{H}}
     \leqslant K (\Xi_{\ell, n}), \]
  with $\Xi_{\ell, n}$ as in Lemma~\ref{lemma_reduced1}.
\end{corollary}

\begin{proof}
  The proof consists simply in noting that
  \[ \| (- \Delta_x + m^2 + \mathfrak{L}) (\bar{\phi}_n) \|_{\mathcal{H}}
     \lesssim \| (- \Delta_x + m^2 + \mathfrak{L}_n) (\bar{\psi}_n)
     \|^{\mathcal{A}}_{\ell, - 1}, \]
  indeed
  \begin{eqnarray*}
    \| (- \Delta_x + m^2 + \mathfrak{L}) (\bar{\phi}_n) \|_{\mathcal{H}}^2 & =
    & \int_{M \times \mathbb{R}^2} (- \Delta_x + m^2 + \mathfrak{L})
    (\mathcal{A}^{- 1} (\bar{\phi}_n) (x, z))^2 \mathd x \mathd z\\
    & = & \sum_{k = 1}^n \frac{1}{\sigma_k^2} \int_{\mathbb{R}^2} (- \Delta_x
    + m^2 + \lambda_k) (\bar{\psi}_n)^2 (x) \mathd x\\
    & \leqslant & (\| (- \Delta_x + m^2 + \mathfrak{L}_n) (\bar{\psi}_n)
    \|^{\mathcal{A}}_{\ell, - 1})^2 \int_{\mathbb{R}^2} (1 + | x |)^{2 \ell}
    \mathd x
  \end{eqnarray*}
  which is finite whenever $\ell < - 1$.
\end{proof}

An important consequence of Corollary~\ref{corollary_reduced1} is the
following lemma.

\begin{lemma}
  \label{lemma_reduced2}Let $K \subset \mathcal{W}$ be a compact set, then
  $\mathfrak{K} = \cup_{n \in \mathbb{N}} T^{- 1}_n (K)$ is precompact in
  $\mathcal{W}$.
\end{lemma}

\begin{proof}
  We write $\mathcal{T}_n (\hat{\phi}, w) = \hat{\phi} + U_n (\hat{\phi} + w)
  .$ If $w_{n, y} \in \mathcal{W}$ is solution to the equation $T_n  (w_{n,
  y}) = y$ then $\hat{\phi}_{n, y} = w_{n, y} - y$ will be solution to the
  equation $\mathcal{T}_n (\hat{\phi}_{n, y} + y) = 0$. This will implies that
  $\bar{\phi}_{n, y} = \mathcal{I} (\hat{\phi}_{n, y})$ is a solution to
  equation~{\eqref{eq:reduced2}} for the realization of white noise $\xi^A
  (y)$. By Corollary~\ref{corollary_reduced1} we have that $\| \hat{\phi}_{n,
  y} \|_{\mathcal{H}} \leqslant K (\Xi_{n, \ell, y})$. On the other hand we
  have the estimate
  \[ \sup_{n \in \mathbb{N}} (\Xi_{\ell, n, y}) = \sup_{n \in \mathbb{N}} \|
     P_n \mathcal{I} \xi^{\mathcal{A}} (y) \|_{\nobracket C^0_{\ell}
     (\mathbb{R}^2 \nobracket) \otimes_{\pi} C^0_w (M)} \lesssim \sup_{n \in
     \mathbb{N}} \| P_n \xi^{\mathcal{A}} (y) \|_{\mathcal{W}} \leqslant \|
     \xi^{\mathcal{A}} (y) \|_{\mathcal{W}}, \]
  which implies that
  \[ \sup_{n \in \mathbb{N}, y \in K} \| \hat{\phi}_{n, y} \|_{\mathcal{H}}
     \leqslant \sup_{n \in \mathbb{N}, y \in K} K (\Xi_{n, \ell, y}) \lesssim
     \sup_{y \in K} \| \xi^{\mathcal{A}} (y) \|_{\mathcal{W}} = C_K < + \infty
     . \]
  Letting $\tilde{K} = \{ w \in \mathcal{W}, \| w \| < C_K \}$ we have that
  $\tilde{K}$ is compact in $\mathcal{W}$ since the inclusion map $i :
  \mathcal{H} \rightarrow \mathcal{W}$ is compact. This fact implies that
  $\mathfrak{K} \subset K + \tilde{K}$ is precompact since it is contained in
  the sum of two compact sets.
\end{proof}

\begin{lemma}
  \label{lemma_reduced3}The family $(\nu_n)_n$ of measures such that $T_{n,
  \asterisk} (\nu_n) = \mu^{\mathcal{A}}$ is tight.
\end{lemma}

\begin{proof}
  Let $\tilde{K}$ be a compact set such that $\mu (\tilde{K}) \geqslant 1 -
  \epsilon$ for a fixed $0 < \epsilon < 1$, then, by
  Lemma~\ref{lemma_reduced2}, $\mathfrak{K} \assign \overline{\cup_{n \in
  \mathbb{N}} T^{- 1}_n (\tilde{K})}$ is a compact set in $\mathcal{W}$. This
  implies
  \[ \nu_n (\mathfrak{K}) \geqslant \nu_n (\cup_i T^{- 1}_i (\tilde{K}))
     \geqslant \nu_n (T^{- 1}_n (\tilde{K})) \geqslant \mu^{\mathcal{A}}
     (\tilde{K}) \geqslant 1 - \epsilon, \]
  for any $n \in \mathbb{N}$. Thus the sequence $\nu_n$ is tight and the lemma
  is proven. 
\end{proof}

\begin{lemma}
  \label{lemma_reduced4}Suppose that $(\nu_n)_n$ is a sequence of weak
  solutions to equation~{\eqref{eq:reduced1}} weakly converging to $\nu$, then
  $\nu$ is a solution to~{\eqref{eq:main1}}, i.e. $T_{\asterisk} (\nu) =
  \mu^{\mathcal{A}}$.
\end{lemma}

\begin{proof}
  Proving the claim is equivalent to prove that for any function $h :
  \mathcal{W} \rightarrow \mathbb{R}$ which is bounded with continuous and
  bounded Fr{\'e}chet derivative we have $\int h \circ T \mathd \nu = \int h
  \mathd \mu^{\mathcal{A}}$. In order to prove this we note that
  \[ \| h \circ T (w) - h \circ T_n (w) \|_{\mathcal{W}} \leqslant \| h
     \|_{C^1 (\mathcal{W})} \cdot \| U (w) - P_n U (P_n (w)) \|_{\mathcal{W}}
     . \]
  On the other hand we have
  \[ \| U (w) - P_n U (P_n (w)) \|_{\mathcal{W}} \leqslant \sqrt{\sum_{k =
     n}^{\infty} \sigma_k} \cdot \| U (w) \|_{\mathcal{H}} + \sup_{w \in [w,
     P_n (w)]} \| \nabla U (y) \|_{\mathcal{L} (\mathcal{W})} \| P_n (w) - w
     \|_{\mathcal{W}} . \]
  Finally we observe that
  \[ \| P_n (w) - w \|_{\mathcal{W}} = \| (I - P_n) (T_n (w)) \|_{\mathcal{W}}
     . \]
  Let $K_{\epsilon} \subset \mathcal{W}$ be a compact set such that $\nu_n
  (\mathcal{W} \backslash K_{\epsilon}), \nu (\mathcal{W} \backslash
  K_{\epsilon}) < \epsilon$ then we have
  \begin{eqnarray}
    \left| \int_{\mathcal{W}} h \circ T \mathd \nu - \int_{\mathcal{W}} h
    \mathd \mu \right| & \leqslant & \left| \int_{\mathcal{W}} h \circ T
    \mathd \nu - \int_{\mathcal{W}} h \circ T \mathd \nu_n \right| + \left|
    \int_{\mathcal{W}} (h \circ T - h \circ T_n) \mathd \nu_n \right|
    \nonumber\\
    &  & \quad + \left| \int_{\mathcal{W}} h \circ T_n \mathd \nu_n -
    \int_{\mathcal{W}} h \mathd \mu \right| \nonumber\\
    & \leqslant & \left| \int_{\mathcal{W}} h \circ T \mathd \nu -
    \int_{\mathcal{W}} h \circ T \mathd \nu_n \right| + \epsilon \| h \|_{C^0
    (\mathcal{W})} \nonumber\\
    &  & \quad + \left| \int_{K_{\epsilon}} (h \circ T - h \circ T_n) \mathd
    \nu_n \right| .  \label{eq:reduced5}
  \end{eqnarray}
  But
  \begin{eqnarray}
    \left| \int_{K_{\epsilon}} (h \circ T - h \circ T_n) \mathd \nu_n \right|
    & \lesssim_h & \sqrt{\sum_{k = n}^{\infty} \sigma_k} \cdot \sup_{w \in
    K_{\epsilon}} \| U (w) \|_{\mathcal{H}} + \sup_{w \in K_{\epsilon}} \|
    \nabla U (w) \|_{\mathcal{L} (\mathcal{W})} \nonumber\\
    &  & \quad \times \int_{\mathcal{W}} \| (I - P_n) (T_n (w))
    \|_{\mathcal{W}} \mathd \nu_n .  \label{eq:reduced6}
  \end{eqnarray}
  On the other hand
  \begin{equation}
    \int_{\mathcal{W}} \| (I - P_n) (T_n (w)) \|_{\mathcal{W}} \mathd \nu_n =
    \int_{\mathcal{W}} \| (I - P_n) (w) \|_{\mathcal{W}} \mathd
    \mu^{\mathcal{A}} \lesssim \sum_{k = n}^{\infty} \sigma_k
    \label{eq:reduced7}
  \end{equation}
  where we use that $\mathbb{E} [\| (- \Delta_x + 1)^{- 1} (\xi^k)
  \|^2_{C^0_{\ell} (\mathbb{R}^2) \cap W^{1 -, p}_{\ell} (\mathbb{R}^2)}]$ is
  equal to a finite constant independent on $k$. Since $\sum \sigma_k$ is
  convergent, we have that the right hand side of~{\eqref{eq:reduced7}}
  converges to 0 as $n \rightarrow + \infty$. Since $U (w), \nabla U (w)$ are
  continuous with respect to $w$, this implies that the right hand side of
  equation~{\eqref{eq:reduced6}}, converges to 0 as $n \rightarrow + \infty$.
  Exploiting this fact, the fact that $\nu_n$ weakly converges to $\nu$ and
  equation~{\eqref{eq:reduced5}} we obtain that $\left| \int_{\mathcal{W}} h
  \circ T \mathd \nu - \int_{\mathcal{W}} h \mathd \mu^{\mathcal{A}} \right|
  \leqslant \epsilon \| h \|_{C^0 (\mathcal{W})}$. Since $\epsilon > 0$ is
  arbitrary the lemma is proved.
\end{proof}

\begin{proof*}{Proof of Theorem~\ref{theorem_main}}
  Thanks to Lemma~\ref{lemma_reduced4} we have that the sequence of weak
  solutions $\nu_n$, satisfying the dimensional reduction principle (whose
  existence is proven by Proposition~\ref{proposition_reduced1}) converges
  weakly to a solution $\nu$ of the equation~{\eqref{eq:main1}}. If we are
  able to prove that $\Upsilon_{n, f} \nu_n \rightarrow \Upsilon_f \nu$
  weakly, then the theorem would follow.
  
  First of all we note that $\Upsilon_{n, f} (w) = \Upsilon_f (P_n (w))$ and
  that $\Upsilon_f$ is a bounded Fr{\'e}chet $C^1 (\mathcal{W})$ function. We
  want to prove that $\nabla \Upsilon_f (w)$ is bounded when $w$ is in a
  bounded set of $\mathcal{W}$. We have that
  \[ \nabla \Upsilon_f (w) [h] = - \Upsilon_f (w) \cdot \int_{\mathbb{R}^2
     \times M} f' (x) g (z) V' (\mathcal{I} (w) (x, z)) (\mathcal{I} h) (x, z)
     \mathd x \mathd z. \]
  In particular we have that
  \[ \| \nabla \Upsilon_f (w) \|_{\mathcal{W}^{\asterisk}} \lesssim
     \int_{\mathbb{R}^2 \times M} f' (x) g (z) \exp (\alpha \| \mathcal{I} (w)
     \|_{C^0_{\ell} (\mathbb{R}^2) \otimes_{\pi} C^0_{\omega} (M)} (1 + | x
     |)_{}^{\ell} \omega (z)) \mathd x \mathd z, \]
  for a suitable positive constant $\alpha$. On the other hand the linear map
  $\mathcal{I} : \mathcal{W} \rightarrow C^0_{\ell} (\mathbb{R}^2)
  \otimes_{\pi} C^0_{\omega} (M)$ is continuous which implies that if $B$ is a
  bounded set of $\mathcal{W}$ given by $\sup_{w \in B} \| \mathcal{I} (w)
  \|_{C^0_{\ell} (\mathbb{R}^2) \otimes_{\pi} C^0_{\omega} (M)} < + \infty$.
  This observation proves that $\sup_{w \in B} \| \nabla \Upsilon_f (w)
  \|_{\mathcal{W}^{\asterisk}} < + \infty$ for any bounded set $B \subset
  \mathcal{W}$.
  
  Let for any given $\epsilon > 0$, $K_{\epsilon} \subset \mathcal{W}$ be a
  compact set such that $\nu_n (\mathcal{W} \backslash K_{\epsilon}), \nu
  (\mathcal{W} \backslash K_{\epsilon}) < \epsilon$, then there exists a ball
  $B_{\epsilon} \subset \mathcal{W}$ such that $K_{\epsilon} \cup (\cup_{n \in
  \mathbb{N}} P_n (K_{\epsilon})) \subset B_{\epsilon}$. Let $F$ be a
  continuous and bounded function on $\mathcal{W}$ then
  \begin{eqnarray}
    \left| \int_{\mathcal{W}} F (w) \Upsilon_{f, n} \mathd \nu_n -
    \int_{\mathcal{W}} F (w) \Upsilon_f \mathd \nu \right| & \leqslant &
    \left| \int_{\mathcal{W}} F (w) (\Upsilon_f (P_n (w)) - \Upsilon_f (w))
    \mathd \nu_n \right| + \nonumber\\
    &  & \quad + \left| \int_{\mathcal{W}} F (w) \Upsilon_f \mathd \nu_n -
    \int_{\mathcal{W}} F (w) \Upsilon_f \mathd \nu \right| \nonumber\\
    & \lesssim & \epsilon + \| \nabla \Upsilon_f (w) \|_{C^1 (B)}
    \int_{\mathcal{W}} \| P_n (w) - w \|_{\mathcal{W}} \mathd \nu_n
    \nonumber\\
    &  & \quad + \left| \int_{\mathcal{W}} F (w) \Upsilon_f \mathd \nu_n -
    \int_{\mathcal{W}} F (w) \Upsilon_f \mathd \nu \right| \nonumber\\
    & \lesssim & \epsilon + \| \nabla \Upsilon_f (w) \|_{C^1 (B_{\epsilon})}
    \int_{\mathcal{W}} \| P_n (T_n (w)) - T_n (w) \|_{\mathcal{W}} \mathd
    \nu_n \nonumber\\
    &  & \quad + \left| \int_{\mathcal{W}} F (w) \Upsilon_f \mathd \nu_n -
    \int_{\mathcal{W}} F (w) \Upsilon_f \mathd \nu \right| \rightarrow
    \epsilon, \nonumber
  \end{eqnarray}
  as $n \rightarrow \infty$, where the constants implied in the symbol
  $\lesssim$ do not depend on $\epsilon$. For this reason since $\epsilon$ is
  arbitrary positive, $\int_{\mathcal{W}} F (w) \Upsilon_{f, n} \mathd \nu_n
  \rightarrow \int_{\mathcal{W}} F (w) \Upsilon_f \mathd \nu$ and the theorem
  is proven.
\end{proof*}

\subsection{Dimensional reduction in the full
space}\label{subsection_extension1}

In this section we consider the following equation
\begin{equation}
  (- \Delta_x - \Delta_z + m^2) (\phi) (x, z) + \mathcal{A}^2 [f (x') g (z')
  V' (\phi (x', z'))] (x, z) = \xi^A (x, z) \label{eq:extension1},
\end{equation}
where $x, x' \in \mathbb{R}^2$ and $z, z' \in M$. We consider the following
hypotheses on $\mathfrak{L}$ and $\mathcal{A}$ (in addition to the previous
H$f$ and H$g$ on the cut-offs):

{\descriptionparagraphs{\item[Hypothesis~H$\mathfrak{L}$1] $\mathfrak{L} = -
\Delta_z$ and $M =\mathbb{R}^d$

\item[Hypothesis~H$\mathcal{A}$1] $\mathcal{A} (h (z')) = \mathfrak{a}
\asterisk h$ where $\mathfrak{a} (z)$ is a $\mathcal{C}^{1 + d / 2 +}
(\mathbb{R}^2)$ H{\"o}lder continuous function with compact support.

\item[Hypothesis~H$\xi 1$] The noise $\xi^A$ is such that if $h_1, h_2 \in
C^{\infty}_0 (\mathbb{R}^2 \times \mathbb{R}^d)$ we have
\[ \mathbb{E} [\langle h_1, \xi^A \rangle \langle h_2^{}, \xi^A \rangle] =
   \int_{\mathbb{R}^2 \times \mathbb{R}^d} (I \otimes \mathcal{A}) (h_1) (x,
   z) (I \otimes \mathcal{A}) (h_2) (x, z) \mathd x \mathd z. \]
}}

In this subsection we assume Hypotheses~H$\mathfrak{L}$1,~H$\mathcal{A}$1
and~H$\xi$1 instead of the corresponding assumptions
H$\mathfrak{L}$,~H$\mathcal{A}$ and~H$\xi$ respectively we had before. The
assumptions~H$f$ and~H$g$ on the space cut-offs however remain.

\

We define the abstract Wiener space $(\mathcal{W}, \mathcal{H},
\mu^{\mathcal{A}})$ relative to equation~{\eqref{eq:extension1}} as follows
\begin{eqnarray}
  \mathcal{H} & = & L^2 (\mathbb{R}^2) \otimes_H \mathcal{A} (L^2
  (\mathbb{R}^d)) \nonumber\\
  \mathcal{W} & = & (- \Delta_x + 1) (C^0_{\ell} (\mathbb{R}^2) \cap W^{1 -,
  p}_{\ell} (\mathbb{R}^2)) \otimes_{\pi} C^0_{\ell'} (\mathbb{R}^d) \nonumber
\end{eqnarray}
where we suppose $\ell, \ell' > 0$. The maps $T$ and $U$, and so the concept
of weak and strong solution to equation~{\eqref{eq:extension1}}, are defined
in the same way as in the previous section. We shall prove the following
theorem.

\begin{theorem}
  \label{theorem_extension1}Suppose that $V$ satisfies Hypothesis~QC (in
  Section \ref{section_discrete1}), then there exists a weak solution $\nu$ to
  equation~{\eqref{eq:extension1}} such that for any Borel measurable function
  $F : \mathbb{W} \rightarrow \mathbb{R}$
  \[ \int_{\mathcal{W}} F (\mathcal{} \mathcal{I} w (0)) \Upsilon_f (w) \mathd
     \nu (w) = Z_f \int_{\mathbb{W}} F (\omega) \mathd \kappa (\omega), \]
  as in~{\eqref{eq:kappa2}}. 
\end{theorem}

It is clear that the conditions H$\mathfrak{L}$1, H$\mathcal{A}$1 on
$\mathfrak{L}, M, \mathcal{A}$ are incompatible with the previous
Hypotheses~H$\mathfrak{L}$ and~H$\mathcal{A}$. Thus
Theorem~\ref{theorem_extension1} does not follow directly from
Theorem~\ref{theorem_main}. The rest of this subsection will focus on the
proof of Theorem~\ref{theorem_extension1}.

\

In order to achieve the proof we first replace $M$ by the manifolds $M_R
=\mathbb{T}^d_R$, i.e. $M_R$ is a torus of radius $R > 0$. If $h$ is a
real-valued function on $M =\mathbb{R}^d$ with sufficient decay at infinity we
can easily define a function $h_R$ on $M_R$ in the following way
\[ h_R (z_R) = \sum_{j \in \mathbb{Z}^d} h (z_R + R j) \]
where $z_R \in [- R / 2, R / 2]^d$. There is also a standard way of defining a
function $h_r$ on $M_R$ given one on $M$, that is the following
\[ \tilde{h}_R (z_R) = h (z_R) . \]
If the function $h$ has support contained in $[- R / 2, R / 2]^d$ then
$\tilde{h}_R = h_R$.

On the other hand if we consider a function $k : M_R \rightarrow \mathbb{R}$
we can associated with it a function $k^p$ defined on all of $M$ in the
following way
\[ k^p (z) = k \left( - R + 2 R \left( z - \left\lfloor \frac{z}{2 R}
   \right\rfloor \right) \right), \]
$z \in M =\mathbb{R}^d$. In general we have that $\widetilde{(k^p)}_R = k$,
and so if $k$ has compact support contained in $(- R / 2, R / 2)^d$ we have $k
= (k^p)_R$. Furthermore if $\mathfrak{b}$ is a function with compact support
on $\mathbb{R}^d$ and $k$ is a function on $M_R$ we have $(\mathfrak{b}_R
\asterisk k)^p = \mathfrak{b} \asterisk k^p$. Furthermore it is important to
note that if $\mathfrak{G}_R$ is the Green function associated with the
operator $(- \Delta_x - \Delta_z + m^2)^{- 1}$ on $\mathbb{R}^2 \times M_R$
and $\mathcal{G}$ is the Green function of the same operator on $\mathbb{R}^2
\times M$ we have
\[ \mathcal{G}_R = \mathfrak{G}_R \]
We define the operator $\mathcal{A}_R$ on $L^2 (M_R)$ as $\mathcal{A}_R (k) =
\mathfrak{a}_R \asterisk k$. In this way we can define the abstract Wiener
space $(\mathcal{W}_R, \mathcal{H}_R, \mu^{\mathcal{A}_R}_R)$ as follows
\begin{eqnarray}
  \mathcal{H}_R & = & L^2 (\mathbb{R}^2) \otimes_H \mathcal{A}_R (L^2 (M_R))
  \nonumber\\
  \mathcal{W}_R & = & (- \Delta_x + 1) (C^0_{\ell} (\mathbb{R}^2) \cap W^{1 -,
  p}_{\ell} (\mathbb{R}^2)) \otimes_{\pi} \mathcal{A}_R^{1 / 2} (L^2 (M_R))
  \nonumber
\end{eqnarray}
and $\mu^{\mathcal{A}_R}_R$ is the law of the noise $\xi^A_R$ on
$\mathcal{W}_R$ with covariance $\mathcal{A}_R$.

Using the map $(\cdot)^p$ we can extend the noise $\xi^A_R$, defined on
$M_R$, to the noise $\xi^{A, p}_R$ defined on all $M$. This means that the law
$\mu^{\mathcal{A}_R}_R$ of $\xi^{\mathcal{A}}_R$ on $\mathcal{W}_R$, thanks to
the map $(\cdot)^p,$induces a Gaussian measure $\mu^{\mathcal{A}_R, p}_R$ on
$\mathcal{W}$, which is the measure of the noise $\xi^{A, p}_R$. It is simple
to prove that $\mu^{\mathcal{A}_R, p}_R$ weakly converges to
$\mu^{\mathcal{A}}$ when $R \rightarrow + \infty$ and the support of
$\mathfrak{a}$ is compact.

Let $U_R : \mathcal{W}_R \rightarrow \mathcal{H}_R$ be the function defined
by
\[ U_R (w_R) : = f (x) \mathcal{A}_R (g_R (z_R) V' (\mathfrak{G}_R \asterisk
   w_R)) = f (x) \mathcal{A}_R (g_R (z_R) V' (\mathcal{G}_R \asterisk w_R)) .
\]
We set $T_R (w_R) = w_R + U_R (w_R)_{}$. The map $U_R$ induces a map $U_R^p$
on $\mathcal{W}$ in the following way
\[ U_R^p (w) = f (x) \mathcal{A}_R (g_R^p (z_R) V' (\mathcal{G} \asterisk
   (\tilde{w}_R)^p)), \]
and $T^p_R (w) = w + U^p_R (w)$. Let $\nu_R$ be a probability law on
$\mathcal{W}_R$ such that $T_{R, \asterisk} (\nu_R) = \mu^{\mathcal{A}_R}_R$,
then it induces a probability law $\nu^p_R$ on $\mathcal{W}$ such that $T^p_R
(\nu_R^p) = \mu^{\mathcal{A}_R, p}_R$.

\begin{lemma}
  \label{lemma_extension1}Let $R_n \in \mathbb{R}_+$ be a sequence such that
  $R_n \rightarrow + \infty$, then the sequence $\nu_{R_n}^p$ is tight on
  $\mathcal{W}$ and if $\nu_{R_n}^p \rightarrow \nu$, as $n \rightarrow
  \infty$ and $R_n \rightarrow \infty$, then $T_{\asterisk} (\nu) =
  \mu^{\mathcal{A}}$.
\end{lemma}

\begin{proof}
  We note that the support of $\nu_{R_n}^p$ is contained in the set of $R_n$
  periodic distributions contained in $\mathcal{W}$ and furthermore the map
  $T_{R_n}^p$ sends $R_n$ periodic distributions into $R_n$ periodic
  distributions.
  
  Using the methods of Lemma~\ref{lemma_reduced1} and
  Corollary~\ref{corollary_reduced1}, it is simple to prove that if $y_n \in
  \mathcal{W}$ is an $R_n$ periodic distribution and if $w_{y_n} \in
  \mathcal{W}$ is an $R_n$ periodic distribution such that $w_{y_n} \in T^{p,
  - 1}_{R_n} (y_n)$ then
  \begin{equation}
    \| w_{y_n} - y_n \|_{\mathcal{H}} \leqslant K_{f, g} (\Xi_{\ell, \ell'}
    (y_n)), \label{eq:extension2}
  \end{equation}
  where
  \[ \Xi_{\ell, \ell'} (y_n) = \sup_{(x, z) \in \mathbb{R}^2 \times M} |
     \mathcal{G} \asterisk y_n (x, z) | \rho_{\ell} (x) \rho_{\ell'} (z), \]
  and $K_{f, g}$ is a positive increasing continuous function depending only
  on the functions $f$ and $g$. It is simple to prove that the map $\Xi_{\ell,
  \ell'}$ is continuous on $\mathcal{W}$. Using the fact that, as $n
  \rightarrow \infty$ and $R_n \rightarrow \infty$, $\mu^{\mathcal{A}_R,
  p}_{R_n}$ converges to $\mu^{\mathcal{A}}$ weakly and so the sequence
  $\mu^{\mathcal{A}_R, p}_{R_n}$ is tight, we can use the
  bound~{\eqref{eq:extension2}} and the same methods of
  Lemma~\ref{lemma_reduced3} to prove that $\nu_{R_n}^p$ is tight.
  
  Suppose that $\nu_{R_n}^p$ weakly converges to $\nu$, we want to prove that
  $T_{\asterisk} (\nu) = \mu^{\mathcal{A}}$. Let $F_{R_n}$ be a function of
  the form $F_{R_n} (w) = G (\langle h_{1,} w_{} \rangle, \ldots, \langle h_r,
  w \rangle)$, where $G : \mathbb{R}^r \rightarrow \mathbb{R}$ is \ a
  continuous and bounded function and $h_1, \ldots, h_r$ are smooth functions
  with support in $(- R / 2 + \mathfrak{r}, R / 2 - \mathfrak{r})^n$ and
  $\mathfrak{r} = \tmop{diam} (\tmop{supp} (\mathfrak{a}))$. For this kind of
  function we have that $F_{R_n} \circ T_{R_k}^p = F_{R_n} \circ T$ for $k
  \geqslant n$. From this observation we get
  \[ \int F_{R_n} \circ T \mathd \nu = \lim_k \int F_{R_n} \circ T \mathd
     \nu_{R_k}^p = \lim_k \int F_{R_n} \circ T_{R_k}^p \mathd \nu_{R_k}^p =
     \lim_k \int F_{R_n} \mathd \mu^{\mathcal{A}_{R_k}, p}_{R_k} = \int
     F_{R_n} \mathd \mu^{\mathcal{A}}, \]
  where the limit is taken for $k \rightarrow \infty$ and $R_k \rightarrow
  \infty$. Since the functions of the form $F_{R_n}$, for $n \in \mathbb{N}$,
  generate the space of all $\mathcal{W}$ Borel measurable functions the Lemma
  is proved.
\end{proof}

\begin{proof*}{Proof of Theorem~\ref{theorem_extension1}}
  By Theorem~\ref{theorem_main} there exists a probability law $\nu_{R_n}$ on
  $\mathcal{W}_{R_n}$ such that $T_{R_n, \asterisk} (\nu_{R_n}) =
  \mu^{\mathcal{A}_{R_n}}_{R_n}$. On the other hand this implies that
  $T^p_{R_n} (\nu_{R_n}^p) = \mu^{\mathcal{A}_{R_n}, p}_{R_n}$ and so, by
  Lemma~\ref{lemma_extension1}, there exists at last one probability measure
  $\nu$ on $\mathcal{W}$ such that $T_{\asterisk} (\nu) = \mu^{\mathcal{A}}$
  and $\nu^p_{R_n} \rightarrow \nu$ weakly, as $n \rightarrow \infty$.
  
  Then, using the notations of the proof of Lemma~\ref{lemma_extension1}, we
  have that for any continuous function $F_{R_n} : \mathbb{W} \rightarrow
  \mathbb{R}$
  \begin{equation}
    \int_{\mathcal{W}} F_{R_n} (\mathcal{G} \asterisk w (0, z)) \Upsilon_f (w)
    \mathd \nu^p_{R_k} (w) = \int_{\mathbb{W}} F_{R_n} (\omega) \mathd
    \kappa_{R_k}^p (\omega), \label{eq:extension3}
  \end{equation}
  where we used that
  \[ \int_{\mathcal{W}} F_{R_n} (\mathcal{G}_{R_k} \asterisk w_{R_k} (0, z))
     \Upsilon_f (w_{R_k}) \mathd \nu_{R_k} (w_{R_k}) = \int_{\mathcal{W}}
     F_{R_n} (\mathcal{G} \asterisk w (0, z)) \Upsilon_f (w) \mathd
     \nu^p_{R_k} (w), \]
  for $k \geqslant n$ and a similar equality for $\kappa_{R_k}$. Since
  $F_{R_k} (\mathcal{G} \asterisk \cdot)$ and $\Upsilon_f$ are continuous on
  $\mathcal{W}$, the left hand side of~{\eqref{eq:extension3}}, converges to
  $\int_{\mathcal{W}} F_{R_n} (\mathcal{G} \asterisk w (0, z)) \Upsilon_f (w)
  \mathd \nu (w)$ as $k \rightarrow + \infty$.
  
  Furthermore, since
  \[ \left. \frac{\mathd \kappa_{R_k}^p}{\mathd \mu^{\mathcal{A}_{R_k}
     \mathfrak{, L}, p}} = Z^{- 1}_{\kappa_{R_k}^p} \right. \left. \exp \left(
     - 4 \pi \int_{\mathbb{R}^2} g (z) V (\omega (z)) \mathd z \right) \right.
  \]
  and since $\mu^{\mathcal{A}_{R_k} \mathfrak{, L}, p}$ weakly converges to
  $\mu^{\mathcal{A}, \mathfrak{L}}$, we have that $\kappa^p_{R_k}$ weakly
  converges to $\kappa$, as $k, R_k \rightarrow \infty$. This proves that the
  right hand side of~{\eqref{eq:extension3}} converges to $\int_{\mathbb{W}}
  F_{R_n} (\omega) \mathd \kappa (\omega)$, as $k, R_k \rightarrow \infty$.
  Since the functions of the form $F_{R_n}$, for $n \in \mathbb{N}$, generate
  the space of $\mathbb{W}$ measurable functions, the theorem is proved.
\end{proof*}

\subsection{Cut-off removal with convex
potential}\label{subsection_extension2}

Hereafter we denote by $\omega_{\beta} (x)$ the function
\begin{equation}
  \omega_{\beta} (x) \assign \exp (- \beta \sqrt{(1 + | x |^2)}),
  \label{eq:exponentialweight}
\end{equation}
$\beta > 0$, $x \in \mathbb{R}^2$, and introduce the space
$\mathcal{W}_{\beta}$ in the following way
\[ \mathcal{W}_{\beta} : = (- \Delta + 1) C^0_{\exp \beta} (\mathbb{R}^2)
   \otimes_{\pi} \mathcal{A}^{1 / 2} (L^2 (M)) \]
where $C^0_{\exp \beta}$ is the space of continuous functions with respect to
the weighted $L^{\infty}$ norm
\[ \| g \|_{\infty, \exp \beta} \assign \sup_{x \in \mathbb{R}^2} |
   \omega_{\beta} (x) g (x) | . \]
In this section we want to prove the following theorem.

\begin{theorem}
  \label{theorem_extension2}Suppose that $V$ is a convex function, and suppose
  that $\mathcal{A}$, and $\mathfrak{L}$ satisfy Hypotheses~H$\mathcal{A}$,
  H$\mathfrak{L}$ and H$\xi$ (or Hypotheses~H$\mathcal{A}$1, H$\mathfrak{L}$1
  and H$\xi$1) then there exists a unique strong solution $\phi (x, z)$ to
  equation~{\eqref{eq:main1}} (or to equation~{\eqref{eq:extension1}}) with $f
  \equiv 1$ taking values on $\mathcal{W}_{\beta}$ (for a $\beta \leqslant
  \beta_0$ which depends only on $m^2$) such that for any $\mathbb{W}$
  measurable bounded function $F$ we have
  \begin{equation}
    \mathbb{E} [F (\phi (0, z))] = \int_{\mathbb{W}} F (\omega) \mathd \kappa
    (\omega) \label{eq:extension4}
  \end{equation}
  ($\kappa$ as in equation~{\eqref{eq:kappa2}}).
\end{theorem}

The proof is very similar to those of Theorem~\ref{theorem_main} and
Theorem~\ref{theorem_extension1}. For this reason we report here only the main
differences. First of all we need a replacement for
Proposition~\ref{proposition_reduced1}.

\begin{proposition}
  \label{proposition_extension1}Let $V$ (and so $V_n$) be a convex function
  then under assumptions H$\mathfrak{L}$, H$\mathcal{A}$ and H$\xi$ \ there
  exists a unique strong solution $\phi_n$ to equation~{\eqref{eq:reduced1}}
  such that for all $\mathbb{W}$ measurable bounded function $F$ we have
  \begin{equation}
    \mathbb{E} [F (\phi_n (0, z))] = \int_{\mathbb{W}} F (\omega) \mathd
    \kappa_n (\omega) \label{eq:reduced9}
  \end{equation}
  where $\kappa_n$ is given by expression~{\eqref{eq:kappa1}}.
\end{proposition}

\begin{proof}
  The proof is given in~{\cite{Albeverio2018elliptic}}, Theorem~2 in the case
  $\mathfrak{L}_n = 0$. The case considered here is a trivial extension. 
\end{proof}

In order to pass from equation~{\eqref{eq:reduced9}} to
equation~{\eqref{eq:extension4}} we need a generalisation of
Lemma~\ref{lemma_reduced1}. We denote by $\| \cdot \|_{\exp \beta,
\ell}^{\mathcal{A}}$ and $\| \cdot \|^{\mathcal{A}}_{U, k}$ the following
norms
\[ \begin{array}{lll}
     \| h \|^{\mathcal{A}}_{\exp \beta, k} & \assign & \sup_{x \in
     \mathbb{R}^2} \sqrt{\sum_{j = 1}^n \sigma_j^{2 k} (h^k (x))^2
     \omega_{\beta} (x)^2},\\
     \| h \|^{\mathcal{A}}_{U, k} & \assign & \sup_{x \in U} \sqrt{\sum_{j =
     1}^n \sigma_j^{2 k} (h^k (x))^2} .
   \end{array} \]


\begin{lemma}
  \label{lemma_extension2}There exists a number $\beta_0 > 0$ such that for
  any $0 < \beta \leqslant \beta_0$, and for any open bounded set $U \subset
  \mathbb{R}^2$ and under
  Hypotheses~C,~H$g$,~H$\mathcal{A}$,~H$\mathfrak{L}$,~H$\xi$ and $f \equiv 1$
  we have
  \begin{eqnarray}
    \| \bar{\psi}_n \|^{\mathcal{A}}_{\exp \beta, - 1} & \lesssim & \left\|
    \exp (\alpha \Xi_{\ell', n} \rho_{\ell'} (x) w (z)) \mathbbm{1}_{g (z)
    \not{=} 0} \omega_{\beta} (x) \right\|_{\infty} \label{eq:extension5} \\
    \| (- \Delta_x + m^2 + \mathfrak{L}_n) (\bar{\psi}_n) \|^{\mathcal{A}}_{U,
    - 1} & \lesssim & \| \exp (\alpha' (\Xi_{\ell', n} w (z) + \| \bar{\psi}_n
    \|^{\mathcal{A}}_{\exp \beta, - 1}) \omega_{\beta} (x)) \|_{C^0 (U)}
    \label{eq:extension6} 
  \end{eqnarray}
  uniformly in $n$ (where $\Xi_{\ell, n}$ is defined as in Lemma
  \ref{lemma_reduced1}).
\end{lemma}

\begin{proof}
  The proof is verbatim the same of Lemma~\ref{lemma_reduced1} where we
  replace the function $\rho_{\theta, \ell} (x) = (1 + \theta | x |)^{\ell}$
  by the function $\omega_{\beta}$, defined as in the beginning of the
  section, and we use the fact for $\beta$ small enough we have
  \[ \left( \frac{| \nabla \omega_{\beta} |^2 + \Delta
     \omega_{\beta}}{\omega_{\beta}^2} \right) < m^2 . \]
\end{proof}

The inequality~{\eqref{eq:extension6}} implies that, for any bounded open
subset $U$ of $\mathbb{R}^{2 + d}$ we have
\begin{equation}
  \| (- \Delta_x + m^2 + \mathfrak{L}) (\bar{\phi}_n) \|_{\mathcal{H}_U}
  \leqslant K_U (\Xi_{\ell, n}) \label{eq:extension7}
\end{equation}
where $\| \cdot \|_{\mathcal{H}_U}$ is the natural norm of the Hilbert space
$\mathcal{H}_U = L^2 (U) \otimes_H \mathcal{A} (L^2 (M))$, and $K_U$ is a
positive increasing continuous function depending only on $U$.
Inequality~{\eqref{eq:extension7}} guarantees us enough compactness to
generalise \ Lemma~\ref{lemma_reduced2}, Lemma~\ref{lemma_reduced3} and
Lemma~\ref{lemma_reduced4} in order to prove the existence of a weak solution
to equation~{\eqref{eq:main1}} satisfying the dimensional reduction principle
when $f \equiv 1$ and under Hypotheses~C, H$g$, H$\mathcal{A}$,
H$\mathfrak{L}$, H$\xi$. We can generalise the described result under
Hypotheses~C, H$g$, H$\mathcal{A}$1, H$\mathfrak{L}$1, H$\xi$1 using a similar
strategy.

It remains to prove the uniqueness of the solution to
equation~{\eqref{eq:main1}} (or~{\eqref{eq:extension1}}) when $V$ is convex.
This can be done using the methods of Lemma~\ref{lemma_reduced1} and
Lemma~\ref{lemma_extension2}. We give here only a sketch of the proof.

Let $\phi_1$ and $\phi_2$ be two strong solutions to
equation~{\eqref{eq:main1}}. It is simple to prove that $\phi_1 (x, \cdot) -
\phi_2 (x, \cdot) \in \mathcal{A} (L^2 (M))$, so the following function is
well defined
\[ \Psi_{\beta} (x)^2 = \omega_{2 \beta} (x) \int_M (\mathcal{A}^{- 1} (\phi_1
   (x, z') - \phi_2 (x, z')))^2 \mathd z, \]
where $\omega_{2 \beta}$ defined as in equation
{\eqref{eq:exponentialweight}}. The function $\Psi^2_{\beta}$ belongs to $C^2
(\mathbb{R}^2)$ \ and goes to zero at infinity. This means that the maximum is
realised in at least one point $\bar{x} \in \mathbb{R}^2 .$ Making some
computations similar to the ones of Lemma~\ref{lemma_reduced1} we obtain
\[ \begin{array}{lll}
     m^2 \Psi_{\beta} (\bar{x})^2 & \leqslant & - \int_M g (z) (\phi_1
     (\bar{x}, z) - \phi_2 (\bar{x}, z))_{} (V' (\phi_1 (\bar{x}, z)) - V'
     (\phi_2 (\bar{x}, z))) \mathd z\\
     &  & \quad - \left( \frac{| \nabla \omega_{\beta} |^2 + \Delta
     \omega_{\beta}}{\omega_{\beta}^2} \right) \Psi_{\beta} (\bar{x})^2 .
   \end{array} \]
On the other hand there exists a point $\theta_{\bar{x}, z} \in [\phi_1
(\bar{x}, z), \phi_2 (\bar{x}, z)]$ such that
\[ V' (\phi_1 (\bar{x}, z)) - V' (\phi_2 (\bar{x}, z)) \leqslant V''
   (\theta_{\bar{x}, z}) (\phi_1 (\bar{x}, z) - \phi_2 (\bar{x}, z))_{} . \]
Choosing $\beta > 0$ small enough we obtain
\[ \sup_{x \in \mathbb{R}^2} \Psi_{\beta} (x)^2 \lesssim - \int_M g (z) V''
   (\theta_{\bar{x}, z}) (\phi_1 (\bar{x}, z) - \phi_2 (\bar{x}, z))^2_{}
   \mathd z \leqslant 0, \]
since $V$ is convex. This implies that $\Psi_{\beta} (x) \equiv 0$ and so that
$\phi_1 (x, z) = \phi_2 (x, z)$.

\section{The exponential interaction on
$\mathbb{R}^2$}\label{section_exponential}

In this section we want to consider the following elliptic SPDE
\begin{equation}
  (- \Delta_x - \Delta_z + m^2) (\phi) + g (z) \alpha \exp (\alpha \phi -
  \infty) = \xi \label{eq:exponential1}
\end{equation}
where $x \in \mathbb{R}^2$, $z \in M =\mathbb{R}^2$ and $\xi = \xi (x, z)$ is
a standard Gaussian white noise on $\mathbb{R}^4$ $| \alpha | < 4 \sqrt{2}
\pi$, and where $- \infty$ means that the equation should be properly
renormalised. In order to give a meaning to the previous equation we formally
subtract from the solution $\phi$ the solution to the linear equation (i.e.
equation {\eqref{eq:exponential1}} with $g = 0$) which means that we consider
the equation for the unknown $\bar{\phi}$:
\begin{equation}
  (- \Delta_x - \Delta_z + m^2) (\bar{\phi}) + g (z) \alpha \exp (\alpha
  \bar{\phi}) \eta (x, z) = 0 \label{eq:exponential2},
\end{equation}
where $\eta (x, z) = \exp^{\diamond} (\mathcal{I} \xi)$ is a renormalised
version of the distribution $\exp (\alpha \mathcal{I} \xi - \infty)$, where
$\exp^{\diamond}$ denotes the Wick exponential of the Gaussian distribution
$\mathcal{I} \xi$. Hereafter we denote by $B^s_{p, q, \ell} (\mathbb{R}^{d +
2})$ the weighted Besov space of indices $1 \leqslant p \leqslant \infty$ and
$1 \leqslant q \leqslant \infty$ and weight given by $\rho_{\ell} (x, z) =
\left( \sqrt{| x |^2 + | z |^2} + 1 \right)^{- \ell}$
(see~{\cite{Triebel2006}}). It is well known that $\mathcal{I} \xi \in B^{-
\delta}_{p, p}$ for any $1 \leqslant p \leqslant \infty$ and $\delta > 0$ (see
e.g.~{\cite{GH18}}).

\

In the following we shall give a rigorous meaning to
equation~{\eqref{eq:exponential2}} (and so to
equation~{\eqref{eq:exponential1}}) when the exponent $| \alpha | < 4 \sqrt{2}
\pi$ and, when $\alpha \bar{\phi} \leqslant 0$, prove that there exists only
one solution to equation~{\eqref{eq:exponential1}}.

Furthermore we want to prove that dimensional reduction holds for the unique
solution to equation~{\eqref{eq:exponential1}}, namely that, if we consider
the measure $\kappa_g$ given by
\begin{equation}
  \frac{\mathd \kappa_g}{\mathd \mu} = \exp \left( - 4 \pi \int_{\mathbb{R}^2}
  g (z) \exp^{\diamond} (\alpha \omega) (\mathd z) \right)
  \label{eq:exponentiallaw1},
\end{equation}
where $\omega \in B^{- \delta}_{p, p, \ell} (\mathbb{R}^2)$ and where $\mu$ is
the law of $(- \Delta_z + m^2)^{- 1 / 2} (\xi)$ on $B^{- \delta}_{p, p, \ell}
(\mathbb{R}^2)$, we have
\[ \mathbb{E} [F (\phi (0, \cdot))] = \int F (\omega) \mathd \kappa_g (\omega)
   \label{eq:exponentialreduction1}, \]
for any bounded measurable function defined on $B^{- \delta}_{p, p, \ell}$.

In order to prove the existence and uniqueness of the solution $\bar{\phi}$ to
equation~{\eqref{eq:exponential2}} and dimensional
reduction~{\eqref{eq:exponentialreduction1}} for the random field $\phi =
\bar{\phi} + \mathcal{I} \xi$, we need to introduce the following two
approximate equations
\begin{eqnarray}
  & (- \Delta_x - \Delta_z + m^2) (\phi_{\epsilon}) +
  \mathfrak{a}_{\epsilon}^{\asterisk 2} \asterisk [g (z) \alpha \exp (\alpha
  \phi_{\epsilon} - C_{\epsilon})] = \mathfrak{a}_{\epsilon} \asterisk \xi & 
  \label{eq:exponential3}\\
  & (- \Delta_x - \Delta_z + m^2) (\bar{\phi}_{\epsilon}) +
  \mathfrak{a}_{\epsilon}^{\asterisk 2} \asterisk [g (z) \alpha \exp (\alpha
  \bar{\phi}_{\epsilon}) \eta_{\epsilon} (x, z)] = 0 & 
  \label{eq:exponential4}
\end{eqnarray}
where $\mathfrak{a}$ is a positive function satisfying
Hypothesis~H$\mathcal{A} 1$ (see Section~\ref{subsection_extension1}) and
$C_{\epsilon} : = \frac{\alpha^2}{2} \mathbb{E} [\mathcal{I}
(\mathfrak{a}_{\epsilon} \asterisk \xi)]$, $\bar{\phi}_{\epsilon} : =
\phi_{\epsilon} - \mathcal{I} (\mathfrak{a}_{\epsilon} \asterisk \xi)$,
$\eta_{\epsilon} : = \exp^{\diamond} (\alpha \mathcal{I}
(\mathfrak{a}_{\epsilon} \asterisk \xi))$ and $\bar{\phi}_{\epsilon} =
\phi_{\epsilon} - \mathfrak{a}_{\varepsilon} \asterisk \mathcal{I} \xi$. The
constant $C_{\epsilon}$ is chosen in such a way that $\eta_{\epsilon}
\rightarrow \eta$ in $B^s_{p, p, \ell}$ for suitable $s < 0$, $1 < p \leqslant
2$ and $\ell > 0$ (see Section~\ref{subsection_probabilistic}).

For equations~{\eqref{eq:exponential3}} and~{\eqref{eq:exponential4}} we have
the following fundamental dimensional reduction result.

\begin{theorem}
  \label{theorem_regularexponential}Equation~{\eqref{eq:exponential4}} admits
  a unique solution in $C^0_{\ell} (\mathbb{R}^4)$ for $\ell$ big enough.
  Furthermore we have that that, for any $\epsilon > 0$,
  \begin{equation}
    \mathbb{E} [F (\phi_{\epsilon} (0, z))] = \int_{\mathcal{C}^{0 -}_{\ell}
    (\mathbb{R}^2)} F (\omega) \mathd \kappa_{\epsilon} (\omega)
    \label{eq:exponential5},
  \end{equation}
  for all $z \in \mathbb{R}^2$, where
  \[ \frac{\mathd \kappa_{\epsilon}}{\mathd \mu_{\epsilon}} = \exp \left( - 4
     \pi \int_{\mathbb{R}^2} g (z) \exp (\alpha \omega (z) - \alpha
     C_{\epsilon}) \mathd z \right) \]
  and $\mu_{\epsilon}$ is the Gaussian measure on $\mathcal{C}^{0 -}_{\ell}$
  with covariance $\mathbb{E} [\omega (z) \omega (z')] =
  \mathfrak{a}^{\asterisk 2}_{\epsilon} \asterisk \mathcal{G}_z (z - z')$,
  where $\mathcal{G}_z$ is the Green function of the operator $(- \Delta_z +
  m^2)$. Finally the unique solution to equation~{\eqref{eq:exponential4}}
  satisfies $\alpha \bar{\phi}_{\varepsilon} \leqslant 0$.
\end{theorem}

\begin{proof}
  The proof is an application of Theorem~\ref{theorem_extension1} to
  equation~{\eqref{eq:exponential3}} using the fact that $\exp (\alpha y -
  C_{\epsilon})$ satisfies Hypothesis~C in Section~\ref{section_discrete}.
  
  The fact that $\alpha \bar{\phi}_{\varepsilon} \leqslant 0$ follows from an
  application of the maximum principle to the function $\hat{z} \mapsto
  \bar{\phi}_{\epsilon} (\hat{z}) \rho_{\ell} (\lambda \hat{z})$ where
  $\hat{z} = (x, z) \in \mathbb{R}^4$ and $\lambda > 0$ small enough.
\end{proof}

In order to prove existence, uniqueness and the reduction principle for
equation~{\eqref{eq:exponential1}} we have now to study the behaviour of the
regularised noises $\eta_{\epsilon}$ and their convergence to $\eta$ in the
Besov spaces of the form $B^s_{p, p, \ell}$, as $\varepsilon \rightarrow 0$.
Once we have established this convergence, we can give a meaning to
equation~{\eqref{eq:exponential2}} and we are able to prove that there exists
a subsequence of $\epsilon_n \rightarrow 0$ such that $\phi_{\epsilon_n}
\rightarrow \phi$ in probability in $B^{- \delta}_{p, p, \ell}$. This fact
will permit us to prove equality~{\eqref{eq:exponentialreduction1}}.

\subsection{Probabilistic analysis}\label{subsection_probabilistic}

In this subsection we propose an analysis of the regularity of the noise
$\eta$ and $\eta_{\epsilon}$, for $\epsilon > 0$. First of all we note that
\begin{eqnarray}
  \eta_{\epsilon} & = & \sum_{k = 0}^{\infty} \frac{\alpha^k}{k!} (\mathcal{I}
  \xi_{\epsilon})^{\diamond k},  \label{eq:noise1}\\
  \eta & = & \sum_{k = 0}^{\infty} \frac{\alpha^k}{k!} (\mathcal{I}
  \xi)^{\diamond k},  \label{eq:noise2}
\end{eqnarray}
where
\[ (\mathcal{I} \xi_{\epsilon})^{\diamond k} = \underbrace{\mathcal{I}
   \xi_{\epsilon} \diamond \ldots \diamond \mathcal{I} \xi_{\epsilon}}_{k
   \quad \tmop{times}} {}^{}, \]
the symbol $\diamond$ denotes the Wick product, $\xi_{\epsilon} =
\mathfrak{a}_{\epsilon} \asterisk \xi$ and $\mathcal{I} \xi_{\epsilon}$ is
defined correspondingly. The previous expressions are well defined as $L^2
(\mu)$ convergent series for $| \alpha | < 4 \sqrt{2} \pi$. Indeed we have
that for any \ smooth function $g$ exponentially decaying at infinity \
\begin{equation}
  \mathbb{E} [| \langle \eta_{\epsilon}, g \rangle |^2] = \int_{\mathbb{R}^8}
  g (\hat{z}) g (\hat{z}') \exp (\alpha^2 \mathcal{G}_{\epsilon} (\hat{z} -
  \hat{z}')) \mathd \hat{z} \mathd \hat{z}' \label{eq:noise3}
\end{equation}
where, hereafter, we write $\mathcal{G}_{\epsilon} = \mathfrak{a}^{\asterisk
2}_{\epsilon} \asterisk \mathcal{G}$, where $\mathcal{G}$ is the Green
function associated with the operator $(- \Delta_{\hat{z}} + m^2)^{- 2}$,
$\hat{z} = (x, z) \in \mathbb{R}^4$. It is well known (see Proposition A.1
of~{\cite{Albeverio2002}}) that for $\hat{z} \in \mathbb{R}^4$ such that $|
\hat{z} | \leqslant 1$ there exists a constant $C_1 > 0$ for which the
following inequality holds
\begin{equation}
  \mathcal{G} (\hat{z}) \leqslant - \frac{2}{(4 \pi)^2} \log_+ (| \hat{z} |) +
  C_1 \label{eq:probabilistic1} .
\end{equation}
Furthermore for $| \hat{z} | \geqslant 1$ there exists two constants $C_2, C_3
> 0$ for which we get
\begin{equation}
  \mathcal{G} (\hat{z}) \leqslant C_2 \exp (- C_3 | \hat{z} |)
  \label{eq:probabilistic2} .
\end{equation}
An easy consequence of the inequalities~{\eqref{eq:probabilistic1}}
and~{\eqref{eq:probabilistic2}} are the following inequalities
\begin{equation}
  \mathfrak{a}_{\epsilon}^{\asterisk 2} \asterisk \mathcal{G} (\hat{z})
  \leqslant - \frac{2}{(4 \pi)^2} \log_+ (| \hat{z} |) + C_4,
  \label{eq:probabilistic3}
\end{equation}
when $4 \tmop{diam} (\tmop{supp} (\mathfrak{a})) \epsilon \leqslant | \hat{z}
| \leqslant 1$ and for some constant $C_4 > 0$; \
\begin{equation}
  \mathfrak{a}_{\epsilon}^{\asterisk 2} \asterisk \mathcal{G} (\hat{z})
  \leqslant - \frac{2}{(4 \pi)^2} \log (\epsilon) + C_5,
  \label{eq:probabilistic4}
\end{equation}
when $| \hat{z} | < 2 \tmop{diam} (\tmop{supp} (\mathfrak{a})) \epsilon$ and
for some constant $C_5 > 0$; and finally
\begin{equation}
  \mathfrak{a}_{\epsilon}^{\asterisk 2} \asterisk \mathcal{G} (\hat{z})
  \leqslant C_6 \exp (- C_7 | \hat{z} |) \label{eq:probabilistic5},
\end{equation}
when $| \hat{z} | \geqslant 1$ and for some constants $C_6, C_7 > 0$. It is
important to note that the constants $C_4, C_5, C_6, C_7$ are independent of
$\epsilon$. The previous inequalities and eq.~{\eqref{eq:noise3}} imply that
$\mathbb{E} [| \langle \eta_{\epsilon}, g \rangle |^2] < + \infty$ for any
$\epsilon \geqslant 0$ and $| \alpha | < 4 \sqrt{2} \pi$ (with $\eta_0 =
\eta$).

This means that, if we denote by $D_k$ the $k \geqslant - 1$ Littlewood-Paley
block (see Appendix \ref{appendix_besov} for the definition of this concept),
using the fact that $D_k$ exponentially decreases at infinity, we get
\[ \mathbb{E} [\| \langle \eta_{\epsilon}, D_k (\hat{z} - \cdot) \rangle
   \|_{L^2_{\rho_{\ell}}}^2] = \int_{\mathbb{R}^4} \mathbb{E} [\langle
   \eta_{\epsilon}, D_k \rangle^2] (1 + | \hat{z} |)^{2 \ell} \mathd \hat{z}
   \lesssim \mathbb{E} [\langle \eta, D_k \rangle^2] < + \infty, \]
whenever $\alpha < 4 \sqrt{2} \pi$. Since $D_k (x, z) = 2^{4 k} D_0 (2^k x,
2^k z)$ and using~{\eqref{eq:probabilistic1}} and~{\eqref{eq:probabilistic2}},
we obtain
\begin{eqnarray}
  \mathbb{E} [\langle \eta_{\epsilon}, D_k \rangle^2] & \lesssim &
  \int_{\mathbb{R}^4} D_k (\hat{z}) D_k (\hat{z}') \exp (\alpha^2 \cdot
  \mathfrak{a}_{\epsilon}^{\asterisk 2} \asterisk \mathcal{G} (\hat{z} -
  \hat{z}')) \mathd \hat{z} \mathd \hat{z}' \nonumber\\
  & \lesssim & \int_{\mathbb{R}^8} D_k (\hat{z}) D_k (\hat{z}') | \hat{z} -
  \hat{z}' |^{- \frac{2 \alpha^2}{(4 \pi)^2}} \mathd \hat{z} \mathd \hat{z}' +
  \nonumber\\
  &  & + \left( 1 + \epsilon^{4 - \frac{2 \alpha^2}{(4 \pi)^2}} \right)
  \int_{\mathbb{R}^8} D_k (\hat{z}) D_k (\hat{z}') \mathd \hat{z} \mathd
  \hat{z}'  \nonumber\\
  & \lesssim & 2^{\frac{2 \alpha^2}{(4 \pi)^2} k} \left( \int_{\mathbb{R}^8}
  D_0 (\hat{z}) D_0 (\hat{z}') \mathd \hat{z} \mathd \hat{z}' \right) \lesssim
  2^{\frac{2 \alpha^2}{(4 \pi)^2} k} \nonumber
\end{eqnarray}
where all the constants implied in the symbol $\lesssim$ are independent of
$\epsilon$. This means that
\[ \mathbb{E} [\| \eta_{\epsilon} \|^2_{B^s_{2, 2, \ell}}] \lesssim \sum^{+
   \infty}_{k = 1} 2^{\frac{2 \alpha^2}{(4 \pi)^2} k + 2 s k}, \]
which is finite and uniformly bounded in $\epsilon$ whenever $s < -
\frac{\alpha^2}{(4 \pi)^2}$ and for $\ell > 0$ large enough. Furthermore we
have that
\begin{eqnarray}
  \mathbb{E} [\| \eta - \eta_{\epsilon} \|_{B^s_{2, 2, \ell}}^2]^{1 / 2} &
  \lesssim & \left( \sum_{k = - 1}^{+ \infty} 2^{2 s k} \mathbb{E}[| \langle
  \eta - \eta_{\epsilon}, D_k \rangle |^2] \right)^{1 / 2} \nonumber\\
  & \lesssim & \sum_{k = - 1}^{+ \infty} 2^{s k} \sum^{+ \infty}_{n = 0}
  \frac{\alpha^n}{n!} \mathbb{E} [| \langle \mathcal{I} \xi^{\diamond n} -
  \mathcal{I} \xi_{\epsilon}^{\diamond n}, D_k \rangle |^2]^{1 / 2} \nonumber
\end{eqnarray}
for $\ell$ big enough. On the other hand, from the previous estimate, we get
\begin{eqnarray}
  \sum_{k = - 1}^{+ \infty} 2^{s k} \sum^{+ \infty}_{n = 0}
  \frac{\alpha^n}{n!} \mathbb{E} [| \langle \mathcal{I} \xi^{\diamond n} -
  \mathcal{I} \xi_{\epsilon}^{\diamond n}, D_k \rangle |^2]^{1 / 2} &
  \leqslant & \sum_{k = - 1}^{+ \infty} 2^{s k} \sum^{+ \infty}_{n = 0}
  \frac{\alpha^n}{n!} \mathbb{E} [| \langle \mathcal{I} \xi^{\diamond n}, D_k
  \rangle |^2]^{1 / 2} + \nonumber\\
  &  & + \sum_{k = - 1}^{+ \infty} 2^{s k} \sum^{+ \infty}_{n = 0}
  \frac{\alpha^n}{n!} \mathbb{E} [| \langle \mathcal{I}
  \xi_{\epsilon}^{\diamond n}, D_k \rangle |^2]^{1 / 2} < C, \nonumber
\end{eqnarray}
for some constant $C > 0$ independent of $\epsilon$, whenever $s < -
\frac{\alpha^2}{(4 \pi)^2}$. Furthermore, since $\mathfrak{a}_{\epsilon}$ is a
regular mollifier and by the properties of Wick product, we obtain
\[ \mathbb{E} [| \langle \mathcal{I} \xi^{\diamond n} - \mathcal{I}
   \xi_{\epsilon}^{\diamond n}, D_k \rangle |^2] \rightarrow 0, \]
as $\varepsilon \rightarrow 0$ and for any $k \geq - 1$. The above inequality,
together with the Lebesgue dominated convergence theorem, implies that
\[ \mathbb{E} [\| \eta - \eta_{\epsilon} \|_{B^s_{2, 2, \ell}}^2]^{1 / 2}
   \rightarrow 0, \]
when $\epsilon \rightarrow 0$ for any $\ell > 0$ large enough. We have thus
proven the following lemma.

\begin{lemma}
  \label{lemma_L2}For $| \alpha | < 4 \sqrt{2} \pi$ and $s < -
  \frac{\alpha^2}{(4 \pi)^2}$ and $\ell$ large enough we have that
  $\eta_{\epsilon} \rightarrow \eta$ as $\epsilon \rightarrow 0$ in $L^2
  (\mathcal{W} ; B^s_{2, 2, \ell} (\mathbb{R}^4), \mathd \mu)$ and thus in
  probability in $B^s_{2, 2, \ell} (\mathbb{R}^4)$.
\end{lemma}

We want to use the previous lemma to prove the following theorem.

\begin{theorem}
  \label{theorem_Lp}For $| \alpha | < 4 \sqrt{2} \pi$, $1 < p \leqslant 2$, $s
  < - \frac{\alpha^2 (p - 1)}{(4 \pi)^2}$ and $\ell$ large enough we have that
  $\eta_{\epsilon} \rightarrow \eta$, as $\epsilon \rightarrow 0$, in $L^p
  (\mathcal{W} ; B^s_{p, p, \ell} (\mathbb{R}^4), \mathd \mu)$ and thus in
  probability in $B^s_{p, p, \ell} (\mathbb{R}^4)$.
\end{theorem}

In order to prove Theorem \ref{theorem_Lp} we introduce the following lemma.

\begin{lemma}
  \label{lemma_garban}Let $B_r (\hat{z})$ be the ball of radius $r$ and centre
  in $\hat{z} \in \mathbb{R}^4$ then for any $r < R$, and for any $1 < p < 2$
  \[ \mathbb{E} \left[ \left( \int_{B_r (\hat{z})} \mathd \eta_{\epsilon}
     \right)^p \right] \lesssim r^{- \frac{\alpha^2}{(4 \pi)^2} p (p - 1) + 4
     p} \]
  where the constants depend only $R$ and it is uniform on $\epsilon
  \rightarrow 0$ and $| \alpha | < 4 \sqrt{2} \pi$.
\end{lemma}

\begin{proof}
  The proof can be found in Proposition 2.7 of~{\cite{Vargas2010}}.
\end{proof}

\begin{proof*}{Proof of Theorem~\ref{theorem_Lp}}
  First of all we have that, for any $\varepsilon > 0$ and $k \geqslant 0$, \
  \
  \[ \mathbb{E} [\| \langle \eta_{\epsilon}, D_k (\hat{z} - \cdot) \rangle
     \|_{L^p_{\ell}}^p] = \int_{\mathbb{R}^4} \mathbb{E} [\langle
     \eta_{\epsilon}, D_k \rangle^p] (1 + | \hat{z} |)^{p \ell} \mathd \hat{z}
     \lesssim \mathbb{E} [\langle \eta_{\epsilon}, D_k \rangle^p] . \]
  Then
  \[ \mathbb{E} [\langle \eta_{\epsilon}, D_k \rangle^p] \lesssim \mathbb{E}
     \left[ \left( \int_{B_{r_k} (0)} D_k (\hat{z}) \mathd \eta_{\epsilon}
     (\hat{z}) \right)^p \right] +\mathbb{E} \left[ \left( \int_{\mathbb{R}^4
     \setminus B_{r_k} (0)} D_k (\hat{z}) \mathd \eta_{\epsilon} (\hat{z})
     \right)^2 \right]^{2 / p} . \]
  If we choose $r_k = 2^{- k} \log (2^{\beta k}),$from
  Lemma~\ref{lemma_garban} we have, for any $k \geqslant 0$, \
  \[ \mathbb{E} \left[ \left( \int_{B_{r_k} (0)} D_k (\hat{z}) \mathd
     \eta_{\epsilon} (\hat{z}) \right)^p \right] \lesssim 2^{4 k p}
     2^{\frac{\alpha^2}{(4 \pi)^2} p (p - 1) - 4 p k} \log
     (2^{\beta})^{\frac{\alpha^2}{(4 \pi)^2} p (p - 1) - 4 p k} \]
  Furthermore, if we denote by $\gamma > 0$ the real number such that $D_0
  (\hat{z}) \lesssim \exp (- \gamma | \hat{z} |)$, we obtain
  \[ \begin{array}{lll}
       \mathbb{E} \left[ \left( \int_{\mathbb{R}^4 \setminus B_{r_k} (0)} D_k
       (\hat{z}) \mathd \eta_{\epsilon} (\hat{z}) \right)^2 \right]^{2 / p} &
       \lesssim & 2^{4 k - \gamma \beta k} \int_{\mathbb{R}^8} \exp (- \gamma
       (| \hat{z} | + | \hat{z}' |) + \alpha^2 \mathcal{G}_{\epsilon} (\hat{z}
       - \hat{z}')) \mathd \hat{z} \mathd \hat{z}'\\
       & \lesssim & 2^{4 k - \gamma \beta k} .
     \end{array} \]
  If we choose $\beta$ such that $4 - \gamma \beta < \frac{\alpha^2}{(4
  \pi)^2} p (p - 1)$ we obtain that
  \begin{equation}
    \mathbb{E} [| \langle \eta_{\epsilon}, D_k \rangle |^p] \lesssim
    2^{\frac{\alpha^2}{(4 \pi)^2} p (p - 1)} \log
    (2^{\beta})^{\frac{\alpha^2}{(4 \pi)^2} p (p - 1) - 4 p k} \label{eq:Lp1}
  \end{equation}
  uniformly in $\epsilon$. This implies that $\mathbb{E} [\| \eta_{\epsilon}
  \|_{B^s_{p, p, \ell}}^p]$ is bounded for $s < \frac{\alpha^2}{(4 \pi)^2} (p
  - 1)$ uniformly in $\epsilon > 0$.
  
  Furthermore, by Lemma~\ref{lemma_L2}, $\mathbb{E} [| \langle
  \eta_{\epsilon}, D_k \rangle |^2]$ is uniformly bounded. This means that the
  random variables $| \langle \eta_{\epsilon}, D_k \rangle |^p$ (for $p < 2$)
  are uniformly integrable. On the other hand, by Lemma~\ref{lemma_L2} and
  since $D_k \in B^2_{2, 2, \ell}$, $| \langle \eta_{\epsilon}, D_k \rangle
  |^p$ converges to $| \langle \eta, D_k \rangle |^p$ in probability, we have
  that $\mathbb{E} [| \langle \eta_{\epsilon}, D_k \rangle |^p] \rightarrow
  \mathbb{E} [| \langle \eta, D_k \rangle |^p]$. Thus by the
  bound~{\eqref{eq:Lp1}} we can use the Lebesgue dominated converge theorem
  for computing the limit of $\mathbb{E} [\| \eta_{\epsilon} \|_{B^s_{p, p,
  \ell}}^p]$, obtaining
  \[ \mathbb{E} [\| \eta_{\epsilon} \|_{B^s_{p, p, \ell}}^p] \rightarrow
     \mathbb{E} [\| \eta \|_{B^s_{p, p, \ell}}^p], \]
  as $\varepsilon \rightarrow 0$. But, since $\mathbb{E} [\| \eta_{\epsilon}
  \|_{B^s_{p, p, \ell}}^p]$ is uniformly bounded and since, by
  Lemma~\ref{lemma_L2}, as $\varepsilon \rightarrow 0$, $\eta_{\epsilon}
  \rightarrow \eta$ in probability in $B^{- 2}_{2, 2, \ell}$, we have that
  $\eta_{\epsilon}$ converges weakly to $\eta$ in $L^p (\mathcal{W}, B^s_{p,
  p, \ell}, \mathd \mu)$. Since $L^p (\mathcal{W}, B^s_{p, p, \ell}, \mathd
  \mu)$ is a uniformly convex space, being the space of $L^p$ functions taking
  values in the uniformly convex space $B^s_{p, p, \ell}$
  (see~{\cite{Boyko2009}}), by the weak convergence $\eta_{\epsilon}
  \rightarrow \eta$ and the convergence of the $L^p$ norm $\mathbb{E} [\|
  \eta_{\epsilon} \|_{B^s_{p, p, \ell}}^p] \rightarrow \mathbb{E} [\| \eta
  \|_{B^s_{p, p, \ell}}^p]$, we obtain that $\eta_{\epsilon}$ converges
  strongly to $\eta$ in $B^s_{p, p, \ell} (\mathbb{R}^4)$, as $\varepsilon
  \rightarrow 0$. 
\end{proof*}

\subsection{Analysis of the elliptic SPDE}

In this section we want to prove the following theorem

\begin{theorem}
  \label{theorem_exponentialmain}For any $| \alpha | < 4 \sqrt{2} \pi$, $1 < p
  \leqslant 2$ such that $p < \frac{2 (4 \pi)^2}{\alpha^2}$ and $s < 0$ such
  that $- \frac{2 (p - 1)}{p} < s < - \frac{\alpha^2 (p - 1)}{(4 \pi)^2}$ we
  have that equation~{\eqref{eq:exponential2}} admits a unique solution in
  $B^{s + 2}_{p, p, \ell}$ for $\ell$ large enough and such that $\alpha
  \bar{\phi} \leqslant 0$. \ Furthermore we have that there exists a
  subsequence $\varepsilon_n \rightarrow 0$ such that \
  $\bar{\phi}_{\epsilon_n} \rightarrow \bar{\phi}$ in $B^{s + 2 - \delta}_{p,
  p, \ell + \delta'}$ (for some $\delta, \delta' > 0$ small enough), as $n
  \rightarrow \infty$, almost surely.
\end{theorem}

Hereafter, for $t > 0$, we write $\mathbb{B}^t_{p, p, \ell} (\mathbb{R}^4) =
B^t_{p, p, \ell} (\mathbb{R}^4) \cap L^{\infty} (\mathbb{R}^4)$. The space
$\mathbb{B}^t_{p, p, \ell}$ has a natural norm given by the sum of the norms
of $B^t_{p, p, \ell} (\mathbb{R}^4)$ and $L^{\infty} (\mathbb{R}^4)$.
Furthermore we can equip the space $\mathbb{B}^t_{p, p, \ell}$ with a
different notion of convergence of sequences: we say that a sequence $h_n \in
\mathbb{B}^t_{p, p, \ell}$ pseudoconverges to $h \in \mathbb{B}^t_{p, p,
\ell}$ if $h_n$ converges to $h$ strongly in $B^t_{p, p, \ell}$ and $\sup_n \|
h_n \|_{\infty} < + \infty$ (it is important to note that $\| h \|_{\infty}
\leqslant \liminf_{n \rightarrow + \infty} \| h_n \|_{\infty}$).

\begin{lemma}
  \label{lemma_interpolation}For any $t > 0$ and $1 < p \leqslant 2$, we have
  that $\mathbb{B}^t_{p, p, \ell} \subset B^{t (p - 1)}_{q, q, \ell}$ (where
  $1 / p + 1 / q = 1$) and the immersion is continuous with respect to the
  pseudoconvergence in $\mathbb{B}^t_{p, p, \ell}$. Furthermore
  $\mathbb{B}^t_{p, p, \ell}$ is a Banach algebra.
\end{lemma}

\begin{proof}
  By Proposition \ref{proposition_interpolation}, the interpolation space
  $(B^0_{\infty, \infty}, B^t_{p, p, \ell})_{\frac{p}{q}, q}$ between
  $L^{\infty} \subset B^0_{\infty, \infty}$ and $B^t_{p, p, \ell}$ is exactly
  $B^{t \frac{p}{q}}_{q, q, \ell} = B^{t (p - 1)}_{q, q, \ell}$, where $q =
  \frac{p - 1}{p}$. This means that $\mathbb{B}^t_{p, p, \ell} = L^{\infty}
  \cap B^t_{p, p, \ell} \subset B^0_{\infty, \infty} \cap B^s_{p, p, \ell}$ is
  continuously embedded in $B^{t (p - 1)}_{q, q, \ell}$. The fact that
  $\mathbb{B}_{p, p, \ell}^t$ is an algebra is proven in Corollary 2.86
  of~{\cite{Bahouri2011}}.
\end{proof}

We introduce the space $\mathfrak{M}^s_{p, p, \ell} \subset B^s_{p, p, \ell}$
(with $s$ and $p$ as in Theorem \ref{theorem_exponentialmain}) which is the
set of Radon $\sigma$-finite measures $\mu$ contained in $B^s_{p, p, \ell}$
such that $| \mu | \in B^s_{p, p, \ell}$, where $| \mu | = \mu_+ + \mu_-$, and
where $\mu_+$ and $\mu_-$ are the unique positive measures such that $\mu =
\mu_+ - \mu_-$. We can define a notion of pseudoconvergence on
$\mathfrak{M}^s_{p, p, \ell}$ in the following way: a sequence $\mu_n \in
\mathfrak{M}^s_{p, p, \ell}$ pseudoconverges to $\mu \in \mathfrak{M}^s_{p, p,
\ell}$ if $\mu_n$ converges strongly in $B^s_{p, p, \ell}$ and $\sup_n \| |
\mu_n | \|_{B^s_{p, p, \ell}} < + \infty$.

We now consider the natural product $\cdot$ between smooth functions which
are in $\mathbb{B}^{s + 2 - \delta'}_{p, p, \ell + \delta}$ and and smooth
measures in $\mathfrak{M}^s_{p, p, \ell}$.

\begin{lemma}
  \label{lemma_multiplication}If $s, p$ satisfy the conditions of
  Theorem~\ref{theorem_exponentialmain} and considering $\delta, \delta'
  \geqslant 0$ small enough, the product $\cdot$ can be extended in a unique
  and associative way from $\mathbb{B}^{s + 2 - \delta'}_{p, p, \ell + \delta}
  \times \mathfrak{M}^s_{p, p, \ell}$ into $\mathfrak{M}^s_{p, p, \ell}$. This
  extension is (weakly) continuous with respect to the pseudoconvergence in
  $\mathbb{B}^{s + 2 - \delta'}_{p, p, \ell + \delta}$ and $\mathfrak{M}^s_{p,
  p, \ell}$. Furthermore for any $h \in \mathbb{B}^{s + 2 - \delta'}_{p, p,
  \ell + \delta}$ and $\mu \in \mathfrak{M}^s_{p, p, \ell}$ we have
  \begin{equation}
    \| | h \cdot \mu | \|_{B^s_{p, p, \ell}} \lesssim \| h \|_{\infty} \cdot
    \| | \mu | \|_{B^s_{p, p, \ell}} . \label{eq:exponentialmain1}
  \end{equation}
\end{lemma}

\begin{proof}
  First of all we note that the product $\cdot$ is well defined and (strongly)
  continuous as bilinear functional from $B^{(s + 2 - \delta') (p - 1)}_{q, q,
  \ell + \delta} \times B^s_{p, p, \ell}$ into $B^s_{1, 1, 2 \ell}$. Indeed
  this is a consequence of the fact that $(s + 2) (p - 1) + s = p s + 2 (p -
  1) = p (s + 2 (1 - 1 / p)) > 0$, if $s$ and $p$ satisfy the hypotheses of
  Theorem~\ref{theorem_exponentialmain}, and of Proposition
  \ref{propisition_product}.
  
  The only thing that remains to prove is
  inequality~{\eqref{eq:exponentialmain1}}, since the other statements of the
  lemma easily follow from it. First of all, using the equivalent norm on
  Besov spaces of Proposition \ref{proposition_equivalentnorm}, we note that
  if $\mu$ is a Radon measure we have
  \[ \| | \mu | \|_{B^s_{p, p, \ell}} \sim \left( \int_{\mathbb{R}^4 \times
     (0, 1]} \left( \frac{| \mu | (B_{\lambda} (\hat{z}))}{\lambda^s}
     \right)^p \rho_{\ell} (\hat{z})^p \mathd \hat{z} \mathd \lambda
     \right)^{1 / p} . \]
  Using again Proposition \ref{proposition_equivalentnorm}, if $h$ is a
  continuous bounded function we have that
  \[ \| | h \cdot \mu | \|_{B^s_{p, p, \ell}} \lesssim \left(
     \int_{\mathbb{R}^4 \times (0, 1]} \left( \frac{| h \cdot \mu |
     (B_{\lambda} (\hat{z}))}{\lambda^s} \right)^p \rho_{\ell} (\hat{z})^p
     \mathd \hat{z} \mathd \lambda \right)^{1 / p} \lesssim \| h \|_{\infty}
     \cdot \| | \mu | \|_{B^s_{p, p, \ell}} . \]
  If $h$ is a generic function $\mathbb{B}^{s + 2 - \delta'}_{p, p, \ell +
  \delta}$ there exists a sequence of smooth functions such that, as $n
  \rightarrow \infty$, $h_n \rightarrow h$ in $B^{s + 2 - \delta'}_{p, p, \ell
  + \delta}$ and $\| h_n \|_{\infty} \leqslant \| h \|_{\infty}$. By the first
  part of the lemma we have that $| h_n \cdot \mu |$ converges weakly in
  $\mathcal{S}' (\mathbb{R}^4)$ to $| h \cdot \mu |$ from which we obtain
  \[ \| | h \cdot \mu | \|_{B^s_{p, p, \ell}} \leqslant \liminf \| | h_n \cdot
     \mu | \|_{B^s_{p, p, \ell}} \lesssim \| | \mu | \|_{B^s_{p, p, \ell}}
     \liminf \| h_n \|_{\infty} \lesssim \| h \|_{\infty} \| | \mu |
     \|_{B^s_{p, p, \ell}} . \]
\end{proof}

We introduce the following map
\[ \mathcal{K}_g (\mu, \bar{\varphi}) : = - \alpha (- \Delta_{\hat{z}} +
   m^2)^{- 1} (g (z) G (\alpha \bar{\varphi} (\hat{z})) \cdot \mu (\hat{z})),
\]
for any $| \alpha | < 4 \sqrt{2} \pi$, where $G : \mathbb{R} \rightarrow
\mathbb{R}_+$ is an increasing smooth bounded function with all bounded
derivatives such that $G (x) = \exp (x)$ for $x \leqslant 0$. It is clear that
if $\alpha \bar{\phi} \leqslant 0$ and if $\bar{\phi}$ solves
equation~{\eqref{eq:exponential2}} then we have
\[ \bar{\phi} = \mathcal{K}_g (\eta, \bar{\phi}) . \]
\begin{lemma}
  \label{lemma_exponentialexistence}For $p, s$ satisfying the hypotheses of
  Theorem~\ref{theorem_exponentialmain}, any $\mu \in \mathfrak{M}^s_{+, p, p,
  \ell}$ fixed and $\delta, \delta' > 0$ small enough, there exists at least a
  solution $\bar{\varphi} \in B^{s + 2 - \delta'}_{p, p, \ell + \delta}$ to
  the equation
  \begin{equation}
    \bar{\varphi} = \mathcal{K}_g (\mu, \bar{\varphi}) . \label{eq:K1}
  \end{equation}
\end{lemma}

\begin{proof}
  We want to use Schaefer's fixed-point theorem (see Theorem~4 Section~9.2
  Chapter 9 of~{\cite{Evans1998}}) to prove the lemma. In order to do this we
  have to prove that $\mathcal{K}_g$ is continuous in $\varphi$, that it maps
  any bounded set into a compact set and that the set of solutions to the
  equations
  \[ \bar{\varphi} = \lambda \mathcal{K}_g (\mu, \bar{\varphi}) \]
  are bounded uniformly for all $0 \leqslant \lambda \leqslant 1$.
  
  The continuity of $\mathcal{K}_g$ (in both $\mu$ and $\varphi$) is a
  consequence of Lemma~\ref{lemma_multiplication}. Indeed in the cone
  $\mathfrak{M}^s_{+, p, p, \ell}$ the pseudoconvergence coincides with the \
  convergence in $B^s_{p, p, \ell}$ since if $\mu \in \mathfrak{M}^s_{+, p, p,
  \ell}$ then $| \mu | = \mu$. Furthermore, the map $\varphi \rightarrow G
  (\varphi)$ is continuous (actually $C^1$) from $B^{s + 2 - \delta'}_{p, p,
  \ell + \delta}$ into $\mathbb{B}^{s + 2 - \delta'}_{p, p, \ell + \delta}$.
  Finally the linear operator $(- \Delta_{\hat{z}} + m^2)^{- 1}$ is continuous
  from $B^s_{p, p, \ell}$ into $B^{s + 2}_{p, p, \ell}$ and thus compact from
  $B^s_{p, p, \ell}$ into $B^{s + 2 - \delta'}_{p, p, \ell + \delta}$ (since,
  by Proposition \ref{proposition_inclusion}, the immersion $B^{s + 2}_{p, p,
  \ell} \hookrightarrow B^{s + 2 - \delta'}_{p, p, \ell + \delta}$ is
  compact). Thus by Lemma~\ref{lemma_multiplication}, the map $\mathcal{K}_g$
  is weakly continuous as a map into $B^{s + 2}_{p, p, \ell}$ and strongly
  continuous as a map into $B^{s + 2 - \delta'}_{p, p, \ell + \delta}$.
  
  The map $\mathcal{K}_g$ is compact since the following inequality holds
  \begin{equation}
    \| \mathcal{K}_g (\mu, \varphi) \|_{B^{s + 2}_{p, p, \ell}} \lesssim \| G
    \|_{\infty} \| \mu \|_{B^s_{p, p, \ell}} \label{eq:inequalityboundeness}
  \end{equation}
  and, by Proposition \ref{proposition_inclusion}, we have the compact
  immersion $B^{s + 2}_{p, p, \ell} \hookrightarrow B^{s + 2 - \delta'}_{p, p,
  \ell + \delta}$.
  
  Finally the uniform boundedness in $\lambda$ follows from
  inequality~{\eqref{eq:inequalityboundeness}}. This proves the thesis of the
  lemma.
\end{proof}

\begin{lemma}
  \label{lemma_exponentialuniqueness}Under the hypotheses of
  Theorem~\ref{theorem_exponentialmain} the solution to
  equation~{\eqref{eq:K1}} is unique in $B^{s + 2 - \delta'}_{p, p, \ell +
  \delta}$ for $\delta, \delta' \geqslant 0$ small enough.
\end{lemma}

\begin{proof}
  Let $J : \mathbb{R} \rightarrow \mathbb{R}$ be a smooth, bounded, strictly
  increasing function such that $J (0) = 0$ and let $\bar{\varphi}_1$ and
  $\bar{\varphi}_2$ be two solutions to equation~{\eqref{eq:K1}}. By
  Lemma~\ref{lemma_interpolation}, $J (\bar{\varphi}_1 - \bar{\varphi}_2) \in
  B^{(s + 2 - \delta') (p - 1)}_{q, q, \ell + \delta}$ which implies that
  $\rho_{\ell'} (\lambda \hat{z}) J (\bar{\varphi}_1 - \bar{\varphi}_2) \in
  (B^s_{p, p, \ell})^{\asterisk}$ for $\delta, \delta' \geqslant 0$ small
  enough, $\ell'$ large enough and any $\lambda > 0$. This means that
  \[ \langle \rho_{\ell'} (\lambda \hat{z}) J (\bar{\varphi}_1 -
     \bar{\varphi}_2), (- \Delta_{\hat{z}} + m^2) (\bar{\varphi}_1 -
     \bar{\varphi}_2 - \mathcal{K}_g (\mu, \bar{\varphi}_1) + \mathcal{K}_g
     (\mu, \bar{\varphi}_2)) \rangle = 0. \]
  We are going to see that the inequality $\langle \rho_{\ell'} (\lambda
  \hat{z}) J (\bar{\varphi}_1 - \bar{\varphi}_2), (- \Delta_{\hat{z}} + m^2)
  (\bar{\varphi}_1 - \bar{\varphi}_2) \rangle \geqslant C \int \rho_{\ell'}
  (\lambda \hat{z}) J (\bar{\varphi}_1 - \bar{\varphi}_2) (\bar{\varphi}_1 -
  \bar{\varphi}_2) \mathd \hat{z}$ holds for $\lambda$ small enough and some
  constant $C > 0$. Indeed let $f_1, f_2$ be two smooth functions then, for
  $\lambda$ small enough,
  \[ \begin{array}{l}
       \langle \rho_{\ell'} (\lambda \hat{z}) J (f_1 - f_2), (-
       \Delta_{\hat{z}} + m^2) (f_1 - f_2) \rangle\\
       \qquad \qquad = \int \rho_{\ell'} (\lambda \hat{z}) J (f_1 - f_2) (-
       \Delta_{\hat{z}} + m^2) (f_1 - f_2) \mathd \hat{z}\\
       \qquad \qquad = \int \rho_{\ell'} (\lambda \hat{z}) J' (f_1 - f_2) |
       \nabla f_1 - \nabla f_2 |^2 \mathd \hat{z} + \lambda \int J (f_1 - f_2)
       \cdot (\nabla f_1 - \nabla f_2) \mathd \hat{z}\\
       \qquad \hspace{4em} + m^2 \int \rho_{\ell'} (\lambda \hat{z}) J (f_1 -
       f_2) (f_1 - f_2) \mathd \hat{z}\\
       \qquad \qquad \geqslant - \lambda^2 \int (\Delta \rho_{\ell'} (\lambda
       \hat{z})) J^{- 1} (f_1 - f_2) \mathd \hat{z} + m^2 \int \rho_{\ell'}
       (\lambda \hat{z}) J (f_1 - f_2) (f_1 - f_2) \mathd \hat{z}\\
       \qquad \qquad \geqslant \int \left( m^2 - \left| \frac{\lambda^2 \Delta
       \rho_{\ell'}}{\rho_{\ell'}} \right| \right) \rho_{\ell'} (\lambda
       \hat{z}) J (f_1 - f_2) (f_1 - f_2) \mathd \hat{z}\\
       \hspace{4em} \geqslant C \int \rho_{\ell'} (\lambda \hat{z}) J (f_1 -
       f_2) (f_1 - f_2) \mathd \hat{z}
     \end{array} \]
  where $J^{- 1} (t) = \int_0^t J (\tau) \mathd \tau$ and we have used the
  fact that $J^{- 1} (t) \leqslant J (t) t$ since $J$ is increasing and the
  fact that we can choose $\lambda$ small enough such that the last inequality
  holds. This proves that
  \begin{equation}
    \langle \rho_{\ell'} (\lambda \hat{z}) J (f_1 - f_2), (- \Delta_{\hat{z}}
    + m^2) (f_1 - f_2) \rangle \geqslant C \int \rho_{\ell'} (\lambda \hat{z})
    J (f_1 - f_2) (f_1 - f_2) \mathd \hat{z} \label{eq:inequalityuniqueness},
  \end{equation}
  for $\lambda$ small enough, and some $C > 0$, for smooth function $f_1, f_2
  \in B^{2 + 2 - \delta'}_{p, p, \ell - \delta}$. Since the expressions
  $\langle \rho_{\ell'} (\lambda \hat{z}) J (f_1 - f_2), (- \Delta_{\hat{z}} +
  m^2) (f_1 - f_2) \rangle$ and $\int \rho_{\ell'} (\lambda \hat{z}) J (f_1 -
  f_2) (f_1 - f_2) \mathd \hat{z}$ are continuous for $f_1, f_2 \in B^{2 + 2 -
  \delta'}_{p, p, \ell - \delta}$ (with respect to the $B^{2 + 2 -
  \delta'}_{p, p, \ell - \delta}$-natural norm) we can extend
  inequality~{\eqref{eq:inequalityuniqueness}} to the case of general
  functions $f_1, f_2 \in B^{2 + 2 - \delta'}_{p, p, \ell - \delta}$ and so it
  holds in particular for $f_1 = \bar{\varphi}_1$ and $f_2 = \bar{\varphi}_2$.
  Furthermore we have that $\langle \rho_{\ell'} (\lambda \hat{z}) J
  (\bar{\varphi}_1 - \bar{\varphi}_2), (- \Delta_{\hat{z}} + m^2) (-
  \mathcal{K}_g (\mu, \bar{\varphi}_1) + \mathcal{K}_g (\mu, \bar{\varphi}_2))
  \rangle \geqslant 0$. Indeed by Lemma~\ref{lemma_multiplication} the product
  is associative and thus we obtain
  \begin{eqnarray}
    & \langle \rho_{\ell'} (\lambda \hat{z}) J (\bar{\varphi}_1 -
    \bar{\varphi}_2), (- \Delta_{\hat{z}} + m^2) (- \mathcal{K}_g (\mu,
    \bar{\varphi}_1) + \mathcal{K}_g (\mu, \bar{\varphi}_2)) \rangle = & 
    \nonumber\\
    & = \int \rho_{\ell'} (\lambda \hat{z}) (\alpha G (\alpha
    \bar{\varphi}_1) - \alpha G (\alpha \bar{\varphi}_1)) J (\bar{\varphi}_1 -
    \bar{\varphi}_2) \mathd \mu (\hat{z}) \geqslant 0, &  \nonumber
  \end{eqnarray}
  where we use that $\mathd \mu (\hat{z})$ is a positive measure and $(\alpha
  G (\alpha t_1) - \alpha G (\alpha t_2)) \cdot J (t_1 - t_2)$ is positive
  since both $G$ and $J$ are increasing functions and $J (0) = 0$. Using the
  previous inequalities we obtain that
  \[ \int \rho_{\ell'} (\lambda \hat{z}) J (\bar{\varphi}_1 -
     \bar{\varphi}_2) (\bar{\varphi}_1 - \bar{\varphi}_2) \mathd \hat{z}
     \leqslant 0 \]
  which holds only if $\bar{\varphi}_1 - \bar{\varphi}_2 = 0$ since $J$ is a
  strictly increasing function. 
\end{proof}

\begin{remark}
  Combining Lemma~\ref{lemma_exponentialexistence} and
  Lemma~\ref{lemma_exponentialuniqueness} we deduce that the map $\mu
  \longmapsto \bar{\varphi}$, associating with the measure $\mu$ the unique
  solution $\bar{\varphi}$ to equation~{\eqref{eq:K1}}, is continuous with
  respect to $\mu$. Indeed suppose that $\mu_n \rightarrow \mu$ in
  $\mathfrak{M}^s_{+, p, p, \ell}$ (and so in $B^s_{p, p, \ell}$), then by
  Lemma~\ref{lemma_multiplication}, we have that $\| \bar{\varphi}_n \|_{B^{s
  + 2}_{p, p, \ell}} \lesssim \sup \| \mu_n \|_{B^s_{p, p, \ell}}$, and so
  there exists a converging subsequence $\bar{\varphi}_{k_n}$, as $n
  \rightarrow \infty$. On the other hand, since $\mathcal{K}_g$ is continuous
  in both $\mu$ and $\varphi$, we get that $\bar{\varphi}_{k_n}$ converges to
  the unique solution $\varphi$ to equation~{\eqref{eq:K1}} associated with
  $\mu .$ Since the limit does not depend on the subsequence we have that
  $\varphi_n \rightarrow \varphi$ strongly in $B^{s + 2 - \delta'}_{p, p, \ell
  + \delta}$, which proves the continuity of the solution map.
  
  This continuity of the solution map with respect to the measure $\mu$ is
  similar to the continuity result obtained for solutions of singular SPDEs
  defined by the methods of paracontrolled calculus or regularity structure
  theory. 
\end{remark}

\begin{proof*}{Proof of Theorem~\ref{theorem_exponentialmain}}
  The existence and uniqueness of the solution to
  equation~{\eqref{eq:exponential2}} are proved in
  Lemma~\ref{lemma_exponentialexistence} and
  Lemma~\ref{lemma_exponentialuniqueness} considering the equation $\bar{\phi}
  = \mathcal{K} (\eta, \bar{\phi})$. What remains to prove is the convergence
  of $\bar{\phi}_{\varepsilon_n}$ to $\bar{\phi}$, as $\varepsilon_n
  \rightarrow 0$. Let $\epsilon_n \rightarrow 0$ be a sequence of positive
  numbers such that $\eta_{\epsilon_n} \rightarrow \eta$ almost surely in
  $\mathfrak{M}^s_{+, p, p, \ell}$ and let $w \in \mathcal{W}$ be such that
  $\eta_{\epsilon_n} (w) \rightarrow \eta (w)$ in $\mathfrak{M}^s_{+, p, p,
  \ell}$, as $\varepsilon \rightarrow 0$. We note that
  \begin{equation}
    \bar{\phi}_{\epsilon_n} (w) = \mathfrak{a}^{\asterisk 2}_{\epsilon_n}
    \asterisk (\mathcal{K}_g (\eta_{\epsilon_n} (w),
    \bar{\phi}_{\varepsilon_n} (w))) . \label{eq:exponentialmain2}
  \end{equation}
  From the equality~{\eqref{eq:exponentialmain2}} and
  Lemma~\ref{lemma_multiplication} we obtain that
  \[ \| \bar{\phi}_{\epsilon_n} (w) \|_{B^{s + 2}_{p, p, \ell}} \lesssim
     \sup_{n \in \mathbb{N}} \| \eta_{\varepsilon_n} \|_{B^s_{p, p, \ell}}, \]
  uniformly in $n$. This means that that there exists a subsequence
  $\bar{\phi}_{\epsilon_{k_n}} (w)$ converging to some $\bar{\varphi}$ in
  $B^{s + 2 - \delta'}_{p, p, \ell + \delta}$. On the other hand we have that
  \[ \bar{\varphi} = \lim_{n \rightarrow + \infty}
     \bar{\phi}_{\epsilon_{k_n}} (w) = \lim_{n \rightarrow + \infty}
     \mathfrak{a}^{\asterisk 2}_{\epsilon_n} \asterisk \left( \mathcal{K}_g
     \left( \eta_{\epsilon_{k_n}} (w), \bar{\phi}_{\varepsilon_{k_n}} (w)
     \right) \right) = \mathcal{K}_g (\eta (w), \bar{\varphi}) . \]
  Since, by Lemma~\ref{lemma_exponentialuniqueness}, equation~{\eqref{eq:K1}}
  has a unique solution we have that all the subsequence
  $\bar{\phi}_{\epsilon_{k_n}} (w)$ converge to the same $\bar{\varphi}$ and
  so $\bar{\varphi} = \bar{\phi} (w)$ (which is the unique solution to
  equation~{\eqref{eq:exponential2}} evaluated in $w \in \mathcal{W}$). Since
  the previous reasoning holds for almost every $w \in \mathcal{W}$ and since
  $\eta_{\epsilon_n} \rightarrow \eta_{}$ almost surely, we have that
  $\bar{\phi}_{\epsilon_n} \rightarrow \bar{\phi}$ in $B^{s + 2 - \delta'}_{p,
  p, \ell - \delta}$ almost surely, as $\varepsilon_n \rightarrow 0$.
\end{proof*}

\subsection{Dimensional reduction}

In this section we want to prove the reduction principle for
equation~{\eqref{eq:exponential1}}, i.e. that the random field $\phi (0, z) =
\mathcal{I} \xi (0, z) + \bar{\phi} (0, z)$ has the law $\kappa_g$ given by
expression~{\eqref{eq:exponentiallaw1}}.

Before proving the dimensional reduction for
equation~{\eqref{eq:exponential1}} we have to prove that the restriction of a
solution to a two dimensional hyperplane is a well defined operation.

\begin{lemma}
  \label{lemma_xi}There exists a version of the functional $x \longmapsto
  \mathcal{I} \xi (x, \cdot)$ which is continuous as a function from
  $\mathbb{R}^2$ into $\mathcal{C}^{0 -}_{\ell} (\mathbb{R}^2) \subset B^{0
  -}_{p, p, \ell + \delta}$. Furthermore for any $h \in C^{\infty}_0
  (\mathbb{R}^2)$ the sequence of random variables $\int_{\mathbb{R}^2}
  \mathfrak{a}_{\epsilon} \asterisk (\mathcal{I} \xi) (z) h (z) \mathd z$
  converges to $\langle \mathcal{I} \xi (0, \cdot), h \rangle$ in $L^2 (\mathd
  \mu)$, as $\varepsilon \rightarrow 0$.
\end{lemma}

\begin{proof}
  In order to prove that $x \longmapsto \mathcal{I} \xi (x, \cdot)$ admits a
  continuous version we prove that for any smooth function $h : \mathbb{R}^2
  \rightarrow \mathbb{R}$ such that $(- \Delta_z + m^2)^{- 1 / 2 + \epsilon}
  (h) \in L^2 (\mathbb{R}^2)$, with $\varepsilon > 0$, we have
  \[ \mathbb{E} [| \langle \mathcal{I} \xi (x, \cdot) - \mathcal{I} \xi (y,
     \cdot), h \rangle |^2] \lesssim | x - y |^{\epsilon} \| (- \Delta_z +
     m^2)^{- 1 / 2 + \epsilon} (h) \|_{L^2 (\mathbb{R}^2)}, \]
  where the constants implied by $\lesssim$ do not depend on $h, x, y$. This
  result, exploiting hypercontractivity and a version of Kolmogorov continuity
  criterion for multidimensional random fields, implies the continuity of
  $\mathcal{I} \xi (x, \cdot)$ with respect to $x \in \mathbb{R}^2$. We note
  that, for any $x, y \in \mathbb{R}^2$:
  \[ \mathbb{E} [| \langle \mathcal{I} \xi (x, \cdot) - \mathcal{I} \xi (y,
     \cdot), h \rangle |^2] = 2 \int_{\mathbb{R}^2} \frac{1 - e^{i k_1 \cdot
     (x - y)}}{(k_1^2 + k_2^2 + m^2)^{1 + 2 \epsilon}} \frac{\hat{h}
     (k_2)^2}{(k_1^2 + k_2^2 + m^2)^{1 - 2 \epsilon}} \mathd k_1 \mathd k_2 .
  \]
  Now we observe that $| 1 - e^{i k_1 \cdot (x - y)} | \leqslant 2 | x - y
  |^{\frac{p - 1}{p}} | k_1 |^{\frac{p - 1}{p}}$ for any $p > 1$. If we choose
  $p = \frac{1}{1 - \epsilon}$, $0 < \epsilon < 1$, we obtain the claim.
  
  In order to prove that $\int_{\mathbb{R}^2} \mathfrak{a}_{\epsilon}
  \asterisk (\mathcal{I} \xi) (z) h (z) \mathd z$ converges, as $\varepsilon
  \rightarrow 0$, to $\langle \mathcal{I} \xi (0, \cdot), h \rangle$ in $L^2
  (\mathd \mu)$, we note that the distribution $\mathcal{I} (\delta_0 (x)
  \cdot h (z))$ belongs to $L^2 (\mathbb{R}^4)$, implying that
  \[ \int_{\mathbb{R}^2} \mathfrak{a}_{\epsilon} \asterisk (\mathcal{I} \xi)
     (z) h (z) \mathd z = \delta [\mathcal{I} (\delta_0 (x) \cdot
     \mathfrak{a}_{\epsilon} \asterisk h (z))] \qquad \langle \mathcal{I} \xi
     (0, \cdot), h \rangle = \delta [\mathcal{I} (\delta_0 (x) \cdot h (z))],
  \]
  where $\delta$ denote the Skorohod integral with respect to the white noise
  $\xi$. Since $\mathcal{I} (\delta_0 (x) \cdot \mathfrak{a}_{\epsilon}
  \asterisk h (z))$ converges to $\mathcal{I} (\delta_0 (x) \cdot h (z))$ in
  $L^2 (\mathbb{R}^4)$ also this claim follows. 
\end{proof}

\begin{lemma}
  \label{lemma_trace}The operator $T_0 : C^0_{\ell} (\mathbb{R}^4) \rightarrow
  C^0_{\ell} (\mathbb{R}^2)$ given by $T_0 (h) (\cdot) = h (0, \cdot)$ can be
  uniquely extended in a continuous way as an operator from $B^{s + 2 -
  \delta'}_{p, p, \ell + \delta} (\mathbb{R}^4)$ into $B^{s + 2 - 2 / p -
  \delta'}_{p, p, \ell + \delta} (\mathbb{R}^2)$, when $s, p$ satisfies the
  hypotheses of Theorem~\ref{theorem_exponentialmain} and $\delta' > 0$ is
  small enough.
\end{lemma}

\begin{proof}
  For the proof of this fact we refer to Section ~4.4.1 and Section ~4.4.2
  of~{\cite{Triebel1992}}, see also Section~18.1 of~{\cite{Triebel1997}}.
\end{proof}

\begin{theorem}
  \label{theorem_dimensionalreduction1}Let $\phi$be the unique solution to
  equation~{\eqref{eq:exponential1}} then, if $| \alpha | < 4 \sqrt{2} \pi$,
  there exists $\ell > 0$ such that for any measurable and bounded function
  $F$ on $B^{0 -}_{p, p, \ell + \delta} (\mathbb{R}^2)$ we have
  \begin{equation}
    \mathbb{E} [F (\phi (0, \cdot))] = \int_{B^{0 -}_{p, p, \ell + \delta}
    (\mathbb{R}^2)} F (\omega) \mathd \kappa_g (\omega) . \label{eq:final}
  \end{equation}
\end{theorem}

\begin{proof}
  We prove equation~{\eqref{eq:final}} for the case in which $F (\omega) =
  \tilde{F} (\langle \omega, f_1 \rangle, \ldots, \langle \omega, f_n
  \rangle)$ where $\tilde{F}$ is a bounded continuous function, and $f_1,
  \ldots, f_n \in C^{\infty}_0 (\mathbb{R}^2)$. Since the functions of the
  previous form generate all the $\sigma$-algebra of Borel measurable
  functions on $B^{0 -}_{p, p, \ell + \delta}$, proving the theorem for
  functions of the previous form is equivalent to prove the theorem in
  general.
  
  Since equation~{\eqref{eq:exponential5}} holds, we have only to prove that
  $\mathbb{E} [F (\phi_{\epsilon_n} (0, \cdot))] \rightarrow \mathbb{E} [F
  (\phi (0, \cdot))]$, as $\epsilon_n \rightarrow 0$ and where $\epsilon_n$ is
  any subsequence such that $\bar{\phi}_{\epsilon_n} \rightarrow \bar{\phi}$
  almost surely (whose existence is proved in
  Theorem~\ref{theorem_exponentialmain}) and $\int_{B^{0 -}_{p, p, \ell +
  \delta} (\mathbb{R}^2)} F (\omega) \mathd \kappa_{\epsilon} (\omega)
  \rightarrow \int_{B^{0 -}_{p, p, \ell + \delta} (\mathbb{R}^2)} F (\omega)
  \mathd \kappa (\omega)$, as $\epsilon_n \rightarrow 0$.
  
  The first convergence follows from the fact that $\phi_{\epsilon} =
  \bar{\phi}_{\epsilon} + \mathcal{\mathfrak{a_{\epsilon} (\mathcal{I}
  \xi)}}$. Indeed, by Theorem~\ref{theorem_exponentialmain}, as $\epsilon_n
  \rightarrow 0$,
  \[ \langle \bar{\phi}_{\epsilon_n} (0, \cdot), f_i \rangle = \langle
     \bar{\phi}_{\varepsilon_n}, T^{\asterisk}_0 (f_i) \rangle \rightarrow
     \langle \bar{\phi}, T^{\asterisk}_0 (f_i) \rangle = \langle \bar{\phi}
     (0, \cdot), f_i \rangle, \]
  for any $i = 1, \ldots, n$, almost surely and where the continuity of $T_0,$
  proved in Lemma~\ref{lemma_trace}, is used. On the other hand, by
  Lemma~\ref{lemma_xi}, we have $\langle \mathfrak{a}_{\epsilon} (\mathcal{I}
  \xi) (0, \cdot), f_i \rangle \rightarrow \langle \mathfrak{} \mathcal{I} \xi
  (0, \cdot), f_i \rangle$ in probability, as $\varepsilon_n \rightarrow 0$,
  and this implies that $\langle \phi_{\epsilon_n} (0, \cdot), f_i \rangle
  \rightarrow \langle \phi (0, \cdot), f_i \rangle$ in probability, for any $i
  = 1, \ldots, n$.
  
  Since the convergence in probability implies the one in distribution we get
  $\mathbb{E} [F (\phi_{\epsilon_n} (0, \cdot))] \rightarrow \mathbb{E} [F
  (\phi (0, \cdot))]$. Finally since $\kappa_{\epsilon_n}$ converges weakly to
  $\kappa$, as $\epsilon_n \rightarrow 0$, the thesis follows. \ 
\end{proof}

Now we want to discuss what happens if we remove the cut-off $g$ in
equation~{\eqref{eq:exponential1}} so we consider the equation
\begin{equation}
  (- \Delta_{\hat{z}} + m^2) (\phi) + \alpha \exp (\alpha \phi - \infty) = \xi
  . \label{eq:exponentialcutoff}
\end{equation}
In order to distinguish between the solution to
equation~{\eqref{eq:exponential1}} and equation~{\eqref{eq:exponentialcutoff}}
we denote by $\phi_g$ the solution of the former and by $\phi$ the solution of
the latter. We use also the symbols $\bar{\phi}_g = \phi_g - \mathcal{I} \xi$
and $\bar{\phi} = \phi - \mathcal{I} \xi$.

\begin{proposition}
  \label{proposition_g1}There exists an unique solution $\phi$ to
  equation~{\eqref{eq:exponentialcutoff}}. \ Furthermore, for any affine
  transformation $\Phi : \mathbb{R}^4 \rightarrow \mathbb{R}^4$ in the
  Euclidean group of $\mathbb{R}^4$, the random field $\Phi_{\asterisk} (\phi)
  (\hat{z}) = \phi (\Phi (\hat{z}))$ has the same law of $\phi (\hat{z})$.
\end{proposition}

\begin{proof}
  The existence and uniqueness for the solution to
  equation~{\eqref{eq:exponentialcutoff}} can be proven as in
  Lemma~\ref{lemma_exponentialexistence} and
  Lemma~\ref{lemma_exponentialuniqueness}, since the estimates used to prove
  those lemmas do not depend on $g$.
  
  The invariance of the law of the solutions with respect to affine
  transformations in the Euclidean group of $\mathbb{R}^4$ follows from the
  invariance (in law) of the white noise $\xi$ and of the left-hand-side of
  equation~{\eqref{eq:exponentialcutoff}} with respect to translations and
  rotations and from the uniqueness of the solution to
  equation~{\eqref{eq:exponentialcutoff}}.
\end{proof}

\begin{theorem}
  \label{theorem_dimensionalreduction2}If $g_n$ is an (increasing) sequence of
  cut-offs with compact support such that $g_n \uparrow 1$ then
  $\bar{\phi}_{g_n} \rightarrow \bar{\phi}$ in $B^{s + 2 - \delta'}_{p, p,
  \ell + \delta}$ (which means that $\phi_{g_n} \rightarrow \phi$ in $B^{0
  -}_{p, p, \ell + \delta} (\mathbb{R}^4)$). This means that the sequence of
  probability measures $\kappa_{g_n}$ on $B^{0 -}_{p, p, \ell + \delta}
  (\mathbb{R}^2)$ converges weakly, as $n \rightarrow \infty$, to a
  probability measure $\kappa$ which is invariant with respect to the natural
  action of the Euclidean group of $\mathbb{R}^2$ on $B^{0 -}_{p, p, \ell +
  \delta} (\mathbb{R}^2)$ and does not depend on the sequence of cut-offs
  $g_n$.
\end{theorem}

\begin{proof}
  The proof is essentially based on the fact that
  \[ \| \bar{\phi}_g \|_{B^{s + 2}_{p, p, \ell}} \lesssim \| \eta \|_{B^s_{p,
     p, \ell}} \]
  uniformly on $g$. From the previous inequality and some reasoning similar to
  the ones used in the proof of Theorem~\ref{theorem_exponentialmain} the
  convergence follows. The properties of the limit measure $\kappa$ follow
  from the same properties of the solution $\phi$ to
  equation~{\eqref{eq:exponentialcutoff}} proved in
  Proposition~\ref{proposition_g1}.
\end{proof}

\section{Elliptic quantization of the $P (\varphi)_2$
model}\label{section_power}

In this section we discuss the elliptic stochastic quantization of the $P
(\varphi)_2$ model where $P$ is a polynomial of even degree and satisfying
Hypothesis~QC. In order to avoid technical details we consider only the case
$P (\varphi) = \frac{\varphi^{2 n}}{2 n}$ for $n \in \mathbb{N}$, the general
case being then a straightforward generalisation. We consider the equation
\begin{equation}
  (- \Delta_x - \Delta_z + m^2) (\phi) + f (x) \phi^{\diamond 2 n - 1} = \xi
  \label{eq:power1}
\end{equation}
where $z \in M =\mathbb{T}^2$, $\diamond$ stands for the Wick product and
$\xi$ is a $\mathbb{R}^2 \times \mathbb{T}^2$ white noise.
Equation~{\eqref{eq:power1}} can be better understood if we consider the
equation for $\bar{\phi} : = \phi - \mathcal{I} \xi$, the usual Da
Prato-Debussche trick, obtaining the equation that is
\begin{equation}
  (- \Delta_x - \Delta_z + m^2) (\bar{\phi}) + \sum_{k = 0}^{2 n - 1} \binom{2
  n - 1}{k} f (x) \cdot \mathcal{I} \xi^{\diamond k} \cdot \bar{\phi}^{2 n - 1
  - k} = 0. \label{eq:power2}
\end{equation}
Equation~{\eqref{eq:power2}} is expected to be well defined, since we expect
that the solution $\bar{\phi}$ is in $H^1 (\mathbb{R}^2 \times \mathbb{T}^2) =
W^{1, 2} (\mathbb{R}^2 \times \mathbb{T}^2)$, in $L^{2 n}_{f^{1 / 2 n}}
(\mathbb{R}^2 \times \mathbb{T}^2)$ (where $L^{2 n}_{f^{1 / 2 n}}$ is the
weighted $L^{2 n}$ space with respect to the space weight $f^{\frac{1}{2 n}}$)
and in $B^{2 - \delta}_{\mathfrak{p}, \mathfrak{p}, \ell} (\mathbb{R}^2 \times
\mathbb{T}^2)$, where $\mathfrak{p =} \frac{2 n - 1}{2 n}$. Furthermore it is
well known that $\mathcal{I} \xi^{\diamond k} \in \mathcal{C}^{-
\delta}_{\ell} = B^{- \delta}_{\infty, \infty, \ell} (\mathbb{R}^2 \times
\mathbb{T}^2)$ for any $\ell > 0$ (see~{\cite{GH18}}). In general we expect
that equation~{\eqref{eq:power2}}, and so equation~{\eqref{eq:power1}}, for
any realization of the noise $\xi$ admits multiple solutions. So we need a
notion of weak solution to equation~{\eqref{eq:power2}}.

First of all we consider a fixed probability space $\mathcal{W}^e =
(\mathcal{C}^{- \delta}_{\ell} (\mathbb{R}^2 \times \mathbb{T}^2))^{2 n}
\times \mathfrak{W}$, where $\mathfrak{W} = H^{1 - \delta'}_{\ell'}
(\mathbb{R}^2 \times \mathbb{T}^2) \cap L^{2 n - \delta'}_{f^{1 / 2 n +
\delta'}} (\mathbb{R}^2 \times \mathbb{T}^2) \cap B^{2 -
\delta'}_{\mathfrak{p}, \mathfrak{p}, \ell} (\mathbb{R}^2 \times
\mathbb{T}^2)$, and where $\delta, \delta' > 0$ are small enough, that is the
space where $(\mathcal{I} \xi, \mathcal{I} \xi^{\diamond 2}, \ldots,
\mathcal{I} \xi^{\diamond 2 n}, \bar{\phi}) \in \mathcal{W}^e$ are defined.

\begin{definition}
  \label{definition_we}A probability measure $\nu^e$ on $\mathcal{W}^e$ is a
  weak solution to equation~{\eqref{eq:power2}} if the projection of $\nu^e$
  on $[\mathcal{C}^{- \delta}_{\ell} (\mathbb{R}^2 \times \mathbb{T}^2)]^{2
  n}$ gives the law of $(\mathcal{I} \xi, \mathcal{I} \xi^{\diamond 2},
  \ldots, \mathcal{I} \xi^{\diamond 2 n - 1})$ (which is the law of a Gaussian
  noise with covariance $(- \Delta_x - \Delta_z + m^2)^{- 2}$ and its Wick
  powers) and it is supported on the set of solutions to the equation
  \begin{equation}
    (- \Delta_x - \Delta_z + m^2) (\bar{\theta}) + \sum_{k = 0}^{2 n - 1}
    \binom{2 n - 1}{k} f (x) \cdot \sigma_k \cdot \bar{\theta}^{2 n - 1 - k} =
    0, \label{eq:power3}
  \end{equation}
  where $x \in \mathbb{R}^2$, $z \in \mathbb{T}^2$, $(\sigma_1, \ldots,
  \sigma_{2 n}, \bar{\theta}) \in \mathcal{W}^e$ and $\sigma_0 = 1$, defined
  on $\mathcal{W}^e$.
  
  The weak solution $\nu$ to equation~{\eqref{eq:power1}} associated with the
  weak solution $\nu^e$ to equation~{\eqref{eq:power2}}, is the probability
  law on $\mathcal{C}^{- \delta}_{\ell} (\mathbb{R}^2 \times \mathbb{T}^2) +
  \mathfrak{W}$ given by the push-forward of $\nu^e$ with respect to the map
  $\sigma_1 + \theta$ defined on $\mathcal{W}^e$.
\end{definition}

We introduce a modified equation on $\mathfrak{W}$ given by
\begin{equation}
  (- \Delta_x - \Delta_z + m^2) (\bar{\theta}) + \mathfrak{a}^{\asterisk
  2}_{\epsilon} \asterisk \left[ \sum_{k = 0}^{2 n - 1} \binom{2 n - 1}{k} f
  (x) \cdot \sigma_k \cdot \bar{\theta}^{2 n - 1 - k} \right] = 0
  \label{eq:power4}
\end{equation}
where $\mathfrak{a}_{\epsilon}$ is some regular enough mollifier on
$\mathbb{T}^2$ such that the operator $\mathcal{A}_{\epsilon} =
\mathfrak{a}_{\epsilon} \asterisk$ satisfies Hypothesis~H$\mathcal{A}$ (for
example we can take $\mathfrak{a}_{\epsilon}$ as the Green function associated
with the operator $(- \epsilon \Delta_z + 1)^{- k}$ for $k$ large enough),
$\varepsilon > 0$. \

\begin{lemma}
  \label{lemma_power1}Let $\bar{\theta} \in \mathfrak{W}$ be a solution to
  equation~{\eqref{eq:power4}} then $\bar{\theta} \in H^1 (\mathbb{R}^2 \times
  \mathbb{T}^2) \cap B^{2 - \delta''}_{\mathfrak{p}, \mathfrak{p}}
  (\mathbb{R}^2 \times \mathbb{T}^2) \cap L^{2 n}_{f^{1 / 2 n}} (\mathbb{R}^2
  \times \mathbb{T}^2)$ (with $\delta'' < \delta'$ being $\delta'$ as in the
  definition of $\mathcal{W}^e$) and
  \begin{eqnarray}
    \| \bar{\theta} \|_{H^1}^2 + \| \bar{\theta} \|_{L^{2 n}_{f^{1 / 2 n}}}^{2
    n} & \lesssim & \left[ \sum_{k = 1}^{2 n} \| \sigma_k \|_{\mathcal{C}^{-
    \delta}_{\ell}} \right]^{\beta_1} \label{eq:P1} \\
    \| \bar{\theta} \|_{B^{2 - \delta''}_{\mathfrak{p}, \mathfrak{p}}} &
    \lesssim & \left[ \| \theta \|_{H^1} + \| \theta \|_{L^{2 n}_f} + \sum_{k
    = 1}^{2 n} \| \sigma_k \|_{\mathcal{C}^{- \delta}_{\ell}}
    \right]^{\beta_2} \label{eq:P2} 
  \end{eqnarray}
  where $\beta_1, \beta_2 \in \mathbb{R}_+$ and where $\beta_1, \beta_2$ and
  the constants implied by the symbol $\lesssim$ do not depend on $\epsilon$
  and $\sigma_k$.
\end{lemma}

\begin{proof}
  Let $p > 1, 0 < \gamma < 1$ and $\epsilon'' > 0$ be such that $\frac{1}{p} =
  \frac{\gamma}{2 n - \delta'} + \frac{1 - \gamma}{2} (1 - \delta')$ and
  $\frac{2 n - 1}{p} + \epsilon'' = 1$. Then by Proposition
  \ref{proposition_inclusion} and Proposition \ref{proposition_interpolation}
  it is simple to see that $\bar{\theta} \in B^{(1 - \gamma) (1 - \delta)}_{p,
  2, f^{1 / 2 n}} (\mathbb{R}^2 \times \mathbb{T}^2)$ and
  \[ \| \bar{\theta} \|_{B^{(1 - \gamma)}_{p, 2, f^{1 / 2 n}}} \lesssim \|
     \bar{\theta} \|_{L^{2 n - \delta'}_{f^{1 / 2 n}}}^{\gamma} \cdot \|
     \bar{\theta} \|^{1 - \gamma}_{H^{1 - \delta'}_{\ell'}} . \]
  This means that
  \begin{eqnarray}
    \| (1 + | x |)^{\ell} (\mathcal{A}_{\epsilon}^{- 2} \bar{\theta}) \|_{B^{2
    - \delta''}_{\mathfrak{p}, \mathfrak{p}}} & \lesssim & \left\| (1 + | x
    |)^{\ell} \left[ \sum_{k = 0}^{2 n} \binom{2 n}{k} f (x) \cdot \sigma_k
    \cdot \bar{\theta}^{2 n - k} \right] \right\|_{B^{-
    \delta''}_{\mathfrak{p, \mathfrak{p}}}} \nonumber\\
    & \lesssim & \sum_{k = 0}^{2 n - 1} \| (1 + | x |)^{\ell} f (x) \cdot
    \sigma_k \cdot \bar{\theta}^{2 n - k} \|_{B^{- \delta''}_{\mathfrak{p,
    \mathfrak{p}}}} \nonumber\\
    & \lesssim & \sum_{k = 0}^{2 n - 1} \| (1 + | x |)^{\ell} f^{1 / 2 n + (2
    n - 1) \delta'} \cdot \sigma_k \|_{\mathcal{C}^{- \delta}} \times
    \nonumber\\
    &  & \times \left[ \| \bar{\theta} \|_{L^{2 n - \delta'}_{f^{1 / 2 n +
    \delta'}}}^{\gamma} \cdot \| \bar{\theta} \|^{1 - \gamma}_{H^{1 -
    \delta'}_{\ell'}} \right]^{2 n - 1 - k} \nonumber\\
    & \lesssim & \left[ \| \bar{\theta} \|_{H^{1 - \delta'}_{\ell'}} + \|
    \bar{\theta} \|_{L^{2 n - \delta'}_{f^{1 / 2 n + \delta'}}} \right.
    \nonumber\\
    &  & \left. + \sum_{k = 1}^{2 n - 1} \| (1 + | x |)^{\ell} f^{1 / 2 n +
    (2 n - 1) \delta'} \sigma_k \|_{\mathcal{C}^{- \delta}} \right]^{\beta_2}
    \nonumber
  \end{eqnarray}
  where $\delta > 0$ such that $\frac{1 - \gamma}{2} (1 - \delta') - \delta >
  0$. Since
  \begin{equation}
    \| \bar{\theta} \|_{B^{2 - \delta''}_{\mathfrak{p}, \mathfrak{p}}}
    \lesssim \| (1 + | x |)^{\ell} (\mathcal{A}_{\epsilon}^{- 2} \bar{\theta})
    \|_{B^{2 - \delta''}_{\mathfrak{p}, \mathfrak{p}}}, \label{eq:p2}
  \end{equation}
  uniformly in $\epsilon$, this implies that $\bar{\theta} \in B^{2 -
  \delta''}_{\mathfrak{p}, \mathfrak{p}}$ and also that $\bar{\theta} \in H^1$
  since $\| \bar{\theta} \|_{H^1}^2 \lesssim \| \bar{\theta} \cdot (- \Delta +
  m^2) (\bar{\theta}) \|_{L^1} \leqslant \| \bar{\theta} \|_{H^{1 -
  \delta'}_{\ell}} \cdot \| (1 + | x |)^{\ell} \bar{\theta} \|_{B^{2 -
  \delta''}_{\mathfrak{p}, \mathfrak{p}}}$. Using the fact that $\bar{\theta}
  \in H^1$ and the previous reasoning we obtain inequality~{\eqref{eq:P2}}, as
  we shall now explain.
  
  Let us multiply both sides equation~{\eqref{eq:power4}} by
  $\mathcal{A}_{\epsilon}^{- 2} (\bar{\theta})$ (operation which is well
  defined since $\| (1 + | x |)^{\ell} (\mathcal{A}_{\epsilon}^{- 2}
  \bar{\theta}) \|_{B^{2 - \delta''}_{\mathfrak{p}, \mathfrak{p}}}$ is
  bounded) and take the integral over $\mathbb{R}^2 \times \mathbb{T}^2$
  obtaining
  \[ \| \mathcal{A}_{\epsilon}^{- 1} (\bar{\theta}) \|_{H^1}^2 \leqslant - \|
     \bar{\theta} \|_{L^{2 n}_{f^{1 / 2 n}}}^{2 n} + C \sum_{k = 0}^{2 n - 2}
     \| f^{1 / 2 n - (2 n - 1) \delta'} \sigma_k \|_{\mathcal{C}^{- \delta}}
     \left[ \| \bar{\theta} \|_{L^{2 n - \delta'}_{f^{1 / 2 n +
     \delta'}}}^{\gamma} \| \bar{\theta} \|^{1 - \gamma}_{H^{1 -
     \delta'}_{\ell'}} \right]^{2 n - k} \]
  for some constant $C > 0$ depending only on $n$. Using that $\frac{2 n -
  1}{p} + \epsilon'' = 1$, the fact that $\| \bar{\theta} \|_{H^1} \lesssim \|
  \mathcal{A}_{\epsilon}^{- 1} (\bar{\theta}) \|_{H^1}^2$, $\| \bar{\theta}
  \|_{L^{2 n - \delta'}_{f^{1 / 2 n + \delta'}}} \lesssim \| \bar{\theta}
  \|_{L^{2 n}_{f^{1 / 2 n}}}$ and Young's inequality for products we obtain
  \begin{equation}
    \| \mathcal{} \theta \|_{H^1}^2 + \| \theta \|_{L^{2 n}_{f^{1 / 2 n}}}^{2
    n} \lesssim \sum_{k = 0}^{2 n - 2} \| f^{1 / 2 n - \delta' (2 n - 1)}
    \sigma_k \|_{\mathcal{C}^{- \delta}}^{1 / \epsilon} . \label{eq:p1}
  \end{equation}
  Inserting now inequality {\eqref{eq:p1}} into inequality {\eqref{eq:p2}} we
  obtain the thesis.
\end{proof}

An easy consequence of the previous lemma is the following one.

\begin{lemma}
  \label{lemma_powercompact}Let $\mathcal{F}_{\epsilon} : (\mathcal{C}^{-
  \delta}_{\ell})^{2 n - 1} \rightarrow \mathcal{P} (\mathfrak{W})$,
  $\varepsilon > 0$, be the set-valued map associating $(\sigma_1, \ldots,
  \sigma_{2 n - 1})$ with the set of solutions to equation~{\eqref{eq:power4}}
  in $\mathfrak{W}$, if $K \subset (\mathcal{C}^{- \delta}_{\ell})^{2 n - 1}$
  is a compact set then $\bigcup_{\epsilon < 1} \mathcal{F}_{\epsilon} (K)$ is
  precompact in $\mathfrak{W}$.
\end{lemma}

\begin{proof}
  The proof of the lemma consists only in noticing that a consequence of
  Proposition \ref{proposition_inclusion} is that the inclusion map $i : H^1
  (\mathbb{R}^2 \times \mathbb{T}^2) \cap B^{2 - \delta''}_{\mathfrak{p},
  \mathfrak{p}} (\mathbb{R}^2 \times \mathbb{T}^2) \cap L^{2 n}_{f^{1 / 2 n}}
  (\mathbb{R}^2 \times \mathbb{T}^2) \rightarrow \mathfrak{W}$ is compact.
  Then the thesis follows from Lemma~\ref{lemma_power1}.
\end{proof}

Hereafter we denote by $\mathfrak{P :} \mathcal{W}^e \rightarrow
(\mathcal{C}^{- \delta}_{\ell} (\mathbb{R}^2 \times \mathbb{T}^2))^{2 n}$ the
natural projection of the Cartesian product $\mathcal{W}^e = (\mathcal{C}^{-
\delta}_{\ell} (\mathbb{R}^2 \times \mathbb{T}^2))^{2 n} \times \mathfrak{W}$
on the first component.

\begin{lemma}
  Let $\nu_{\epsilon}^e$ be a sequence of weak solutions to
  equation~{\eqref{eq:power4}} and suppose that $\mathfrak{P}_{\asterisk}
  (\nu_{\epsilon}^e)$ is tight. Then $\nu_{\epsilon}^e$ is tight. Furthermore
  suppose that $\mathfrak{P}_{\asterisk} (\nu_{\epsilon}^e)$ weakly converges
  to the law of $(\mathcal{I} \xi, \mathcal{I} \xi^{\diamond 2}, \ldots,
  \mathcal{I} \xi^{\diamond 2 n - 1})$ as $\epsilon \rightarrow 0$, then any
  convergent subsequence $\nu_{\epsilon_n}^e$ converges, as $\varepsilon_n
  \rightarrow 0$, to a weak solution $\nu^e$ to equation~{\eqref{eq:power2}}.
\end{lemma}

\begin{proof}
  The proof of the tightness of $\nu_{\epsilon}^e$ is similar to the one of
  Lemma~\ref{lemma_reduced3}, where Lemma~\ref{lemma_reduced2} is replaced by
  Lemma~\ref{lemma_powercompact}.
  
  Since $\mathfrak{P}_{\asterisk} (\nu_{\epsilon}^e)$ weakly converges to the
  law of $(\mathcal{I} \xi, \mathcal{I} \xi^{\diamond 2}, \ldots, \mathcal{I}
  \xi^{\diamond 2 n - 1})$ and so $\mathfrak{P}_{\asterisk} (\nu^e)$ has the
  same law as $(\mathcal{I} \xi, \mathcal{I} \xi^{\diamond 2}, \ldots,
  \mathcal{I} \xi^{\diamond 2 n - 1})$, the proof of the fact that $\nu^e$ is
  a weak solution to equation~{\eqref{eq:power2}} is equivalent to prove that
  $\nu^e$ is such that for any $C^1$ bounded function $F$ \ from $B^{-
  \delta}_{\mathfrak{p}, \mathfrak{p}}$ into $\mathbb{R}$ with bounded
  derivative we have
  \begin{equation}
    \int F \left( (- \Delta_x - \Delta_z + m^2) (\bar{\theta}) + \mathfrak{}
    \sum_{k = 0}^{2 n - 1} \binom{2 n - 1}{k} f (x) \sigma_k \bar{\theta}^{2 n
    - k} \right) \mathd \nu^e = \int F (E) \mathd \nu^e = F (0) .
    \label{eq:weak1}
  \end{equation}
  We know that $\int F (E_{\epsilon}) \mathd \nu_{\epsilon}^e = F (0)$ with
  \[ E_{\epsilon} = (- \Delta_x - \Delta_z + m^2) (\bar{\theta}) +
     \mathfrak{a}^{\asterisk 2}_{\epsilon} \asterisk \left( \sum_{k = 0}^{2 n
     - 1} \binom{2 n - 1}{k} f (x) \cdot \sigma_k \cdot \bar{\theta}^{2 n - k}
     \right) . \]
  On the other hand $E_{\epsilon}$ converges to $E$ uniformly on compact sets
  (since $\mathfrak{a}_{\epsilon}^{\asterisk 2} \asterisk$ strongly converges
  as an operator to the identity on $B^{- \delta}_{\mathfrak{p},
  \mathfrak{p}}$), which implies that $F \circ E_{\epsilon}$ converges to $F
  \circ E$ uniformly on compact sets, as $\varepsilon \rightarrow 0$.
  
  This fact and the boundedness of $F$ and of its derivatives implies (using a
  reasoning similar to the one of Lemma~\ref{lemma_reduced4}, see also Lemma~2
  in~{\cite{Albeverio2018elliptic}}) that $F \circ E_{\epsilon} \mathd
  \nu_{\epsilon}^e$ weakly converges to $F \circ E \mathd \nu^e$, as
  $\varepsilon \rightarrow 0$. This implies equation~{\eqref{eq:weak1}} and
  thus the thesis of the lemma. \ \ 
\end{proof}

\begin{theorem}
  \label{theorem_polynomialreduction}There exists at least one weak solution
  $\nu^e$ to equation~{\eqref{eq:power2}} such that, for any measurable
  bounded function $F : B^{- \delta}_{\mathfrak{p}, \mathfrak{p}, \ell}
  (\mathbb{T}^2) \rightarrow \mathbb{R}$ (where $\delta > 0$ and $\ell > 0$ as
  in the definition of $\mathcal{W}^e$ and \ $\mathfrak{p} = \frac{2 n - 1}{2
  n}$),
  \begin{equation}
    \int_{\mathcal{W}^e} F (\phi (0, \cdot)) \tilde{\Upsilon}_f (\sigma,
    \bar{\theta}) \mathd \nu^e (\sigma, \bar{\theta}) = Z_f \int_{B^{-
    \delta}_{\mathfrak{p}, \mathfrak{p}, \ell} (\mathbb{T}^2)} F (\omega)
    \mathd \kappa (\omega) \label{eq:powermain}
  \end{equation}
  where $\phi = \sigma_1 + \bar{\theta}$ (where $\sigma_1$ and $\bar{\theta}$
  are as in Definition \ref{definition_we}),
  \[ \tilde{\Upsilon}_f (\sigma, \bar{\theta}) = \exp \left( \sum_{k = 0}^{2
     n} \binom{2 n}{k} \langle \sigma_k \bar{\theta}^{2 m - k}, f' \rangle
     \right), \]
  \[ Z_f = \int_{\mathcal{W}^e} \tilde{\Upsilon}_f (\sigma, \bar{\theta})
     \mathd \nu^e (\sigma, \bar{\theta}) \]
  and $\kappa$ is given by
  \[ \frac{\mathd \kappa}{\mathd \mu^{- \Delta_z}} = \frac{\exp \left(
     \int_{\mathbb{T}^2} \frac{\omega^{\diamond 2 m} (z)}{2 m} \mathd z
     \right)}{Z_{\kappa}}, \]
  where $\mu^{- \Delta_z}$ is the Gaussian with covariance given
  by~{\eqref{eq:covariance1}} with $\mathfrak{L} = - \Delta_z$ and
  $\mathcal{A} = \tmop{id}$.
\end{theorem}

\begin{lemma}
  \label{lemma_uniform}Let $\nu^e_{\epsilon}$ be a sequence of weak solutions
  to equation~{\eqref{eq:power4}} such that $\mathfrak{P}_{\asterisk}
  (\nu_{\epsilon}^e) \sim (\mathcal{\mathfrak{a}}_{\epsilon} \asterisk
  \mathcal{I} \xi, \ldots, (\mathcal{\mathfrak{a}}_{\epsilon} \asterisk
  \mathcal{I} \xi)^{\diamond 2 n - 1})$ (where $\sim$ means ``with the same
  law'') then
  \[ \int \exp \left( p \sum_{k = 0}^{2 n} \binom{2 n}{k} \langle \sigma_k,
     f'  \bar{\theta}^{2 n - k} \rangle \right) \mathd \nu_{\epsilon}^e < C_p
  \]
  form some constant $C_p$ uniform in $\epsilon$ small enough and depending on
  $p \geq 1$.
\end{lemma}

\begin{proof*}{Proof}
  We want to use Nelson's trick to prove the theorem (see Chapter~V
  of~{\cite{Simon1974}}). Then we put $E (\sigma_k, \bar{\theta}) = \sum_{k =
  0}^{2 n} \binom{2 n}{k} \langle \sigma_k, f'  \bar{\theta}^{2 n - k}
  \rangle$ and introduce the following expression
  \[ E_{\epsilon, N} (\sigma_1, \bar{\theta}) = \left.
     \left\{\begin{array}{lll}
       \int_{\mathbb{R}^2 \times \mathbb{T}^2} \sum_{k = 0}^{2 n} \binom{2
       n}{k} \left( \mathfrak{a}_{\frac{1}{N} - \epsilon} (\sigma_1) (\hat{z})
       \right)^{\diamond k} f' (x)  \bar{\theta}^{2 n - k} (\hat{z}) \mathd
       \hat{z}, & \quad & \text{for $\epsilon < \frac{1}{N}$,}\\
       \int_{\mathbb{R}^2 \times \mathbb{T}^2} \sum_{k = 0}^{2 n} \binom{2
       n}{k} (\sigma_1 (\hat{z}))^{\diamond k} f' (x)  \bar{\theta}^{2 n - k}
       (\hat{z}) \mathd \hat{z}, &  & \text{for $\epsilon \geqslant
       \frac{1}{N}$.}
     \end{array}\right. \right. \]
  We have to prove that, for any $p \geq 1$,
  \begin{equation}
    \left( \int | E - E_{\epsilon, N} |^p \mathd \nu_{\epsilon}^e \right)^{1 /
    p} \lesssim (p - 1)^u N^{- \alpha} \label{eq:IP1}
  \end{equation}
  for some $u \in \mathbb{N}$ and $\alpha > 0$ and
  \begin{equation}
    E_{\epsilon, N} \lesssim - f' (x) (\log (N))^{\alpha'} \label{eq:IP2},
  \end{equation}
  for some $\alpha' > 0$, and that both inequalities are uniform in
  $\epsilon$. Inequality~{\eqref{eq:IP1}} is obvious when $\epsilon \geqslant
  \frac{1}{N}$, since $E - E_{\epsilon, N} = 0$ $\nu^e_{\epsilon}$-almost
  surely. Consider the case $\epsilon < \frac{1}{N}$ then
  \begin{eqnarray}
    | E - E_{\epsilon, N} | & \lesssim & \left| \sum_{k = 0}^{2 n}
    \int_{\mathbb{R}^2 \times \mathbb{T}^2} f' (x) \left( \sigma_k (z) -
    \left( \mathfrak{a}_{\frac{1}{N} - \epsilon} (\sigma_1) (\hat{z})
    \right)^{\diamond k} \right) \cdot \bar{\theta}^{2 n - k} (\hat{z}) \mathd
    \hat{z} \right| \nonumber\\
    & \lesssim & \sum_{k = 1}^{2 n} \left\| | f' (x) |^{1 / 2 n} \left(
    \sigma_k (z) - \left( \mathfrak{a}_{\frac{1}{N} - \epsilon} (\sigma_1)
    (\hat{z}) \right)^{\diamond k} \right) \right\|_{\mathcal{C}^{- \delta}}
    \cdot \| | f' (x) |^{1 - 1 / 2 n} \bar{\theta}
    \|_{B^{\delta}_{\mathfrak{p, \mathfrak{p}}}}^{2 n - k} \nonumber
  \end{eqnarray}
  where $\mathfrak{p} = \frac{2 n - 1}{2 n}$. On the other hand, by
  Lemma~\ref{lemma_power1}, we have
  \[ \| | f' (x) |^{1 - 1 / 2 n} \bar{\theta} \|_{B^{\delta}_{\mathfrak{p,
     \mathfrak{p}}}}^{2 n - k} \lesssim \left( \| \bar{\theta} \|_{L^{2
     n}_{f^{1 / 2 n}}} + \| \bar{\theta} \|_{H^1} \right)^{2 n - k} \lesssim
     \left( \sum_{k = 1}^{2 n} \| \sigma_k \|_{\mathcal{C}^{- \delta}_{\ell}}
     \right)^{\beta_1 (2 n - 1)} \]
  $\nu^e_{\epsilon}$-almost surely and uniformly in $\epsilon$. Using Holder's
  inequality we obtain, for any $p \geq 1$,
  \begin{eqnarray}
    \left( \int | E - E_{\epsilon, N} |^p \mathd \nu_{\epsilon}^e
    \right)^{\frac{1}{p}} & \lesssim & \sum_{k = 1}^{2 n} \left[ \int \left\|
    | f' (x) |^{1 / 2 n} \left[ \sigma_k (z) - \left(
    \mathfrak{a}_{\frac{1}{N} - \epsilon} (\sigma_1) (\hat{z})
    \right)^{\diamond k} \right] \right\|_{\mathcal{C}^{- \delta}}^{2 p}
    \mathd \nu^e_{\epsilon} \right]^{\frac{1}{2 p}} \times \nonumber\\
    &  & \times \left( \int \left( \sum_{k = 1}^{2 n} \| \sigma_k
    \|_{\mathcal{C}^{- \delta}_{\ell}} \right)^{2 \beta_1 (2 n - 1) p} \mathd
    \nu_{\epsilon}^e \right)^{\frac{1}{2 p}} \nonumber
  \end{eqnarray}
  Since the law of $\sigma_k$ is a multilinear functional of a Gaussian random
  field, by hypercontractivity we have that there exists $m_1, m_2 \in
  \mathbb{N}$ and $\alpha > 0$ such that
  \begin{eqnarray}
    \int \left\| | f' (x) |^{1 / 2 n} \left( \sigma_k (z) - \left(
    \mathfrak{a}_{\frac{1}{N} - \epsilon} (\sigma_1) (\hat{z})
    \right)^{\diamond k} \right) \right\|_{\mathcal{C}^{- \delta}}^{2 p}
    \mathd \nu^e_{\epsilon} & \lesssim & (p - 1)^{p u_1} N^{- \alpha}
    \nonumber\\
    \int \left( \sum_{k = 1}^{2 n} \| \sigma_k \|_{\mathcal{C}^{-
    \delta}_{\ell}} \right)^{2 \beta_1 (2 N - 1) p} \mathd \nu_{\epsilon}^e &
    \lesssim & (p - 1)^{p u_2} \nonumber
  \end{eqnarray}
  uniformly in $\epsilon$. This prove inequality~{\eqref{eq:IP1}}. We note
  that
  \[ \sum_{k = 0}^{2 n} \binom{2 n}{k} \left( \mathfrak{a}_{\frac{1}{N} -
     \epsilon} (\sigma_1) (\hat{z}) \right)^{\diamond k} f' (x) 
     \bar{\theta}^{2 n - k} (\hat{z}) = f' (x) H_{2 n} \left(
     \mathfrak{a}_{\frac{1}{N} - \epsilon} (\sigma_1) (\hat{z}) + \bar{\theta}
     (\hat{z}) ; c_{N, \epsilon} \right) \]
  where $H_{2 n} (x ; \lambda)$ is the Hermite polynomial of degree $2 n$ of a
  Gaussian random variable with variance $\lambda$ and $c_{N, \epsilon}
  =\mathbb{E} \left[ \left( \mathfrak{a}_{\frac{1}{N}} \asterisk (\mathcal{I}
  \xi) \right)^2 \right] \sim \log (N)$, as $N \rightarrow + \infty$ for any
  $\epsilon > 0$. It is simple to see that, for any $x \in \mathbb{R}^2$ and
  $\hat{z} = (x, z) \in \mathbb{R}^2 \times \mathbb{T}^2$,
  \begin{eqnarray}
    f' (x) H_{2 n} \left( \mathfrak{a}_{\frac{1}{N} - \epsilon} (\sigma_1)
    (\hat{z}) + \bar{\theta} (\hat{z}) ; c_{N, \epsilon} \right) & \lesssim &
    a_n f' (x) \left( \mathfrak{a}_{\frac{1}{N} - \epsilon} (\sigma_1)
    (\hat{z}) + \bar{\theta} (\hat{z}) \right)^{2 n} \nonumber\\
    &  & + f' (x) \sum_{k = 1}^m a_k \left( \mathfrak{a}_{\frac{1}{N} -
    \epsilon} (\sigma_1) (\hat{z}) + \bar{\theta} (\hat{z}) \right)^{2 N - 2
    k} c_{N, \epsilon}^k \nonumber\\
    & \lesssim & \tilde{a}_n f' (x) \left( \mathfrak{a}_{\frac{1}{N} -
    \epsilon} (\sigma_1) (\hat{z}) + \bar{\theta} (\hat{z}) \right)^{2 N} - f'
    (x) c_{N, \epsilon}^n \nonumber\\
    & \lesssim & - f' (x) c_{N, \epsilon}^n \lesssim - f' (x) (\log (N))^n,
    \nonumber
  \end{eqnarray}
  where we used that $f' < 0$ and $a_N, \tilde{a}_N' > 0$. This proves
  inequality~{\eqref{eq:IP2}}.
  
  Now we use Lemma~V.5 of~{\cite{Simon1974}} and the fact that
  inequality~{\eqref{eq:IP1}} and~{\eqref{eq:IP2}} are uniform in $\epsilon$,
  to prove that there exist two constants $b, \alpha'' > 0$ depending only on
  $f'$ and $m$ but not on $\epsilon > 0$ such that
  \begin{equation}
    \nu_{\epsilon}^e (E (\sigma, \bar{\theta}) > b \log (K)) \leqslant e^{-
    K^{\alpha''}} . \label{eq:measure}
  \end{equation}
  Inequality~{\eqref{eq:measure}} implies that $\exp (E) \in L^p
  (\nu^e_{\epsilon})$ for any $p \geq 1$, and we have also a uniform bound on
  the $L^p$ norm of $\exp (E)$ with respect to $\epsilon > 0$. This concludes
  the proof of the lemma.
\end{proof*}

\begin{proof*}{Proof of Theorem~\ref{theorem_polynomialreduction}}
  First we introduce the equation
  \begin{equation}
    (- \Delta_x - \Delta_z + m^2) (\phi) +
    \mathfrak{a}_{\varepsilon}^{\asterisk 2} \asterisk (H_{2 n} (\phi ;
    c_{\epsilon})) = \mathfrak{a}_{\epsilon} \asterisk \xi_{}
    \label{eq:powerproof1},
  \end{equation}
  where $c_{\epsilon} =\mathbb{E} [(\mathfrak{a}_{\epsilon} \asterisk
  \xi_{})^2]$. Equation~{\eqref{eq:powerproof1}} is of the
  form~{\eqref{eq:main1}}, and the potential $V_{\epsilon} (y) = H_{2 n} (y ;
  c_{\epsilon})$, where $y \in \mathbb{R}$, satisfies Hypothesis~QC since it
  is an even polynomial with positive leading coefficient. This means that,
  for any $\epsilon > 0$, there exists a weak solution $\nu_{\epsilon}$ to
  equation~{\eqref{eq:powerproof1}}, such that
  \[ \int_{\mathcal{W}} F (\phi (0, \cdot)) \Upsilon_f^{\epsilon} (\phi)
     \mathd \nu_{\epsilon} = \int_{B^{- \delta}_{\mathfrak{p}, \mathfrak{p},
     \ell} (\mathbb{T}^2)} F (\omega) \mathd \kappa_{\epsilon} (\omega) \]
  where $\Upsilon_f^{\epsilon} (\phi) = \int_{\mathbb{R}^2 \times
  \mathbb{T}^2} f' (x) H_{2 n} (\phi (\hat{z}) ; c_{\epsilon}) \mathd
  \hat{z}$. On the other hand if by $\nu^e_{\epsilon}$ we denote the law of
  $(\mathcal{\mathfrak{a}}_{\epsilon} \asterisk \mathcal{I} \xi, \ldots,
  (\mathcal{\mathfrak{a}}_{\epsilon} \asterisk \mathcal{I} \xi)^{\diamond 2
  n}, \phi - \mathcal{\mathfrak{a}}_{\epsilon} \asterisk \mathcal{I} \xi)$
  (where $\phi$ has law $\nu$ and $\mathcal{\mathfrak{a}}_{\epsilon} \asterisk
  \mathcal{I} \xi$ is given by equation~{\eqref{eq:powerproof1}}), we have
  that $\nu^e_{\epsilon}$ is a solution to equation~{\eqref{eq:power4}} and
  \[ \int_{\mathcal{W}} F (\phi (0, \cdot)) \Upsilon_f^{\epsilon} (\phi)
     \mathd \nu_{\epsilon} = \int_{\mathcal{W}^e} F (\phi (0, \cdot))
     \tilde{\Upsilon}_f (\sigma, \bar{\theta}) \mathd \nu^e_{\epsilon}
     (\sigma, \bar{\theta}) . \]
  By Lemma~\ref{lemma_xi} and Lemma~\ref{lemma_trace} we have that the measure
  $F (\phi (0, \cdot)) \mathd \nu_{\epsilon}^e$ weakly converges to $F (\phi
  (0, \cdot)) \mathd \nu^e$, as $\varepsilon \rightarrow 0$. Let $G_r :
  \mathcal{W}^e \rightarrow \mathbbm{R}$, $r \in \mathbb{N}$, be a sequence of
  continuous functions such that $0 \leqslant G_r \leqslant 1$, $G_r = 1$ when
  $\tilde{\Upsilon}_f (\sigma, \bar{\theta}) \leqslant r$ and $G_r = 0$ when
  $\tilde{\Upsilon}_f (\sigma, \bar{\theta}) \geqslant 2 r$ (the existence of
  this kind of functions follows from the fact that $\tilde{\Upsilon}_f
  (\sigma, \bar{\theta})$ is continuous on $\mathcal{W}^e$). Since $F (\phi
  (0, \cdot)) \mathd \nu_{\epsilon}^e$ weakly converges to $F (\phi (0,
  \cdot)) \mathd \nu^e$ then $\int_{\mathcal{W}^e} G_r F \tilde{\Upsilon}_f
  \mathd \nu^e_{\epsilon} \rightarrow \int_{\mathcal{W}^e} G_r F
  \tilde{\Upsilon}_f \mathd \nu^e$, as $\varepsilon \rightarrow 0$. On the
  other hand, by Lemma~\ref{lemma_uniform}, for all $\varepsilon > 0$ and any
  $r > 0$,
  \[ \left| \int_{\mathcal{W}^e} (G_r - 1) F \tilde{\Upsilon}_f \mathd
     \nu^e_{\epsilon} \right| \leqslant \int_{\tilde{\Upsilon}_f \geqslant r}
     | F | \tilde{\Upsilon}_f \mathd \nu^e_{\epsilon} \leqslant \frac{C_p \| F
     \|_{\infty}}{r^{p - 1}} \]
  for any $p \geq 1$ (and a similar inequality is true for $\epsilon = 0$).
  Taking $r \rightarrow + \infty$ the thesis follows.
\end{proof*}

\appendix\section{Besov spaces}\label{appendix_besov}

In this appendix we recall some results on weighted Besov spaces used in this
paper. We consider only the case of Besov spaces defined on $\mathbb{R}^n$ but
all what follows holds also for Besov spaces on $\mathbb{T}^n$ or on
$\mathbb{R}^{n_1} \times \mathbb{T}^{n_2}$ (like the ones used in Section
\ref{section_power}).

\

First we recall the definition of Littewood-Paley block: let $\chi, \varphi$
be smooth non-negative functions from $\mathbb{R}^n$ into $\mathbb{R}$ such
that
\begin{itemize}
  \item $\tmop{supp} (\chi) \subset B_1 (0)$ and $\tmop{supp} (\varphi)
  \subset B_2 (0) \setminus B_1 (0)$
  
  \item $\chi, \varphi \leq 1$ and $\chi (y) + \sum_{j \geq 0} \varphi (2^{-
  j} y) = 1$ for any $y \in \mathbb{R}^n$
  
  \item $\tmop{supp} (\varphi (2^{- j} \cdot)) \cap \tmop{supp} (\varphi (2^{-
  i} \cdot)) = \emptyset$ if $| i - j | > 1$,
\end{itemize}
where by $B_r (x)$ we denote the ball of center $x \in \mathbb{R}^n$ and
radius $r > 0$.

We introduce the following notations: $\varphi_{- 1} = \chi$, $\varphi_j
(\cdot) = \varphi (2^{- j} \cdot)$, $D_j = \hat{\varphi}_j$ and for any $f \in
\mathcal{S}' (\mathbb{R}^n)$ we put $\Delta_j (f) = D_j \asterisk f$.
Furthermore we write, for any $\ell > 0$, \ $\rho_{\ell} (y) = (1 + | y
|^2)^{- \ell / 2}$, $L^p_{\ell} (\mathbb{R}^d)$ is the $L^p$ space with
respect to the norm
\[ \| f \|_{p, \ell} = \left( \int_{\mathbb{R}^n} (f (y) \rho_{\ell} (y))^p
   \mathd y \right)^{1 / p}, \]
where $p \in [1, + \infty]$.

\begin{definition}
  Consider $s \in \mathbb{R}, p, q \in [1, + \infty]$ and $\ell \geq 0$. If $f
  \in \mathcal{S} (\mathbb{R}^n)$ we define the norm
  \begin{equation}
    \| f \|_{B^s_{p, q, \ell}} = \left( \sum_{j \geq - 1} \| \Delta_j (f)
    \|_{p, \ell}^q \right)^{1 / q} . \label{eq:norm11}
  \end{equation}
  The space $B^s_{p, q, \ell}$ is the subset of $\mathcal{S}' (\mathbb{R}^n)$
  such that the norm {\eqref{eq:norm11}} is finite. 
\end{definition}

\begin{remark}
  It is important to note that $B^s_{2, 2, \ell}$ is equal to the weighted
  Sobolev space $W^{s, 2}_{\ell} = H^s_{\ell}$ of $L^2$ distributions, and
  that for $s > 0$, $s \not{\in} \mathbb{N}$, $B^s_{\infty, \infty}$ is equal
  to $\mathcal{C}^s$, the space of Holder functions of regularity $s$. 
\end{remark}

We denote by $\mathfrak{F}_r$, with $r \in \mathbb{N}$, the space of
functions$f \in C^r (\mathbb{R}^n)$ with support in $B_1 (0)$ and norm $\| f
\|_{\mathfrak{} C^r (\mathbb{R}^n)} \leq 1$. If $f \in \mathfrak{F}_r$ we
write $f_{y, \lambda} (\cdot) = \lambda^{- n} f \left( \frac{\cdot -
y}{\lambda} \right)$, $y \in \mathbb{R}^n$ and $\lambda > 0$. \

\begin{proposition}
  \label{proposition_equivalentnorm}For any $s \in \mathbb{R}, p, q \in [1, +
  \infty]$ and $\ell \geq 0$ an equivalent norm in the space $B^s_{p, q,
  \ell}$ is given by the following expression
  \begin{equation}
    \| f \|_{B^s_{p, q, \ell}} = \left( \int_0^1 \frac{\| \sup_{g \in
    \mathfrak{F}_r} | \langle f, g_{\cdot, \lambda} \rangle | \|_{p,
    \ell}^q}{\lambda^{s q}} \right)^{1 / q} . \label{eq:equivalentnorm}
  \end{equation}
  where $r \in \mathbb{N}$ is the first integer bugger then $- s$.
\end{proposition}

\begin{proof}
  Theorem 6.15, in {\cite{Triebel2006}}, proves the equivalece between norm
  {\eqref{eq:norm11}} and the norm of $B^s_{p, p, \ell}$ built using wavelets,
  while Proposition 2.4 in {\cite{HaiLabb2017}} proves the equivalence between
  the norm of $B^s_{p, p, \ell}$ built using wavelets and the norm
  {\eqref{eq:equivalentnorm}}. Combining the two results we obtain the thesis.
\end{proof}

\begin{proposition}
  \label{proposition_inclusion}Consider $p_1, p_2, q_1, q_2 \in [1, \infty]$,
  $\ell_{1,} \ell_2 \geq 0$ and $s_1, s_2 \in \mathbb{R}$ such that $s_1 -
  \frac{n}{p_1} > s_2 - \frac{n}{p_2}$ and $\ell_1 > \ell_2$ then
  $B^{s_2}_{p_2, q_2, \ell_2} \subset B^{s_1}_{p_1, q_1, \ell_1}$ and the
  immersion is compact.
  
  Furthermore for $1 \leq p \leq 2$
  \[ B^s_{p, p, \ell} \subset W^{s, p}_{\ell} \subset B^s_{p, 2, \ell} ; \]
  for $2 \leq p < \infty$
  \[ B^s_{p, 2, \ell} \subset W^{s, p}_{\ell} \subset B^s_{p, p, \ell} \]
  and for $p = \infty$ $W^{s, \infty}_{\ell} \subset B^s_{\infty, \infty}$.
  Each of above immersions is continuous. 
\end{proposition}

\begin{proof}
  The proof of the first part of the proposition can be found in Theorem 6.7
  in {\cite{Triebel2006}}.
  
  The second part of the proposition is proved in Theorem 6.4.4 and Theorem
  6.2.4 of {\cite{Joran1976}} for unweighted spaces. The result for weighted
  spaces follows from the fact that $f \in B^s_{p, q, \ell}$ and $g \in W^{s,
  p}_{\ell}$ if and only if $f \cdot \rho_{\ell} \in B^s_{p, q}$ and $g \cdot
  \rho_{\ell} \in W^{s, p}_{\ell}$ (see {\cite{Triebel2006}} Theorem 6.5 point
  ii) ). 
\end{proof}

\begin{proposition}
  \label{propisition_product}Consider $p_1, p_2, p_3, q_1, q_2, q_3 \in [1,
  \infty]$ such that $\frac{1}{p_1} + \frac{1}{p_2} = \frac{1}{p_3}$ and $q_1
  = q_3$ and $q_2 = \infty$, consider $\ell_{1,} \ell_2, \ell_3 \geq 0$ with
  $\ell_1 + \ell_2 = \ell_3$ and consider $s_1 < 0$ $s_2 \geq 0$ and $s_3 =
  s_1 + s_2$. Consider the bilinear functional \ $\Pi (f, g) = f \cdot g$
  defined on $\mathcal{S} (\mathbb{R}^n)$ taking values in $\mathcal{S}
  (\mathbb{R}^n)$. Then there exists a unique continuous extension of $\Pi$ as
  the map
  \[ \Pi : B^{s_1}_{p_1, q_1, \ell_1} \times B^{s_2}_{p_2, q_2, \ell_2}
     \rightarrow B^{s_3}_{p_3, q_3, \ell_3} \]
  and we have, for any $f, g$ for which the norms are defined:
  \[ \| \Pi (f, g) \|_{B^{s_3}_{p_3, q_3, \ell_3}} \lesssim \| f
     \|_{B^{s_1}_{p_1, q_1, \ell_1}} \| g \|_{B^{s_2}_{p_2, q_2, \ell_2}} . \]
\end{proposition}

\begin{proof}
  The proof can be found in {\cite{MW17}} in Section 3.3 for Besov spaces with
  exponential weights. The proof for polynomial weights is similar.
\end{proof}

\begin{proposition}
  \label{proposition_interpolation}Consider $p_1, p_2, q_1, q_2 \in [1,
  \infty]$, $\ell_{1,} \ell_2 \geq 0$ and $s_1, s_2 \in \mathbb{R}$. For any
  $\theta \in [0, 1]$ we write $\frac{\theta}{p_1} + \frac{1 - \theta}{p_2} =
  \frac{1}{p_{\theta}}$, $\frac{\theta}{q_1} + \frac{1 - \theta}{q_2} =
  \frac{1}{q_{\theta}}$, $\theta \ell_1 + (1 - \theta) \ell_2 = \ell_{\theta}$
  and $\theta s_1 + (1 - \theta) s_2 = s_{\theta}$. If $f \in B^{s_1}_{p_1,
  q_1, \ell_1} \cap B^{s_2}_{p_2, q_2, \ell_2}$ then $f \in
  B^{s_{\theta}}_{p_{\theta}, q_{\theta}, \ell_{\theta}}$, and furthermore
  \[ \| f \|_{B^{s_{\theta}}_{p_{\theta}, q_{\theta}, \ell_{\theta}}} \leq \|
     f \|_{B^{s_1}_{p_1, q_1, \ell_1}}^{\theta} \| f \|_{B^{s_2}_{p_2, q_2,
     \ell_2}}^{1 - \theta} . \]
  Furthermore if $f \in W^{s_1, p_1}_{\ell_1} \cap W^{s_2, p_2}_{\ell_2}$ then
  \ $f \in W^{s_{\theta}, p_{\theta}}_{\ell_{\theta}}$, and furthermore
  \[ \| f \|_{W^{s_{\theta}, p_{\theta}}_{\ell_{\theta}}} \leq \| f
     \|_{W^{s_1, p_1}_{\ell_1}}^{\theta} \| f \|_{W^{s_2, p_2}_{\ell_2}}^{1 -
     \theta} . \]
\end{proposition}

\begin{proof}
  The proof is based on the fact that the complex interpolation
  $(B^{s_1}_{p_1, q_1, \ell_1}, B^{s_2}_{p_2, q_2, \ell_2})_{\theta}$ of the
  two spaces $B^{s_1}_{p_1, q_1, \ell_1}$, $B^{s_2}_{p_2, q_2, \ell_2}$ is
  given by $B^{s_{\theta}}_{p_{\theta}, q_{\theta}, \ell_{\theta}}$. This
  interpolation is proved in~{\cite{Joran1976}} Theorem 6.4.5 point (6) for
  unweighted space. For weighted space it follows from the fact that $f \in
  B^{s_i}_{p_i, q_i, \ell_i}$ if and only if $f \cdot \rho_{\ell_i} \in
  B^{s_i}_{p_i, q_i}$.
  
  A similar reasoning, using the interpolation proved in~{\cite{Joran1976}}
  Theorem 6.4.5 point (7), holds also for the second inequality. 
\end{proof}

\begin{thebibliography}{10}
  \bibitem[1]{AGoVe}S.~Albeverio, B.~K.~Driver, M.~Gordina, and 
  A.~M.~Vershik. {\newblock}Equivalence of the Brownian and energy
  representations. {\newblock}\tmtextit{Zap. Nauchn. Sem. S.-Peterburg. Otdel.
  Mat. Inst. Steklov. (POMI)}, 441(Veroyatnostc i Statistika. 22):17--44,
  2015. {\newblock}Reprinted in J. Math. Sci. (N.Y.) \tmtextbf{2}19 (2016),
  no. 5, 612--630.{\newblock}
  
  \bibitem[2]{AGaHK}S.~Albeverio, G.~Gallavotti, and  R.~H{\o}egh-Krohn.
  {\newblock}Some results for the exponential interaction in two or more
  dimensions. {\newblock}\tmtextit{Comm. Math. Phys.}, 70(2):187--192,
  1979.{\newblock}
  
  \bibitem[3]{AHiPRoS2}S.~Albeverio, T.~Hida, J.~Potthoff, M.~R{\"o}ckner, and
  L.~Streit. {\newblock}Dirichlet forms in terms of white noise analysis. I.
  Construction and QFT examples. {\newblock}\tmtextit{Rev. Math. Phys.},
  1(2-3):291--312, 1989.{\newblock}
  
  \bibitem[4]{AHiPRoS1}S.~Albeverio, T.~Hida, J.~Potthoff, M.~R{\"o}ckner, and
  L.~Streit. {\newblock}Dirichlet forms in terms of white noise analysis. II.
  Closability and diffusion processes. {\newblock}\tmtextit{Rev. Math. Phys.},
  1(2-3):313--323, 1989.{\newblock}
  
  \bibitem[5]{Yoshida}S.~Albeverio, T.~Kagawa, Y.~Yahagi, and  M.~Yoshida.
  {\newblock}Non-local markovian symmetric forms on infinite dimensional
  spaces i, ii. {\newblock}2019.{\newblock}
  
  \bibitem[6]{AKaMiRo}S.~Albeverio, H.~Kawabi, S.~Mihalache, and 
  M.~R{\"o}kner. {\newblock}Dirichlet form approach to stochastic quantization
  under exponential interaction in finite volume. {\newblock}In preparation,
  2016.{\newblock}
  
  \bibitem[7]{albeverio_invariant_2017}S.~Albeverio  and  S.~Kusuoka.
  {\newblock}The invariant measure and the flow associated to the
  $\Phi^4_3$-quantum field model. {\newblock}\tmtextit{ArXiv:1711.07108}, nov
  2017. {\newblock}To appear in Ann. Scuola Normale di Pisa.{\newblock}
  
  \bibitem[8]{ARu}S.~Albeverio  and  B.~R{\"u}diger.
  {\newblock}Infinite-dimensional stochastic differential equations obtained
  by subordination and related Dirichlet forms. {\newblock}\tmtextit{J. Funct.
  Anal.}, 204(1):122--156, 2003.{\newblock}
  
  \bibitem[9]{Albeverio2002}S.~Albeverio  and  M.~W.~Yoshida. {\newblock}$H -
  C^1$ maps and elliptic spdes with polynomial and exponential perturbations
  of Nelson's euclidean free field. {\newblock}\tmtextit{Journal of Functional
  Analysis}, 196(2):265--322, 2002.{\newblock}
  
  \bibitem[10]{AHKMT}S.~A.~Albeverio, R.~J.~H{\o}egh-Krohn, J.~A.~Marion,
  D.~H.~Testard, and  B.~S.~Torr{\'e}sani. {\newblock}\tmtextit{Noncommutative
  distributions},  volume  175  of \tmtextit{Monographs and Textbooks in Pure
  and Applied Mathematics}. {\newblock}Marcel Dekker, Inc., New York, 1993.
  {\newblock}Unitary representation of gauge groups and algebras.{\newblock}
  
  \bibitem[11]{Albeverio2018elliptic}S.~Albeverio, F.~C.~De Vecchi, and 
  M.~Gubinelli. {\newblock}Elliptic stochastic quantization.
  {\newblock}\tmtextit{ArXiv preprint arXiv:1812.04422}, 2018.{\newblock}
  
  \bibitem[12]{AHaRu}S.~Albeverio, Z.~Haba, and  F.~Russo. {\newblock}A
  two-space dimensional semilinear heat equation perturbed by (Gaussian) white
  noise. {\newblock}\tmtextit{Probab. Theory Related Fields}, 121(3):319--366,
  2001.{\newblock}
  
  \bibitem[13]{AHK3}S.~Albeverio  and  R.~H{\o}egh-Krohn.
  {\newblock}Uniqueness of the physical vacuum and the Wightman functions in
  the infinite volume limit for some non polynomial interactions.
  {\newblock}\tmtextit{Comm. Math. Phys.}, 30:171--200, 1973.{\newblock}
  
  \bibitem[14]{Krohn1974}S.~Albeverio  and  R.~H{\o}egh-Krohn. {\newblock}The
  Wightman axioms and the mass gap for strong interactions of exponential type
  in two-dimensional space-time. {\newblock}\tmtextit{J. Functional Analysis},
  16:39--82, 1974.{\newblock}
  
  \bibitem[15]{AlHKa}S.~Albeverio  and  R.~H{\o}egh-Krohn.
  {\newblock}Uniqueness and the global Markov property for Euclidean fields.
  The case of trigonometric interactions. {\newblock}\tmtextit{Comm. Math.
  Phys.}, 68(2):95--128, 1979.{\newblock}
  
  \bibitem[16]{AHK2}S.~Albeverio  and  R.~H{\o}egh-Krohn.
  {\newblock}Martingale convergence and the exponential interaction in
  $\tmmathbf{R}^n$. {\newblock}In \tmtextit{Quantum fields---algebras,
  processes (Proc. Sympos., Univ. Bielefeld, Bielefeld, 1978)},  pages 
  331--353. Springer, Vienna, 1980.{\newblock}
  
  \bibitem[17]{AFeHKL}S.~Albeverio, R.~H{\o}egh-Krohn, J.~E.~Fenstad, and 
  T.~Lindstr{\o}m. {\newblock}\tmtextit{Nonstandard methods in stochastic
  analysis and mathematical physics},  volume  122  of \tmtextit{Pure and
  Applied Mathematics}. {\newblock}Academic Press, Inc., Orlando, FL,
  1986.{\newblock}
  
  \bibitem[18]{AlKrPaSc1986}S.~Albeverio, R.~H{\o}egh-Krohn, S.~Paycha, and 
  S.~Scarlatti. {\newblock}Path space measure for the Liouville quantum field
  theory and the construction of relativistic strings.
  {\newblock}\tmtextit{Phys. Lett. B}, 174(1):81--86, 1986.{\newblock}
  
  \bibitem[19]{AlKrPaSc1992}S.~Albeverio, R.~H{\o}egh-Krohn, S.~Paycha, and 
  S.~Scarlatti. {\newblock}A global and stochastic analysis approach to
  bosonic strings and associated quantum fields. {\newblock}\tmtextit{Acta
  Appl. Math.}, 26(2):103--195, 1992.{\newblock}
  
  \bibitem[20]{AJPS}S.~Albeverio, J.~Jost, S.~Paycha, and  S.~Scarlatti.
  {\newblock}\tmtextit{A mathematical introduction to string theory},  volume 
  225  of \tmtextit{London Mathematical Society Lecture Note Series}.
  {\newblock}Cambridge University Press, Cambridge, 1997.
  {\newblock}Variational problems, geometric and probabilistic
  methods.{\newblock}
  
  \bibitem[21]{AKaRo}S.~Albeverio, H.~Kawabi, and  M.~R{\"o}ckner.
  {\newblock}Strong uniqueness for both Dirichlet operators and stochastic
  dynamics to Gibbs measures on a path space with exponential interactions.
  {\newblock}\tmtextit{J. Funct. Anal.}, 262(2):602--638, 2012.{\newblock}
  
  \bibitem[22]{AlLi}S.~Albeverio  and  S.~Liang. {\newblock}A new lattice
  approximation for the H{\o}egh-Krohn quantum field model.
  {\newblock}\tmtextit{Rep. Math. Phys.}, 54(2):149--157, 2004.{\newblock}
  
  \bibitem[23]{Yoshida2008}S.~Albeverio  and  M.~W.~Yoshida.
  {\newblock}Reflection positive random fields and Dirichlet spaces.
  {\newblock}In \tmtextit{Proceedings of RIMS Workshop on Stochastic Analysis
  and Applications}, RIMS K{\^o}ky{\^u}roku Bessatsu, B6,  pages  15--29. Res.
  Inst. Math. Sci. (RIMS), Kyoto, 2008.{\newblock}
  
  \bibitem[24]{Yoshida2010}S.~Albeverio  and  M.~W.~Yoshida. {\newblock}Some
  abstract considerations on the homogenization problem of infinite
  dimensional diffusions. {\newblock}In \tmtextit{Applications of
  renormalization group methods in mathematical sciences}, RIMS
  K{\^o}ky{\^u}roku Bessatsu, B21,  pages  183--192. Res. Inst. Math. Sci.
  (RIMS), Kyoto, 2010.{\newblock}
  
  \bibitem[25]{AnKa}S.~Andres  and  N.~Kajino. {\newblock}Continuity and
  estimates of the Liouville heat kernel with applications to spectral
  dimensions. {\newblock}\tmtextit{Probab. Theory Related Fields},
  166(3-4):713--752, 2016.{\newblock}
  
  \bibitem[26]{Bahns2018}D.~Bahns  and  K.~Rejzner. {\newblock}The quantum
  sine-Gordon model in perturbative AQFT. {\newblock}\tmtextit{Comm. Math.
  Phys.}, 357(1):421--446, 2018.{\newblock}
  
  \bibitem[27]{Bahouri2011}H.~Bahouri, J.-Y.~Chemin, and  R.~Danchin.
  {\newblock}\tmtextit{Fourier analysis and nonlinear partial differential
  equations},  volume  343  of \tmtextit{Grundlehren der Mathematischen
  Wissenschaften [Fundamental Principles of Mathematical Sciences]}.
  {\newblock}Springer, Heidelberg, 2011.{\newblock}
  
  \bibitem[28]{GaNi}G.~Benfatto, G.~Gallavotti, and  F.~Nicol{\`o}.
  {\newblock}On the massive sine-Gordon equation in the first few regions of
  collapse. {\newblock}\tmtextit{Comm. Math. Phys.}, 83(3):387--410,
  1982.{\newblock}
  
  \bibitem[29]{BeShXi2014}N.~Berestycki, S.~Sheffield, and  X.~Sun.
  {\newblock}Equivalence of liouville measure and gaussian free field.
  {\newblock}\tmtextit{ArXiv preprint arXiv:1410.5407}, 2014.{\newblock}
  
  \bibitem[30]{Joran1976}J.~Bergh  and  J.~L{\"o}fstr{\"o}m.
  {\newblock}\tmtextit{Interpolation spaces. An introduction}.
  {\newblock}Springer-Verlag, Berlin-New York, 1976. {\newblock}Grundlehren
  der Mathematischen Wissenschaften, No. 223.{\newblock}
  
  \bibitem[31]{Boyko2009}N.~Boyko  and  V.~Kadets. {\newblock}Uniform
  $G$-convexity for vector-valued $L_p$ spaces. {\newblock}\tmtextit{Serdica
  Math. J.}, 35(1):1--14, 2009.{\newblock}
  
  \bibitem[32]{CaCh}R.~Catellier  and  K.~Chouk. {\newblock}Paracontrolled
  distributions and the 3-dimensional stochastic quantization equation.
  {\newblock}\tmtextit{Ann. Probab.}, 46(5):2621--2679, 2018.{\newblock}
  
  \bibitem[33]{ChaHaiShe}A.~Chandra, M.~Hairer, and  H.~Shen. {\newblock}The
  dynamical sine-gordon model in the full subcritical regime.
  {\newblock}\tmtextit{ArXiv preprint arXiv:1808.02594}, 2018.{\newblock}
  
  \bibitem[34]{da_prato_strong_2003}G.~Da Prato  and  A.~Debussche.
  {\newblock}Strong solutions to the stochastic quantization equations.
  {\newblock}\tmtextit{The Annals of Probability}, 31(4):1900--1916,
  2003.{\newblock}
  
  \bibitem[35]{DaPraDe}G.~Da Prato  and  A.~Debussche. {\newblock}Strong
  solutions to the stochastic quantization equations.
  {\newblock}\tmtextit{Ann. Probab.}, 31(4):1900--1916, 2003.{\newblock}
  
  \bibitem[36]{DuSh2019}J.~Dub{\'e}dat  and  H.~Shen. {\newblock}Stochastic
  Ricci Flow on Compact Surfaces. {\newblock}\tmtextit{ArXiv preprint
  arXiv:1904.10909}, 2019.{\newblock}
  
  \bibitem[37]{VaRhShe}B.~Duplantier, R.~Rhodes, S.~Sheffield, and  V.~Vargas.
  {\newblock}Log-correlated Gaussian fields: an overview. {\newblock}In
  \tmtextit{Geometry, analysis and probability},  volume  310  of
  \tmtextit{Progr. Math.},  pages  191--216. Birkh{\"a}user/Springer, Cham,
  2017.{\newblock}
  
  \bibitem[38]{Evans1998}L.~C.~Evans. {\newblock}\tmtextit{Partial
  differential equations},  volume~19  of \tmtextit{Graduate Studies in
  Mathematics}. {\newblock}American Mathematical Society, Providence, RI,
  1998.{\newblock}
  
  \bibitem[39]{FroPa}J.~Fr{\"o}hlich  and  Y.~M.~Park. {\newblock}Remarks on
  exponential interactions and the quantum sine-Gordon equation in two
  space-time dimensions. {\newblock}\tmtextit{Helv. Phys. Acta},
  50(3):315--329, 1977.{\newblock}
  
  \bibitem[40]{FrSei}J.~Fr{\"o}hlich  and  E.~Seiler. {\newblock}The massive
  Thirring-Schwinger model (\tmtextrm{qed}2): convergence of perturbation
  theory and particle structure. {\newblock}\tmtextit{Helv. Phys. Acta},
  49(6):889--924, 1976.{\newblock}
  
  \bibitem[41]{Garban2018}C.~Garban. {\newblock}Dynamical Liouville.
  {\newblock}\tmtextit{ArXiv preprint arXiv:1805.04507}, 2018.{\newblock}
  
  \bibitem[42]{Gie}R.~Gielerak. {\newblock}Verification of the global Markov
  property in some class of strongly coupled exponential interactions.
  {\newblock}\tmtextit{J. Math. Phys.}, 24(2):347--355, 1983.{\newblock}
  
  \bibitem[43]{Glimm1987}J.~Glimm  and  A.~Jaffe. {\newblock}\tmtextit{Quantum
  physics}. {\newblock}Springer-Verlag, New York, Second  edition, 1987.
  {\newblock}A functional integral point of view.{\newblock}
  
  \bibitem[44]{Gottschalk2013}H.~Gottschalk  and  H.~Thaler. {\newblock}A
  triviality result in the AdS/CFT correspondence for Euclidean quantum fields
  with exponential interaction. {\newblock}\tmtextit{Comm. Math. Phys.},
  324(1):63--75, 2013.{\newblock}
  
  \bibitem[45]{gubinelli_pde_2018}M.~Gubinelli  and  M.~Hofmanov{\'a}.
  {\newblock}A PDE construction of the Euclidean $\Phi^4_3$ quantum field
  theory. {\newblock}\tmtextit{ArXiv preprint arXiv:1810.01700}, oct
  2018.{\newblock}
  
  \bibitem[46]{GH18}M.~Gubinelli  and  M.~Hofmanov{\'a}. {\newblock}Global
  Solutions to Elliptic and Parabolic $\Phi^4$ Models in Euclidean Space.
  {\newblock}\tmtextit{Communications in Mathematical Physics}, Mar
  2019.{\newblock}
  
  \bibitem[47]{GuImPe}M.~Gubinelli, P.~Imkeller, and  N.~Perkowski.
  {\newblock}Paracontrolled distributions and singular PDEs.
  {\newblock}\tmtextit{Forum Math. Pi}, 3:0, 2015.{\newblock}
  
  \bibitem[48]{Hairer2014}M.~Hairer. {\newblock}A theory of regularity
  structures. {\newblock}\tmtextit{Invent. Math.}, 198(2):269--504,
  2014.{\newblock}
  
  \bibitem[49]{HaiLabb2017}M.~Hairer  and  C.~Labb{\'e}. {\newblock}The
  reconstruction theorem in Besov spaces. {\newblock}\tmtextit{J. Funct.
  Anal.}, 273(8):2578--2618, 2017.{\newblock}
  
  \bibitem[50]{HaiShe}M.~Hairer  and  H.~Shen. {\newblock}The dynamical
  sine-Gordon model. {\newblock}\tmtextit{Comm. Math. Phys.}, 341(3):933--989,
  2016.{\newblock}
  
  \bibitem[51]{HiKPS}T.~Hida, H.-H.~Kuo, J.~Potthoff, and  L.~Streit.
  {\newblock}\tmtextit{White noise},  volume  253  of \tmtextit{Mathematics
  and its Applications}. {\newblock}Kluwer Academic Publishers Group,
  Dordrecht, 1993. {\newblock}An infinite-dimensional calculus.{\newblock}
  
  \bibitem[52]{Krohn1971}R.~H{\o}egh-Krohn. {\newblock}A general class of
  quantum fields without cut-offs in two space-time dimensions.
  {\newblock}\tmtextit{Comm. Math. Phys.}, 21:244--255, 1971.{\newblock}
  
  \bibitem[53]{Iwata1987}K.~Iwata. {\newblock}Reversible measures of a $P
  (\phi)_1$-time evolution. {\newblock}In \tmtextit{Probabilistic methods in
  mathematical physics (Katata/Kyoto, 1985)},  pages  195--209. Academic
  Press, Boston, MA, 1987.{\newblock}
  
  \bibitem[54]{Ka}J.-P.~Kahane. {\newblock}Sur le chaos multiplicatif.
  {\newblock}\tmtextit{Ann. Sci. Math. Qu{\'e}bec}, 9(2):105--150,
  1985.{\newblock}
  
  \bibitem[55]{Klein1984}A.~Klein, L.~J.~Landau, and  J.~F.~Perez.
  {\newblock}Supersymmetry and the Parisi-Sourlas dimensional reduction: a
  rigorous proof. {\newblock}\tmtextit{Comm. Math. Phys.}, 94(4):459--482,
  1984.{\newblock}
  
  \bibitem[56]{Ku}S.~Kusuoka. {\newblock}H{\o}egh-Krohn's model of quantum
  fields and the absolute continuity of measures. {\newblock}In
  \tmtextit{Ideas and methods in quantum and statistical physics (Oslo,
  1988)},  pages  405--424. Cambridge Univ. Press, Cambridge, 1992.{\newblock}
  
  \bibitem[57]{Liang}S.~Liang. {\newblock}Different lattice approximations for
  H{\o}egh-Krohn's quantum field model. {\newblock}In \tmtextit{Stochastic
  analysis and applications},  volume~2  of \tmtextit{Abel Symp.},  pages 
  493--499. Springer, Berlin, 2007.{\newblock}
  
  \bibitem[58]{Mi}S.~Mihalache. {\newblock}Stochastische quantisierung bei
  exponentieller wechselwirkung. {\newblock}Master's thesis, Universi{\"a}t
  Bonn, 2006.{\newblock}
  
  \bibitem[59]{MW17}J.-C.~Mourrat  and  H.~Weber. {\newblock}Global
  well-posedness of the dynamic $\Phi^4$ model in the plane.
  {\newblock}\tmtextit{The Annals of Probability}, 45(4):2398--2476, jul
  2017.{\newblock}
  
  \bibitem[60]{MoWe}J.-C.~Mourrat  and  H.~Weber. {\newblock}The dynamic
  $\Phi^4_3$ model comes down from infinity. {\newblock}\tmtextit{Comm. Math.
  Phys.}, 356(3):673--753, 2017.{\newblock}
  
  \bibitem[61]{parisi_random_1979}G.~Parisi  and  N.~Sourlas.
  {\newblock}Random Magnetic Fields, Supersymmetry, and Negative Dimensions.
  {\newblock}\tmtextit{Physical Review Letters}, 43(11):744--745, sep
  1979.{\newblock}
  
  \bibitem[62]{parisi_perturbation_1981}G.~Parisi  and  Y.~S.~Wu.
  {\newblock}Perturbation theory without gauge fixing.
  {\newblock}\tmtextit{Scientia Sinica. Zhongguo Kexue}, 24(4):483--496,
  1981.{\newblock}
  
  \bibitem[63]{RhVa}R.~Rhodes  and  V.~Vargas. {\newblock}Gaussian
  multiplicative chaos and applications: a review.
  {\newblock}\tmtextit{Probab. Surv.}, 11:315--392, 2014.{\newblock}
  
  \bibitem[64]{Vargas2010}R.~Robert  and  V.~Vargas. {\newblock}Gaussian
  multiplicative chaos revisited. {\newblock}\tmtextit{Ann. Probab.},
  38(2):605--631, 2010.{\newblock}
  
  \bibitem[65]{Ryan2002}R.~A.~Ryan. {\newblock}\tmtextit{Introduction to
  tensor products of Banach spaces}. {\newblock}Springer Monographs in
  Mathematics. Springer-Verlag London, Ltd., London, 2002.{\newblock}
  
  \bibitem[66]{Sha2016}A.~Shamov. {\newblock}On Gaussian multiplicative chaos.
  {\newblock}\tmtextit{J. Funct. Anal.}, 270(9):3224--3261, 2016.{\newblock}
  
  \bibitem[67]{Simon1974}B.~Simon. {\newblock}\tmtextit{The $P (\phi)_2$
  Euclidean (quantum) field theory}. {\newblock}Princeton University Press,
  Princeton, N.J., 1974. {\newblock}Princeton Series in Physics.{\newblock}
  
  \bibitem[68]{Triebel1992}H.~Triebel. {\newblock}\tmtextit{Theory of function
  spaces. II},  volume~84  of \tmtextit{Monographs in Mathematics}.
  {\newblock}Birkh{\"a}user Verlag, Basel, 1992.{\newblock}
  
  \bibitem[69]{Triebel1997}H.~Triebel. {\newblock}\tmtextit{Fractals and
  spectra},  volume~91  of \tmtextit{Monographs in Mathematics}.
  {\newblock}Birkh{\"a}user Verlag, Basel, 1997. {\newblock}Related to Fourier
  analysis and function spaces.{\newblock}
  
  \bibitem[70]{Triebel2006}H.~Triebel. {\newblock}\tmtextit{Theory of function
  spaces. III},  volume  100  of \tmtextit{Monographs in Mathematics}.
  {\newblock}Birkh{\"a}user Verlag, Basel, 2006.{\newblock}
\end{thebibliography}

\end{document}
